\section{Non-emptiness of the hypothesis package}\label{sec:pf:stageU}

\providecommand{\stageUA}{\mathcal{A}}
\providecommand{\stageUC}{\mathcal{C}}
\providecommand{\stageUFm}{\mathcal{F}}
\providecommand{\stageULk}{\mathrm{Lk}}
\providecommand{\stageUKP}{\mathrm{KP}}
\providecommand{\stageUBd}{\mathrm{Bd}}
\providecommand{\stageUQO}{\mathrm{QO}}
\providecommand{\stageUDR}{\mathrm{DR}}
\providecommand{\stageUNN}{\mathcal{N}}
\providecommand{\stageUand}{\mathbin{\wedge}}
\providecommand{\stageUEv}{\mathrm{Ev}}
\providecommand{\stageULg}{\mathrm{Lg}}
\providecommand{\stageUBu}{\mathrm{Bu}}
\providecommand{\stageUSel}{\mathrm{Sel}}
\providecommand{\stageULin}{\mathrm{Lin}}
\providecommand{\stageULB}{\mathrm{LB}}
\providecommand{\stageUEx}{\mathrm{Ex}}
\providecommand{\stageUDes}{\mathrm{Des}}
\providecommand{\stageULay}{\mathrm{Lay}}
\providecommand{\stageUPts}{\mathrm{Pts}}
\providecommand{\stageUKits}{\mathrm{Kits}}
\providecommand{\stageUHosts}{\mathrm{Hosts}}
\providecommand{\stageUE}{\mathcal{E}}
\providecommand{\stageUpd}{p_{\mathrm{dep}}}
\providecommand{\stageUpA}{p_{A}}

\subsection{Introduction}\label{sec:pf:stageU:intro}

This section proves \cref{thm:hyps-nonempty}: the hypothesis package $\sfname{HypsAlt}(N,\stageUDR)$ of \cref{def:hyps},
with $N=\sfname{NL}$ and $\stageUDR=\sfname{DefiniteSideCondition}$, is non-empty under the hypotheses
listed there. The proof is an assembly. Each field of the package is produced by a result proved elsewhere, and the
task of this section is to show how these results fit together and to supply the ingredients that are not provided
elsewhere. The proof has three layers.
\begin{enumerate}[label=(\arabic*)]
\item \emph{The main theorem} (\cref{pf:stageU:main}). From the two bundles $P,F$ of classical results and the
inputs $h_A,h_M,h_X,h_L,h_N$ one obtains the separation of rows by theta pairings ($\sfname{FunctionalsSeparate}$),
the rationality of rows ($\sfname{RowRat}$), the slope bound $\sfname{ConstantFunctionFits}$, a grid constant $c^*$, the
open parameter statements, and the functional and covolume rows; then \cref{pf:stageU:unram} applies.
\item \emph{The package from unramified exponents} (\cref{pf:stageU:unram}). The member-bundle theorem, the open
statements, the Galois link and the row package give the twenty-six linked statements of $\sfname{AllRowsAlt}$; then
\cref{pf:stageU:allrows} applies.
\item \emph{The package from the rows} (\cref{pf:stageU:allrows}). The twenty-six linked statements and one further row,
which bounds the sizes of the level primes, are combined into a single linked row; the class-selection theorem
(\cref{pf:stageU:classsel}) turns it into a family class $\stageUA$ with a selection $\sigma$ of bundles satisfying all
rows; eight uniform statements fill the remaining fields.
\end{enumerate}
The mathematical content specific to this section is the following.
\begin{itemize}
\item The class-selection theorem. Its main part is the \emph{non-emptiness field} of the package: designed families
exist in the selected class for every requirement admitting hosts on windows (\cref{def:hostexw}). This is proved by
turning hosts into families of global data. The level primes $T_u$ of the words are chosen in short intervals among
the primes $\pi\bar\pi$ with $\pi\equiv1$ in an imaginary quadratic field; the position of the intervals is governed
by the ratio $D_*/d_u$ of main terms, which is bounded below through the Siegel mass ratio and lower bounds for
Vandermonde products (\cref{pf:stageU:winlow}); an auxiliary prime $q_0\equiv1$ of polynomial size is supplied by a
Linnik-type bound (\cref{pf:stageU:smallprime}).
\item The size row: eventually $\log T_u\le C_{Tu}\,m\,L$ (\cref{pf:stageU:hTu}).
\item Uniform statements of analytic number theory (zero density near $\Re s=1$, the Landau--Page theorem, the
truncated explicit formula, \cref{pf:stageU:litthin}), the hermitian volume ratio deduced from the Siegel--Weil formula
(\cref{pf:stageU:volratio}) and the main-term bound for equal words (\cref{pf:stageU:mainrow}).
\end{itemize}
Several ingredients of these proofs are large constructions of auxiliary records (bases, cores, ledger parts, kits of
points) or long explicit verifications. They are stated precisely below and cited; their proofs are given elsewhere in
this part.

\paragraph{Organisation.} \Cref{sec:pf:stageU:vocab} fixes notation for the eventual quantifiers, rows, linked rows,
member existence and design choices. \Cref{sec:pf:stageU:calculus} collects the calculus of rows.
\Cref{sec:pf:stageU:classsel} proves the class-selection theorem from two propositions which are proved in
\cref{sec:pf:stageU:designs} (existence of designs) and \cref{sec:pf:stageU:layers} (layered data and the window lower
bound). \Cref{sec:pf:stageU:size} proves the size row and \cref{sec:pf:stageU:uniform} the uniform statements.
\Cref{sec:pf:stageU:assembly} assembles the package and proves the main theorems. \Cref{sec:pf:stageU:other} proves results
of the formal development which are used in other sections of this part: properties of downward theta lifts, the Schr\"odinger
model of the doubled space, Hecke point laws and third-layer statements, Galois-theoretic conditions at members, the calculus of
steps, the moving of layers, facts on Hecke operators, class sums and characters, and lemmas on rows, columns and bundles.

\paragraph{Conventions.} We use the notions of \cref{sec:families} freely: global data $D$, size data with scale
$s=(x,\eps,J)$, $L=x^{\eps}$, $b=e^{L}$, families $\stageUFm$ with exponent $\eps$, threshold $x_1$, index sets $I_x$ and
members $(D_{x,i},L_{x,i})$, family classes, and the closedness quantifier $\sfname{Closed}(\stageUC,P)$ of
\cref{def:families:closed}. For $x>x_1$ and $i\in I_x$ we write $O_{x,i}$ for the \emph{object} of the member (its size
data, ties, division data and extra data). Records whose internal structure plays no role in an argument are
described only by their role; the formal counterparts of the definitions are indicated at their end.

% ---------------------------------------------------------------------------
\subsection{Notation: eventual statements, rows and designs}\label{sec:pf:stageU:vocab}

\paragraph{Contexts.}
A \emph{constant package} $C$ is a record of real constants and lists of real constants; it is \emph{standard} if these
satisfy a fixed list of numerical conditions, and \emph{large} for a transported constant $K_0$ in the sense of
\cref{def:families:closed}. A \emph{context} $K$ consists of a scale $K.s$, a standard constant package $K.C$, reals
$c_D,c_h$, a \emph{pre-character datum} $K.P$ (ranks $R$, a host $H$, a finite set of \emph{words} $u$ with
Harish-Chandra parameters at the real places), an integer, integer weight vectors $a,b$ on the index set of the host,
a binary code of the words (an injection of the words into $\{0,1\}^{k-2}$), conditions relating these data to the
scale, and a \emph{level choice} (\cref{sec:pf:stageU:designs}).

\begin{definition}[Eventual statements]\label{pf:stageU:ev}
\begin{enumerate}[label=(\alph*)]
\item For a property $p$ of contexts, $\stageUEv_K\,p(K)$ (``eventually in the scale'') means: there is $\eps_0>0$ such
that for all $\eps\in(0,\eps_0]$ with $\eps<\frac1{52}$ and all $C_8\ge1$ there is a transported constant $K_0$ such
that for every standard constant package $C$ with $C.\mathrm{Cpde}=C_8$ which is $K_0$-large there is $c_1$ such that for
every $c_D\ge c_1$ there is $c_2$ such that for every $c_h\ge c_2$ there is $x_0$ with $p(K)$ for all contexts $K$ with
constants $(C,c_D,c_h)$, exponent $\eps$ and $K.s.x\ge x_0$. The statement $\stageUEv^{\mathrm g}_K\,p(K)$ is defined in the
same way, with three further constants $C_5,c_6,C_b$ quantified universally together with $C_8$ and the additional
hypothesis that $K$ is $(C_5,c_6,C_b)$-good.
\item For a context $K$ and a property $p$ of bases $B$, frames $f$ and twisting data $T$ (\cref{def:families:twist}),
$\stageULg_K\,p$ (``for large twisting data'') means: there are functions $S_1,P_1,\mathcal E,S_2$ assigning to a
pre-base, a frame and one or two subfields of $\C$ finite sets of algebraic numbers ($S_1,S_2$) or of rational primes
($P_1$, and $\mathcal E(\ell)$ for every prime $\ell$), such that $p(B,f,T)$ holds for every base $B$ over the pre-character
datum of $K$ whose coefficient fields contain the prescribed numbers, every frame $f$, and every twisting datum $T$
whose prime set contains $P_1$ and whose auxiliary sets avoid $\mathcal E$. We put
$\sfname{EvLgG}(p)=\stageUEv^{\mathrm g}_K\,\stageULg_K\,p_K$, and we also use $\stageUEv_K\,\stageULg_K\,p_K$.
\item A \emph{parameter tuple} is $\kappa=(\eps,C,c_D,c_h,C_5,c_6,C_b)$; it is \emph{positive} if $C_5,c_6,C_b>0$, and
$\kappa\triangleright K$ means that $K$ has the constants $(C,c_D,c_h)$ and the exponent $\eps$ of $\kappa$ and satisfies the
goodness conditions with $(C_5,c_6,C_b)$. For a property $q$ of parameter tuples, $\stageUEv_\kappa\,q(\kappa)$ means:
there is $\eps_0>0$ such that for all $\eps,C_8,C_5,c_6,C_b$ with $0<\eps\le\eps_0$, $\eps<\frac1{52}$, $1\le C_8$ there is
$K_0$ such that for every standard $K_0$-large $C$ with $C.\mathrm{Cpde}=C_8$ there is $c_1$ such that for all $c_D\ge c_1$
there is $c_2$ such that $q(\eps,C,c_D,c_h,C_5,c_6,C_b)$ holds for all $c_h\ge c_2$. For a property $p$ of parameter
tuples and contexts,
\[
\stageUEv_{\kappa,K}\,p(\kappa,K)\ :=\ \stageUEv_\kappa\bigl(\kappa\text{ positive}\Rightarrow\exists x_0\ \forall K,\
(\kappa\triangleright K\wedge K.s.x\ge x_0)\Rightarrow p(\kappa,K)\bigr).
\]
\end{enumerate}
\formal{GlobalTheory.Mem.Top.Ev}\formal{GlobalTheory.Mem.Top.EvG}\formal{GlobalTheory.Mem.Top.large₂}
\formal{GlobalTheory.Fc.ParEv}\formal{GlobalTheory.Fc.EvP}
\end{definition}

\paragraph{Bases, cores, slots and bundles.}
As in \cref{sec:families:nonempty}, a \emph{base} $B$ fixes a pre-base (ranks, host, hermitian spaces, words) and a
coefficient field with its subfields; $f$ denotes a frame and $T$ a twisting datum over $B$. A \emph{base predicate} is a
predicate $\stageUKP(B,f)$; the one used in the package is $\stageUKP_0(B,f)$: \emph{the subfield $K_s$ of the
coefficient field is exact for some field datum and $f$, and the pre-base is pinned}. A tuple $X$ of \emph{stages} over $B$
(columns, row rest, division, and three further stages) determines a global datum $X.D$; a \emph{core} $c$ over
$(K,X,f,T)$ is a record of characters and places together with an \emph{object} $c.O$ over $X.D$. A \emph{slot} is a
predicate $Q(K,X,f,T,c)$ on cores; $Q\stageUand R$ is the pointwise conjunction. A \emph{bundle point} $\beta$ over a
context $K$ is a core together with a layer and a ledger part (a point with size data at the scale $K.s$) and
compatibility conditions with $K$; it has a datum $\beta.D$ and an object $\beta.O$, and it \emph{lies in} $Q$ if its core
satisfies $Q$. A \emph{bundle type} $\stageUBd$ assigns a type $\stageUBd(D,O,c)$ to a datum, an object and a real
number $c$; the one used in the package is the bundle condition of \cref{sec:families:rows}. A \emph{linkage}
$\stageULk$ is a predicate $\stageULk(K,\beta,c,b)$ on contexts, bundle points, reals and bundles. A \emph{raw selection}
$\sigma$ for a class $\stageUA$ chooses a bundle $\sigma(\stageUFm,x,i)\in\stageUBd(D_{x,i},O_{x,i},C_{sz})$ for every
$\stageUFm\in\stageUA$, $x>x_1$ and $i\in I_x$, where $C_{sz}$ is a size constant of $\stageUFm$.

\paragraph{Rows.}
A parameter tuple $\kappa$ is \emph{attached} to a family $\stageUFm$, written $\kappa\sim\stageUFm$, if the exponent of
$\kappa$ is that of $\stageUFm$, $\kappa$ is positive, $\kappa.C$ is standard, and the lists of constants of $\stageUFm$ are
those of $\kappa.C$ (the entry $C_{Tu}$ being set to $\kappa.C.C_{Tu}$). For a property $\stageUQO$ of data and objects,
$\stageUC[\stageUQO]$ denotes the class of families satisfying the page conditions (\cref{eq:families:rankval} and the
fixed numerical constants of the pages) all of whose members $(D_{x,i},O_{x,i})$, $x>x_1$, satisfy $\stageUQO$, the
condition $\mathrm{QD}$ (every tuple is used; the member is twisted at all places with good twists; the lattices
$\widetilde M(\mu)$ are stable under $\OO_{\mathit{LL}}$), and a list of point conditions (twisting, tameness, Minkowski
divisibility, twist goodness, duties and walls).

\begin{definition}[Rows]\label{pf:stageU:rows}
Let $Q$ be a slot, $\stageULk$ a linkage for a bundle type $\stageUBd$, and $E(\kappa,D,O,c,b)$ a property of bundles
depending on a parameter tuple.
\begin{enumerate}[label=(\alph*)]
\item A bundle $b\in\stageUBd(D,O,c)$ is \emph{built} from $Q$ along $\stageULk$ at $\kappa$, written
$b\in\stageUBu_\kappa(Q,\stageULk)$, if there are a context $K$ with $\kappa\triangleright K$, a bundle point $\beta$ over $K$ in
$Q$ and $b'$ with $\stageULk(K,\beta,c,b')$ and $(D,O,b)=(\beta.D,\beta.O,b')$.
\item $\stageUE(Q,\stageULk;E)$ (``$E$ holds eventually at linked bundles'') is the statement
$\stageUEv_{\kappa,K}$ of: for every bundle point $\beta$ over $K$ in $Q$ and all $c,b$ with $\stageULk(K,\beta,c,b)$,
$E(\kappa,\beta.D,\beta.O,c,b)$.
\item $\stageUSel(Q,\stageULk,E;\stageUA,\sigma)$ (``$\sigma$ is an $E$-selection from $(Q,\stageULk)$''): for every
$\stageUFm\in\stageUA$ there is $\kappa\sim\stageUFm$ with $\stageUFm\in\stageUC[\mathrm{True}]$ such that for all $x>x_1$ and
$i\in I_x$ the bundle $\sigma(\stageUFm,x,i)$ lies in $\stageUBu_\kappa(Q,\stageULk)$ and satisfies $E(\kappa,\cdot)$.
\item For a property $H(\stageUA,\sigma)$ of classes with raw selections, $\sfname{Row}(Q,\stageULk,H)$ means: there is $E$
with $\stageUE(Q,\stageULk;E)$ such that $H(\stageUA,\sigma)$ holds for every $E$-selection $\sigma$ from $(Q,\stageULk)$.
\end{enumerate}
\formal{GlobalTheory.Fc.Built}\formal{GlobalTheory.Fc.RowE}\formal{GlobalTheory.Fc.IsSel}\formal{GlobalTheory.Fc.Row}
\end{definition}

\begin{definition}[Linked rows]\label{pf:stageU:linked}
Let $\stageUKP$ be a base predicate, $Q,Q',R$ slots and $\stageULk$ a linkage.
\begin{enumerate}[label=(\alph*)]
\item $\stageULB_{\stageUKP}(Q,R)$ is the statement $\stageUEv_K\,\stageULg_K$ of: if $\stageUKP(B,f)$, then every core over
stages over $B$ which satisfies $Q$ satisfies $R$. $\stageULB^{\mathrm g}_{\stageUKP}(Q,R)$ is the same with
$\stageUEv^{\mathrm g}$.
\item For a property $r$ of slots, $\stageULin_{\stageUKP}(Q,Q';r):=\exists R,\ \stageULB_{\stageUKP}(Q,R)\wedge r(Q'\stageUand R)$;
we abbreviate $\stageULin_{\stageUKP}(Q,Q';\stageULk,H):=\stageULin_{\stageUKP}(Q,Q';S\mapsto\sfname{Row}(S,\stageULk,H))$.
\item $\stageUEx(Q,\stageULk)$ is the statement $\stageUEv_{\kappa,K}$ of: every bundle point $\beta$ over $K$ in $Q$
carries, for every $c\in\R$, a bundle $b$ with $\stageULk(K,\beta,c,b)$. Then
$\sfname{ExBL}(\stageUKP,Q,Q',\stageULk):=\stageULin_{\stageUKP}(Q,Q';S\mapsto\stageUEx(S,\stageULk))$.
\item For a parameter $p$ in the sense of \cref{def:families:rows}, with slots $Q_{\mathrm b}(p)\supseteq Q_{\mathrm g}(p)$
of built and of good points and bundle linkage $\stageULk_p$, and a property $r$ of slots and linkages,
$\sfname{LinkedAt}(p,r):=\stageULin_{\stageUKP_0}(Q_{\mathrm b}(p),Q_{\mathrm g}(p);S\mapsto r(S,\stageULk_p))$.
\end{enumerate}
\formal{GlobalTheory.Fc.LargeBaseCondition}\formal{GlobalTheory.Fc.Lg1}\formal{GlobalTheory.Fc.RowL}
\formallabel{GlobalTheory.Fc.ExB}\formallabel{GlobalTheory.Fc.ExBL}\formal{GlobalTheory.Fc.RL.LinkedAt}
\end{definition}

This agrees with \cref{def:families:rows}. We also use the classes attached to properties of bundles.

\begin{definition}[Classes attached to bundle properties]\label{pf:stageU:classes}
For $Q,\stageULk,E$ as in \cref{pf:stageU:rows} let $\stageUQO_{Q,\stageULk,E}(\kappa)$ be the property of data and objects
\[
\begin{aligned}
\stageUQO_{Q,\stageULk,E}(\kappa)(D,O)\ :\iff\ &\kappa\text{ positive}\Rightarrow\forall c\ \exists b\in
\stageUBu_\kappa(Q,\stageULk)\cap\stageUBd(D,O,c)\\
&\text{with }E(\kappa,D,O,c,b).
\end{aligned}
\]
For a family $\stageUQO$ of properties of data and objects indexed by parameter tuples let $\stageUA[\stageUQO]$ be the
class of families $\stageUFm$ for which there is $\kappa\sim\stageUFm$ with $\stageUFm\in\stageUC[\stageUQO(\kappa)]$, and put
$\stageUA_{Q,\stageULk,E}:=\stageUA[\stageUQO_{Q,\stageULk,E}]$.
\formallabel{GlobalTheory.Fc.QOf}\formallabel{GlobalTheory.Fc.clsP}\formallabel{GlobalTheory.Fc.cls}
\end{definition}

\paragraph{Member existence.}
For a base predicate $\stageUKP$ and properties $p,q$ of $(K,B,f,T)$, the \emph{step} from $p$ to $q$ is the statement
$\sfname{EvLgG}$ of: if $\stageUKP(B,f)$, $T$ is forced (its prime set contains the type primes, the prime $d$ of the host
and all primes $\le X''(n)$) and $p$ holds, then $q$ holds.

\begin{definition}[Member existence]\label{pf:stageU:memberE}
Let $\stageUKP$ be a base predicate and $Q,Q'$ slots.
\begin{enumerate}[label=(M\arabic*),start=0]
\item (bases) $\mathrm{M}_0(\stageUKP)$: (i) $\stageUEv^{\mathrm g}_K$ of: there are a pre-base and a frame over $K$ such that
for all finite sets $S_1,S_2$ of algebraic numbers (depending on subfields of $\C$) some base on this pre-base satisfies
$\stageUKP$ and has coefficient fields containing $S_1$ and $S_2$; and (ii) $\stageUKP$ is stable under enlargement of
the base.
\item (first step) $\mathrm{M}_1(\stageUKP,Q)$: the step from ``$T$ is untwisted'' to ``a core in $Q$ exists over
$(B,f,T)$''.
\item (export) $\mathrm{M}_2(\stageUKP,Q)$: $\sfname{EvLgG}$ of: if $\stageUKP(B,f)$ and $T$ is untwisted, every core in $Q$
satisfies the export condition (existence of good twists for all coprime powers, and the Bloch--Kato set conditions at
the twisted places).
\item (second pass) $\mathrm{M}_3(\stageUKP,Q,Q')$: $\sfname{EvLgG}$ of: if $\stageUKP(B,f)$, $T$ is untwisted and $c$ is a
core in $Q$ satisfying the export condition, then for every finite exclusion function and every finite set of
algebraic numbers there are an enlargement $B'$ of $B$, the same frame, and a twisting datum $T'$ with the same prime
set, avoiding the exclusions and realising the numbers, over which there is a core in $Q$, twisted at all places with
good twists, and in $Q'$.
\end{enumerate}
$\sfname{MemberE}$ of a slot family is the conjunction of $\mathrm{M}_0(\stageUKP_0)$, $\mathrm{M}_1(\stageUKP_0,Q_{\mathrm b})$,
$\mathrm{M}_2(\stageUKP_0,Q_{\mathrm b})$ and $\mathrm{M}_3(\stageUKP_0,Q_{\mathrm b},Q_{\mathrm g})$, as in
\cref{sec:families:nonempty}.
\formallabel{GlobalTheory.Mem.Top.GroupRetwisting.BaseCond}\formallabel{GlobalTheory.Mem.Top.GroupRetwisting.FirstG}
\formallabel{GlobalTheory.Mem.Top.GroupRetwisting.ExportF}\formallabel{GlobalTheory.Mem.Top.GroupRetwisting.SecondG}
\end{definition}

\paragraph{Design choices.}
$\sfname{MinkDiv}$ is the statement $\forall n,\ \mathrm{MinkMult}(n,\mathrm{minkM}(n))$ about Minkowski-type multipliers; it
is the hypothesis of the field $\mathsf{hEx}$ of the package.

\begin{definition}[Designed families]\label{pf:stageU:designchoices}
For a class $\stageUC$, a real $C_0$ and a function $c_T:\R\times\R\to\R$, $\stageUDes(\stageUC,C_0,c_T)$ means: there is
$\eps_0>0$ such that for every $\eps\in(0,\eps_0]$ and every transported constant $K$ there are a family
$\stageUFm\in\stageUC$ with exponent $\eps$ and reals $C_H\ge0$, $C_5>0$, $C_b>0$ such that $\stageUFm$ is large for
$(K,C_H)$, $\stageUFm.C_F=C_H$, $\stageUFm.\mathrm{CYL}\le9$, $c_T(\eps,\stageUFm.\mathrm{CYL})\le\stageUFm.C_{Tu}$, the number of
words is eventually at most $C_0\Gamma^2x^{C_H}+1$, and $\stageUFm$ has the following design properties: thick sets,
chosen sets $T$, chosen bad characters (with $C_b$), chosen $QD$-data (with $C_5$), chosen codes (with $C_0,C_H$), boxes,
$QD$-rows, divisibility of the multipliers $D_\Gamma$, and tame distinctness. For a class $\stageUA$,
\[
\stageUEx^{\mathrm{des}}(\stageUA,c_T):\iff\forall\lambda>1\ \forall\mathcal G,\quad
\sfname{HostExW}(\lambda,\mathcal G)\Longrightarrow
\stageUDes\bigl((\stageUA\cap\mathcal G.\mathrm{cls}(\lambda))\cap\sfname{TwGood},3,c_T\bigr),
\]
where $\mathcal G$ ranges over requirements (\cref{def:requirement}).
\formallabel{GlobalTheory.EX.DC.DesignChoices}\formallabel{UpperTheory.Top.ExD}
\end{definition}

The field $\mathsf{hEx}$ of the package asserts $\sfname{MinkDiv}\Rightarrow\sfname{ExS}(\stageUA,c_T)$, where
$\sfname{ExS}$ is defined as $\stageUEx^{\mathrm{des}}$ with the copies of the design properties of a different module of the
formal development; the two versions are compared in \cref{pf:stageU:frozen}.

% ---------------------------------------------------------------------------
\subsection{The calculus of rows}\label{sec:pf:stageU:calculus}

A row for a slot is a row for every smaller slot: if $Q\subseteq Q'$, then $\stageUE(Q',\stageULk;E)$ implies
$\stageUE(Q,\stageULk;E)$ (fewer bundle points) and $\stageUBu_\kappa(Q,\stageULk)\subseteq\stageUBu_\kappa(Q',\stageULk)$, so that
$E$-selections from $(Q,\stageULk)$ are $E$-selections from $(Q',\stageULk)$; hence $\sfname{Row}(Q',\stageULk,H)$ implies
$\sfname{Row}(Q,\stageULk,H)$ (\lref{GlobalTheory.Fc.Row.strengthen}, which states exactly this). Similarly
$\stageUEx(Q,\stageULk)$ is antitone in $Q$. Two eventual statements of the same kind imply the eventual statement of the
conjunction: take the minimum of the $\eps_0$, the supremum of the transported constants, the maxima of the thresholds
and the unions of the exceptional sets. Finally $\sfname{Row}(Q,\stageULk,H)$ and $\sfname{Row}(Q,\stageULk,H')$ imply
$\sfname{Row}(Q,\stageULk,H\wedge H')$, with the conjunction of the two bundle properties. These facts are used without
further comment.

\begin{lemma}[Calculus of rows]\label{pf:stageU:rowcalc}
Let $\stageUKP$ be a base predicate, $Q,Q',R$ slots, $\stageULk$ a linkage and $E,E'$ properties of bundles.
\begin{enumerate}[label=(\alph*)]
\item If $\stageULB_{\stageUKP}(Q,R)$ and $\sfname{Row}(Q',\stageULk,H)$, then $\stageULin_{\stageUKP}(Q,Q';\stageULk,H)$.
\item If $\stageULin_{\stageUKP}(Q,Q';\stageULk,H)$ and $\stageULin_{\stageUKP}(Q,Q';\stageULk,H')$, then
$\stageULin_{\stageUKP}(Q,Q';\stageULk,H\wedge H')$.
\item If $\sfname{ExBL}(\stageUKP,Q,Q',\stageULk)$ and $\stageULin_{\stageUKP}(Q,Q';\stageULk,H)$, then there is a slot $R$ with
$\stageULB_{\stageUKP}(Q,R)$, $\stageUEx(Q'\stageUand R,\stageULk)$ and $\sfname{Row}(Q'\stageUand R,\stageULk,H)$.
\item If $\stageUEx(Q,\stageULk)$ and $\stageUE(Q,\stageULk;E)$, then $\stageUEv_{\kappa,K}$ of: every bundle point $\beta$ over
$K$ in $Q$ satisfies $\stageUQO_{Q,\stageULk,E}(\kappa)(\beta.D,\beta.O)$.
\item Let $\stageUFm\in\stageUA_{Q,\stageULk,E}$ and $\kappa\sim\stageUFm$ with $\stageUFm\in\stageUC[\stageUQO_{Q,\stageULk,E}(\kappa)]$.
For all $x>x_1$ and $i\in I_x$ there is $b\in\stageUBu_\kappa(Q,\stageULk)\cap\stageUBd(D_{x,i},O_{x,i},C_{sz})$ satisfying
$E(\kappa,\cdot)$. Consequently, if $\sigma_{Q,\stageULk,E}(\stageUFm,x,i)$ denotes such a bundle (for a chosen $\kappa$), then
$\sigma_{Q,\stageULk,E}$ is an $E$-selection from $(Q,\stageULk)$ for $\stageUA_{Q,\stageULk,E}$.
\item If $\stageUEx(Q,\stageULk)$ and $\sfname{Row}(Q,\stageULk,H)$, then there is $E$ satisfying the conclusion of
\textup{(d)} such that $H(\stageUA_{Q,\stageULk,E},\sigma_{Q,\stageULk,E})$.
\item If $E\Rightarrow E'$ pointwise, every $E$-selection from $(Q,\stageULk)$ is an $E'$-selection; every $E$-selection
from $(Q\stageUand R,\stageULk)$ is an $E$-selection from $(Q,\stageULk)$.
\item $\stageUC[\stageUQO]$ is monotone in $\stageUQO$, and every family in a class $\stageUA[\stageUQO]$ fulfils the seven
duties (page conditions, a global twisting fact, tameness, walls, all tuples used, twisted at all places) and the
stability of the lattices $\widetilde M(\mu)$ under $\OO_{\mathit{LL}}$ eventually.
\end{enumerate}
\formallabel{GlobalTheory.Fc.Row.rowL}\formallabel{GlobalTheory.Fc.RowL.elim}\formallabel{GlobalTheory.Fc.QOf₂}
\formallabel{GlobalTheory.Fc.cls_at}\formallabel{GlobalTheory.Fc.pick}\formallabel{GlobalTheory.Fc.pick_isSel}
\formallabel{GlobalTheory.Fc.Row.elim}\formallabel{GlobalTheory.Fc.IsSel.mono}\formallabel{GlobalTheory.Fc.isSel_of_sAnd}
\formallabel{GlobalTheory.Fc.familyClassWithFurtherProperties_mono}\formallabel{GlobalTheory.Fc.sevenDuties_of_class}
\end{lemma}

\begin{proof}
(a) By antitonicity, $\sfname{Row}(Q',\stageULk,H)$ gives $\sfname{Row}(Q'\stageUand R,\stageULk,H)$, and $R$ is the required
witness. (b) This is \lref{GlobalTheory.Fc.RowL.and}: given witnesses $R_1,R_2$, the slot $R_1\stageUand R_2$ satisfies
$\stageULB_{\stageUKP}(Q,R_1\stageUand R_2)$, and both rows hold for the smaller slot $Q'\stageUand R_1\stageUand R_2$; combine them.
(c) Let $R_1$ be a witness of $\sfname{ExBL}$ and $R_2$ one of the linked row; then $R=R_1\stageUand R_2$ works,
$\stageUEx(Q'\stageUand R,\stageULk)$ following from $\stageUEx(Q'\stageUand R_1,\stageULk)$ and the row from
$\sfname{Row}(Q'\stageUand R_2,\stageULk,H)$ by antitonicity. (d) Combine the two eventual statements. Let $\kappa$ be positive
with $\kappa\triangleright K$, let $\beta$ be a bundle point over $K$ in $Q$, and $c\in\R$. Existence gives $b$ with
$\stageULk(K,\beta,c,b)$; then $b\in\stageUBu_\kappa(Q,\stageULk)$ (witnessed by $K,\beta,b$), and $\stageUE(Q,\stageULk;E)$ gives
$E(\kappa,\beta.D,\beta.O,c,b)$. (e) Since $\stageUFm\in\stageUC[\stageUQO_{Q,\stageULk,E}(\kappa)]$, every member satisfies
$\stageUQO_{Q,\stageULk,E}(\kappa)(D_{x,i},O_{x,i})$ by the definition of $\stageUC[\cdot]$, and $\kappa$ is positive because
$\kappa\sim\stageUFm$; apply the definition with $c=C_{sz}$. For the second statement use the chosen $\kappa$ for every
$\stageUFm$: then $\kappa\sim\stageUFm$, $\stageUFm\in\stageUC[\mathrm{True}]$ by the monotonicity in (h), and the chosen bundles
are built and satisfy $E$ by the first statement. (f) Take the $E$ given by the row; (d) holds by existence, and by (e)
$\sigma_{Q,\stageULk,E}$ is an $E$-selection for $\stageUA_{Q,\stageULk,E}$, so the row gives $H$. (g) The first statement is
immediate from the definition; for the second, a bundle built from a point in $Q\stageUand R$ is built from a point in
$Q$. (h) Monotonicity is immediate from the definition of $\stageUC[\cdot]$. If $\stageUFm\in\stageUA[\stageUQO]$, choose $\kappa$
with $\stageUFm\in\stageUC[\stageUQO(\kappa)]$; the duties and the stability are part of the conditions defining
$\stageUC[\cdot]$ (\lref{GlobalTheory.EX.DQ.familyClass_sevenDuties_and_stableLattices}, which extracts them).
\end{proof}

\begin{lemma}[Second pass with a large-base slot]\label{pf:stageU:secondg}
Let $\stageUKP$ be a base predicate and $Q,Q',R$ slots. If $\mathrm{M}_0(\stageUKP)$, $\stageULB^{\mathrm g}_{\stageUKP}(Q,R)$
and $\mathrm{M}_3(\stageUKP,Q,Q')$, then $\mathrm{M}_3(\stageUKP,Q,Q'\stageUand R)$.
\formallabel{GlobalTheory.Mem.Top.GroupRetwisting.SecondG.sAnd_of_largeBase}
\formallabel{GlobalTheory.T.secondG_sAnd_of_largeBase}
\end{lemma}

\begin{proof}
Both $\stageULB^{\mathrm g}_{\stageUKP}(Q,R)$ and $\mathrm{M}_3$ are of the form $\stageUEv^{\mathrm g}_K\stageULg_K(\cdot)$; work at a
context at which both hold, with exceptional data $(S_1^R,P_1^R,\mathcal E^R,S_2^R)$ for the large-base condition and
$(S_1^S,P_1^S,\mathcal E^S,S_2^S)$ for the second pass. For the conclusion use the exceptional data
$(S_1^S\cup S_1^R,\ P_1^S\cup P_1^R,\ \mathcal E^S,\ S_2^S)$. Let $(B,f,T)$ satisfy them, with $\stageUKP(B,f)$ and $T$
untwisted, let $c$ be a core in $Q$ satisfying the export condition, and let an exclusion function $\mathcal E'$ and a set
$S_2'$ of algebraic numbers be given. Apply the second pass with the exclusion function $\mathcal E'\cup\mathcal E^R$ and the
numbers $S_2'\cup S_2^R$: it gives an enlargement $B'$ of $B$, the same frame, a twisting datum $T'$ with the prime set of
$T$, avoiding $\mathcal E'\cup\mathcal E^R$, with $S_2'\cup S_2^R$ realised in the coefficient field of $B'$, and a core $c'$
over $(B',f,T')$ in $Q$, twisted at all places with good twists, and in $Q'$. Since $B'$ enlarges $B$, it has the same
pre-base and the same images of the subfields $K_s$ and $\mathit{LL}_0$ in $\C$, so the exceptional data of the large-base
condition computed from $B'$ are those computed from $B$, and $(B',f,T')$ satisfies them: $S_1^R$ lies in the common fixed
coefficient field, $S_2^R$ is realised, $P_1^R$ is contained in the common prime set and $T'$ avoids $\mathcal E^R$.
Moreover $\stageUKP(B',f)$ by the stability of $\stageUKP$ under enlargement. Hence the large-base condition gives $R$ at
$c'$, so that $c'$ also lies in $Q'\stageUand R$.
\end{proof}

The rows of the package are stated with ``working copies'' of some family predicates. Let $W_1(\stageUA)$ be the class of
$\stageUFm\in\stageUA$ which are twist-good and smooth and have divisible multipliers, tame distinctness, boxes, $QD$-rows and
$\mathrm{CYL}\le9$, and let $W_3(\stageUA)\subseteq W_1(\stageUA)$ be defined by the chosen-constant condition at level $3$,
$\sfname{TuSize}$ and $\sfname{DGamSize}$. The package uses instead the classes $D3(3,\stageUA)$ of
\cref{def:families:closed} and $W'_1(\stageUA)$, the analogue of $W_1(\stageUA)$ built from the copies of these predicates in
another module.

\begin{lemma}[Working and frozen copies]\label{pf:stageU:frozen}
\begin{enumerate}[label=(\alph*)]
\item For every class $\stageUA$, $D3(3,\stageUA)=W_3(\stageUA)$ and $W'_1(\stageUA)=W_1(\stageUA)$.
\item Every family fulfilling the seven duties of \cref{pf:stageU:rowcalc}(h) fulfils the duties required by the
package.
\item For every class $\stageUA$ and every $c_T$: if $\sfname{MinkDiv}$ implies
$\stageUDes(\stageUA\cap\sfname{TwGood},3,c_T)$, then $\sfname{MinkDiv}$ implies the corresponding statement for the copies
used in $\sfname{ExS}$.
\end{enumerate}
\formallabel{GlobalTheory.Fc.working_classes_frozen}\formallabel{GlobalTheory.Fc.copies_frozen}
\formallabel{GlobalTheory.T.workingFamClass_copies_eq_frozen}
\end{lemma}

\begin{proof}
The two copies of divisibility of the multipliers, of tame distinctness and of the bound on the number of words unfold to
the same conditions (for tame distinctness after unfolding the tame list), and the two copies of $\sfname{MinkDiv}$ have
the same definition; (a) and (c) follow by unfolding. (b) is \lref{GlobalTheory.EX.DQ.duty_seven_library}, which
states that the seven duties imply the duties of the package.
\end{proof}

% ---------------------------------------------------------------------------
\subsection{The class-selection theorem}\label{sec:pf:stageU:classsel}

The following theorem produces an actual class with a selection from linked rows, together with the non-emptiness of
the class. Its two inputs on the existence side are proved in \cref{sec:pf:stageU:designs,sec:pf:stageU:layers}:
\cref{pf:stageU:recs} (designs exist for classes $\stageUA[\stageUQO]$ once layered data exist, \cref{pf:stageU:recexpdef})
and \cref{pf:stageU:membertop} (layered data exist under the member-existence statements). Write (SM) for the Siegel mass
ratio of \cref{prop:families:mass}, and (VR) for the hermitian volume ratio of \cref{pf:stageU:volratio}.

\begin{theorem}[Class selection]\label{pf:stageU:classsel}
Assume (SM) and (VR). Let $\stageUKP$ be a base predicate, $Q,Q'$ slots, $\stageUBd$ a bundle type with a linkage
$\stageULk$, $H$ a property of classes with raw selections, and $c_T:\R\times\R\to\R$. Suppose
\begin{enumerate}[label=(\roman*)]
\item $\sfname{MinkDiv}$ implies $\mathrm{M}_0(\stageUKP)$, $\mathrm{M}_1(\stageUKP,Q)$, $\mathrm{M}_2(\stageUKP,Q)$ and
$\mathrm{M}_3(\stageUKP,Q,Q')$;
\item $\sfname{ExBL}(\stageUKP,Q,Q',\stageULk)$;
\item $\stageULin_{\stageUKP}(Q,Q';\stageULk,H)$.
\end{enumerate}
Then there are a class $\stageUA$ and a raw selection $\sigma$ for $\stageUA$ such that $\sfname{MinkDiv}$ implies
$\stageUEx^{\mathrm{des}}(\stageUA,c_T)$; every family in $\stageUA$ fulfils the seven duties and the lattice stability;
$\sigma$ is a $\mathrm{True}$-selection from $(Q',\stageULk)$; and $H(\stageUA,\sigma)$.
\formallabel{UpperTheory.Top.classSel}
\end{theorem}

\begin{proof}
By \cref{pf:stageU:rowcalc}(c) there is a slot $R$ with $\stageULB_{\stageUKP}(Q,R)$, $\stageUEx(Q'\stageUand R,\stageULk)$ and
$\sfname{Row}(Q'\stageUand R,\stageULk,H)$. By \cref{pf:stageU:rowcalc}(f) there is a property $E$ of bundles such that,
with $\stageUQO:=\stageUQO_{Q'\stageUand R,\stageULk,E}$,
\begin{equation}\label{pf:stageU:eq:QO}
\stageUEv_{\kappa,K}\Bigl(\text{every bundle point }\beta\text{ over }K\text{ in }Q'\stageUand R\text{ satisfies }
\stageUQO(\kappa)(\beta.D,\beta.O)\Bigr),
\end{equation}
and $H(\stageUA,\sigma)$ for $\stageUA:=\stageUA[\stageUQO]$ and $\sigma:=\sigma_{Q'\stageUand R,\stageULk,E}$. The duties and the
stability hold by \cref{pf:stageU:rowcalc}(h). By \cref{pf:stageU:rowcalc}(e), $\sigma$ is an $E$-selection from
$(Q'\stageUand R,\stageULk)$; by \cref{pf:stageU:rowcalc}(g) it is a $\mathrm{True}$-selection from $(Q',\stageULk)$.

It remains to show that $\sfname{MinkDiv}$ implies $\stageUEx^{\mathrm{des}}(\stageUA[\stageUQO],c_T)$. By \cref{pf:stageU:recs}
it suffices to show that $\sfname{MinkDiv}$ implies $\stageULay(\kappa\mapsto\stageUQO(\kappa)\wedge\mathrm{QD})$. Assume
$\sfname{MinkDiv}$; by (i) we have $\mathrm{M}_0(\stageUKP)$, $\mathrm{M}_1(\stageUKP,Q)$, $\mathrm{M}_2(\stageUKP,Q)$ and
$\mathrm{M}_3(\stageUKP,Q,Q')$. The large-base condition $\stageULB_{\stageUKP}(Q,R)$ is an $\stageUEv$-statement and therefore
implies the corresponding $\stageUEv^{\mathrm g}$-statement, the goodness hypothesis being simply not used
(\lref{GlobalTheory.Mem.Top.evG_of_ev}, which states this implication); so $\stageULB^{\mathrm g}_{\stageUKP}(Q,R)$ holds and
\cref{pf:stageU:secondg} gives $\mathrm{M}_3(\stageUKP,Q,Q'\stageUand R)$. \Cref{pf:stageU:membertop}, applied to $\stageUKP$,
$Q$, $Q'\stageUand R$, $\stageUQO$ and \eqref{pf:stageU:eq:QO}, gives
$\stageULay\bigl(\kappa\mapsto(\kappa\text{ positive}\Rightarrow\stageUQO(\kappa))\wedge\mathrm{QD}\bigr)$. Since
$\stageUQO(\kappa)$ is itself of the form ``$\kappa$ positive $\Rightarrow\dots$'', the property
``$\kappa$ positive $\Rightarrow\stageUQO(\kappa)$'' coincides with $\stageUQO(\kappa)$, and we obtain
$\stageULay(\kappa\mapsto\stageUQO(\kappa)\wedge\mathrm{QD})$.
\end{proof}

% ---------------------------------------------------------------------------
\subsection{Existence of designs}\label{sec:pf:stageU:designs}

\paragraph{Level choices and layered data.}
Let $P$ be a pre-character datum with ranks $R$ ($k,m,f,n$ and $r=k+m-1$) and host $H$; the host provides the primes
$Q_S,q,d$, the type primes $g_0,g_1$, two further primes $p_{de},p'_{de}$, and the imaginary quadratic field
$K=\Q(\sqrt{-d})\subseteq E$. For a word $u$ and a real place $\tau$ let $\dim(u,\tau)$ be the Weyl dimension attached to
the Harish-Chandra parameter of $u$ at $\tau$, i.e.\ the Vandermonde product of its entries divided by $\prod_{i<k}i!$.
Put
\begin{equation}\label{pf:stageU:eq:dD}
d_P(u)=\mathrm{mass}_k\,\bigl(\iota_{\mathrm{c}}(g_0)\,\iota_{\mathrm{c}}(g_1)\bigr)^{f}\prod_{\tau}\dim(u,\tau),\qquad
D_P=\mathrm{mass}_r\,\bigl(\iota_{\mathrm{r}}(g_0)\,\iota_{\mathrm{r}}(g_1)\bigr)^{f}\prod_\tau d_\sigma(\tau),
\end{equation}
and $\Gamma(C,s,R)=\exp(C_\Gamma\,m^5fL)$. Here $\mathrm{mass}_n$ is the mass of the standard hermitian lattice $\OO_E^n$,
$\iota_{\mathrm{c}}(g)$ and $\iota_{\mathrm{r}}(g)$ are the local indices at a type prime $g$ of the column level subgroup of
$\GL_k(\F_g)$ and of the row level subgroup of $\GL_{k+m-1}(\F_g)$, $d_\sigma(\tau)$ is the degree attached to the row
parameter of the base word $u_0$ at $\tau$, and $C_\Gamma$ is a constant of $C$. These are the main terms $d_u$ and $D_*$
of \cref{thm:theorem-S}. A \emph{level choice} for $P$ consists of primes $T_u\equiv1\pmod{Q_Sqd}$, one for each word
$u$, with $T_u\notin\{p_{de},p'_{de}\}$, elements $\pi_u\in\OO_K$ with $\pi_u\bar\pi_u=T_u$ and
$\pi_u\equiv1\pmod{Q_Sqd(g_0g_1)^2}$, positive integers $d^*_u$ and reals $D^*>0$, $\Gamma>1$ such that
\begin{equation}\label{pf:stageU:eq:Twin}
\Bigl(\frac{D^*}{d^*_u}\Bigr)^{\frac1{fk}}\Gamma^{-\frac2{fk}}\ \le\ T_u\ \le\
\Bigl(1+\frac1{fk}\Bigr)\Bigl(\frac{D^*}{d^*_u}\Bigr)^{\frac1{fk}}\Gamma^{-\frac2{fk}}\qquad\text{for all }u.
\end{equation}
In a context the level choice satisfies $d^*_u=d_P(u)$, $D^*=D_P$ and $\Gamma=\Gamma(C,s,R)$. Thus the level primes are
determined, up to a factor $1+\frac1{fk}$, by the ratio of the main terms.

\begin{definition}[Layered data]\label{pf:stageU:recexpdef}
For a family $\stageUQO$ of properties of data and objects indexed by parameter tuples, $\stageULay(\stageUQO)$ is the
statement $\stageUEv_\kappa$ of: there are $c_s>0$ and $x_0$ such that for every scale $s$ with exponent $\kappa.\eps$ and
$s.x\ge x_0$, every pre-character datum $P$, all weights $a,b$ and every code $c$ satisfying the goodness conditions with
$\kappa$, there are positive integers $d_u$ and a real $D_s>0$ with
\begin{equation}\label{pf:stageU:eq:recexp}
c_s\,m\,L\ \le\ \log\Bigl(\Bigl(\frac{D_s}{d_u}\Bigr)^{\frac1{fk}}\Gamma(\kappa.C,s,R)^{-\frac2{fk}}\Bigr)\qquad\text{for all }u,
\end{equation}
such that for every level choice with $d^*=d$, $D^*=D_s$ and $\Gamma=\Gamma(\kappa.C,s,R)$ there are a datum $D$ and a
\emph{layer} $Y$ of $D$ lying over $(P,\text{level choice})$, satisfying the record, ledger and twist conditions, such that
every \emph{ledger part} of $Y$ (a point of $D$ with size data) at the scale $s$, with weights $(a,b)$, valued and
canonical, satisfies $\stageUQO(\kappa)$ at its datum and object.
\formallabel{GlobalTheory.Fc.RecExP}
\end{definition}

\paragraph{Hosts, points and kits.}
For a requirement $\mathcal G$ let $\mathcal G.\mathrm{req}(C_5,s,R,H,q_0)$ be $\mathcal G$ evaluated at the scale, the frame of
$R$, the host axioms of $H$, its set $T$ and the rank $t$; it depends neither on $C_5$ nor on $q_0$ nor on the level primes
of $H$. Write $\mathrm{HG}(s,C_8,C_5,c_6,C_b,H,q_0)$ for the size and arithmetic conditions on the host in
\cref{def:hostexw} (the \emph{host conditions}), and $q_0(H)$ for the least prime $p\equiv1\pmod{r(H)Q_Sqd}$, where $r(H)$
is the least prime not dividing $Q_Sqd(p_{de}-1)(p'_{de}-1)$ (the conductor ideal of the hosts used here is trivial).

\begin{definition}[Hosts on windows, points and kits]\label{pf:stageU:kitsdef}
Let $\lambda>1$ and let $\mathcal P(C_5,s,R,H,q_0)$ be a host requirement.
\begin{enumerate}[label=(\alph*)]
\item $\sfname{HostExW}(\lambda,\mathcal P)$ is as in \cref{def:hostexw}: for every $\eps\in(0,\frac1{52})$ there are constants
$C_{T2},C_8\ge1,C_5,c_6,C_b>0,x_0$ and a set $S$ of scales with windows of exponent $\lambda$ (for every $y$ there is
$Y\ge y$ with $[Y,Y^\lambda]$ in the window set) such that every $s\in S$ with exponent $\eps$ and $s.x\ge x_0$ admits
ranks $R$ satisfying the rank conditions, a host $H$ with constant $C_{T2}$ and a prime $q_0$ with
$\mathrm{HG}(s,C_8,C_5,c_6,C_b,H,q_0)$ and $\mathcal P(C_5,s,R,H,q_0)$. Put
$\stageUHosts^{\mathrm q}(\lambda,\mathcal P):=\sfname{HostExW}(\lambda,\mathcal P\wedge(q_0=q_0(H)))$.
\item $\stageUHosts^{\mathrm w}(\lambda,\mathcal P)$: there is $\eps_0>0$ such that for all $\eps\in(0,\eps_0]$, $\eps<\frac1{52}$,
there are $C_{T2}$, $C_8\ge1$, $C_5,c_6,C_b>0$ and $K_0$ such that for every standard $K_0$-large constant package $C$ with
these $C_{T2},C_8$, eventually in $c_D$ and then in $c_h$, there are $x_0$ and a set $S$ of scales with windows of exponent
$\lambda$ such that every $s\in S$ with exponent $\eps$ and $s.x\ge x_0$ admits ranks $R$, a host $H$, $q_0=q_0(H)$ and
weights $a,b$ supported on the set $T$ of $H$, with the rank conditions, $\mathrm{HG}(s,C_8,C_5,c_6,C_b,H,q_0)$, the window
condition for $(a,b)$, the \emph{window-share condition} (a weighted sum over the candidate primes $p<X'$ of local shares,
defined through the local lattice statement of \cref{thm:families:local}, is at most $x^{-9}$), and
$\mathcal P(C_5,s,R,H,q_0)$.
\item A \emph{point} $p$ over a constant package is a datum with size data, ties, extra data and an object; a \emph{point
property} is a predicate on these. A point \emph{comes from} $(H,a,b,c,u_0)$ if its host, weights, words and word codes
are those prescribed by $(H,a,b,c)$ and its distinguished index is $u_0$. A \emph{kit} over $(C,c_D,c_h,s)$ is a
non-empty family of points at the scale $s$, valued and with arithmetic hosts, which pairwise have the same data except
the distinguished index, whose word sets are identified by bijections compatible with the codes, and which covers every
index $u_0<n$.
\item For a family of point properties $Q(\kappa)$, $\stageUPts(Q)$ is the statement $\stageUEv_\kappa$ of: there is $x_0$
such that for every scale $s$ with exponent $\kappa.\eps$ and $s.x\ge x_0$, ranks $R$ satisfying the rank conditions, every
host $H$ with constant $\kappa.C.C_{T2}$ and $\mathrm{HG}(s,\kappa.C.\mathrm{Cpde},\kappa.C_5,\kappa.c_6,\kappa.C_b,H,q_0(H))$,
weights $a,b$ with the window condition, every injective code $c$ of $M=\lceil3\Gamma^2x^{C_F}\rceil$ words and every
$u_0<n$, there is a valued canonical point at the scale $s$ with arithmetic host, coming from $(H,a,b,c,u_0)$ and having
$Q(\kappa)$. For a single point property $Q$, $\stageUPts_0(Q)$ is the same statement with the constants quantified as in
\textup{(b)}.
\item $\stageUKits(\lambda,Q)$ is defined as in \textup{(b)}, except that the conclusion at every $s\in S$ with exponent
$\eps$ and $s.x\ge x_0$ is: there is a kit over $(C,c_D,c_h,s)$ all of whose points have
$Q(\eps,C,c_D,c_h,C_5,c_6,C_b)$, the lower bound for the thick sets, and the choices of bad characters (with $C_b$) and of
$QD$-data (with $C_5$).
\end{enumerate}
\formallabel{GlobalTheory.EX.DC.HostExQW}\formallabel{GlobalTheory.EX.DC.HostWW}\formallabel{GlobalTheory.EX.Kit}
\formallabel{GlobalTheory.Fc.ObjCP}\formallabel{GlobalTheory.EX.DC.ObjC}\formallabel{GlobalTheory.Fc.KitCPW}
\end{definition}

\paragraph{Number-theoretic inputs.}

\begin{definition}[Three statements]\label{pf:stageU:threestatements}
\begin{itemize}
\item[(L)] There is $A\ge1$ such that for every integer $M\ge2$ there is a prime $p\equiv1\pmod M$ with $p\le M^A$. (This is
the residue class $1$ case of Linnik's theorem \cite{Linnik1944}.)
\item[(LP)] For all $\eps\in(0,\frac1{52})$ and $c>0$ there is $x_0$ such that for every scale $s$ with exponent $\eps$ and
$s.x\ge x_0$, all ranks $R$ satisfying the rank conditions, all integers $N,d,g_0,g_1$ with $d$ prime, $d\equiv3\pmod4$,
$d\ge100$, $0<N$, $d\mid N$, $\log N\le3x^{\eps/5}(\log x)^2$, $1\le g_i\le b^9$, and every $T_s>0$ with
$c\,m\,L\le\log T_s$, there is a prime $T$ with $T_s\le T\le(1+\frac1{fk})T_s$ of the form $T=\pi\bar\pi$, where
$\pi=\frac{\alpha+\beta\sqrt{-d}}2$ with integers $\alpha,\beta$ such that $N(g_0g_1)^2\mid\beta$ and
$2N(g_0g_1)^2\mid\alpha-2-\beta$ (that is, $\pi\in\OO_K$ and $\pi\equiv1\pmod{N(g_0g_1)^2\OO_K}$).
\item[(W)] For the class of all families, the property ``eventually every member admits a weight vector'' is closed;
here a weight vector is the input of the local lattice statement of \cref{thm:families:local} at the candidate primes.
\end{itemize}
\formallabel{GlobalTheory.EX.SmallPrimeOne}\formallabel{GlobalTheory.EX.LevelPrime}\formallabel{GlobalTheory.EX.DC.WClosed}
\end{definition}

\begin{lemma}[A Linnik-type bound]\label{pf:stageU:smallprime}
The statement \textup{(L)} holds. More precisely, the prime number theorem in arithmetic progressions with relative error
\textup{(\lref{GlobalTheory.Lin.primesRel_holds})} implies \textup{(L)}.
\formallabel{GlobalTheory.EX.smallPrimeOne_of_primesRel}
\end{lemma}

\begin{proof}
By \lref{GlobalTheory.Lin.primesAPInterval_of_primesRel} (a consequence of that prime number theorem), there are
$C_0,C_r,c_r,c>0$ with the following property: for all integers $M\ge3$ and $a$ prime to $M$ and all reals $Y,h$ with
$Y\ge(4M)^{C_0}$, $h\le1$,
\[
\rho(Y):=C_r\exp\Bigl(-c_r\frac{\log Y}{\log4M}\Bigr)+C_r\exp\bigl(-c_r\sqrt{\log Y}\bigr)\le\frac h{100}
\quad\text{and}\quad(4M)^7\le hY^{1/3},
\]
the number of primes $p\equiv a\pmod M$ in $(Y,Y(1+h)]$ is at least $c\,h\,(4M)^{-6}Y/\log Y$. Choose
$s\ge\max(C_0,21,200C_r/c_r)$ and put $A=3s^2+1$. For $M=2$ take $p=3\le2^A$. Let $M\ge3$ and $Y=(4M)^{s^2}$. Then
$\log4M\ge1$ and $\log Y=s^2\log4M\ge s^2$, so both exponentials in $\rho(Y)$ are at most $e^{-c_rs}$, and
$C_re^{-c_rs}\le\frac1{200}$ because $e^{c_rs}\ge1+c_rs>200C_r$; hence $\rho(Y)\le\frac1{100}$. Moreover
$(4M)^7\le(4M)^{s^2/3}=Y^{1/3}$ and $Y\ge(4M)^{C_0}$, since $s^2\ge21$ and $s^2\ge C_0$. With $a=1$ and $h=1$ we obtain a
prime $p\equiv1\pmod M$ in $(Y,2Y]$. Since $4\le M^2$ we have $Y\le M^{3s^2}$, so $p\le2Y\le M^{3s^2+1}=M^A$.
\end{proof}

\begin{lemma}[The auxiliary prime]\label{pf:stageU:q0bound}
Let $A\ge1$ be as in \textup{(L)}, and let $H$ be a host with trivial conductor ideal and $p_{de},p'_{de}\le N^{C}$, where
$N=Q_Sqd$. Then $q_0(H)\le N^{(3+2C)A}$.
\formallabel{UpperTheory.Top.PolynomialHostPrimeBound.smallest_prime_polynomial_bound}
\end{lemma}

\begin{proof}
By Bertrand's postulate there is a prime $r'$ with $Q<r'\le2Q$, where $Q=N(p_{de}-1)(p'_{de}-1)>0$; it does not divide
$Q$, so $r(H)\le r'\le2Q\le N\cdot Np_{de}p'_{de}$, using $N\ge2$. Hence $r(H)N\le N^3p_{de}p'_{de}\le N^{3+2C}$. By (L)
there is a prime $p\equiv1\pmod{r(H)N}$ with $p\le(r(H)N)^A$, and $q_0(H)\le p$ by minimality; thus
$q_0(H)\le N^{(3+2C)A}$.
\end{proof}

\begin{lemma}[Re-levelling the hosts]\label{pf:stageU:hostq}
Assume \textup{(L)}. For every $\lambda$ and every requirement $\mathcal G$, $\sfname{HostExW}(\lambda,\mathcal G.\mathrm{req})$
implies $\stageUHosts^{\mathrm q}(\lambda,\mathcal G.\mathrm{req})$.
\formallabel{UpperTheory.Top.hostQ}
\end{lemma}

\begin{proof}
Let $A$ be the constant of (L) and fix $\eps$. Take the constants $C_{T2},C_8,C_5,c_6,C_b,x_0$ and the set of scales given by
$\sfname{HostExW}$, and replace $C_8$ by $(3+2C_8)A\ge C_8$. For a window scale $s$ beyond $x_0$ let $(R,H,q_0)$ be the given
data. Replace $H$ by the host $H'$ whose level primes are $\{g_0,g_1,p_{de},p'_{de},q_0(H)\}$, all other data being
unchanged, and replace $q_0$ by $q_0(H)$; note that $q_0(H)$ is prime to $Q_S$ because $q_0(H)\equiv1\pmod{Q_Sqd}$. The host
conditions for $H'$ with the constant $(3+2C_8)A$ follow from those for $H$: the size conditions are monotone in the
constant; the conductor ideal is trivial and $p_{de},p'_{de}\le(Q_Sqd)^{C_8}$ are among the conditions on $H$, so
\cref{pf:stageU:q0bound} gives $q_0(H)\le(Q_Sqd)^{(3+2C_8)A}$; moreover $q_0(H)\equiv1\pmod{Q_Sqd}$ and
$q_0(H)\ne p_{de},p'_{de}$. The level primes $g_0,g_1,p_{de},p'_{de}$ split in $K$ as before, and $q_0(H)$ splits because it
is $\equiv1\pmod{Q_Sqd}$ (\lref{GlobalTheory.HostSplit.split_of_one_mod}, which states that such primes split in the
host). Finally $\mathcal G.\mathrm{req}$ does not depend on the level primes or on $q_0$.
\end{proof}

\begin{lemma}[Level primes]\label{pf:stageU:levelprime}
The statement \textup{(LP)} holds.
\formallabel{GlobalTheory.EX.levelPrime_unconditional}\formallabel{UpperTheory.T3.levelPrimeN}
\end{lemma}

\begin{proof}
We use the prime number theorem for norms in short intervals (\lref{UpperTheory.T3.shortNorm}): there are
$A',c,c'>0$ such that for every prime $d\equiv3\pmod4$ with $d\ge100$, every integer $M\ge3$ and all reals $X\ge1$ and $h$
with $\log(4dM^2)\le c\sqrt{\log X}$ and $e^{-c\sqrt{\log X}}\le h\le1$, the number of primes $p\in(X,X(1+h)]$ of the form
$p=\pi\bar\pi$ with $\pi\equiv1\pmod{M\OO_K}$ is at least $c'hX/((4dM^2)^{A'}\log X)$.

Let $\eps$ and $c_s>0$ be given and put $M=N(g_0g_1)^2$, $X=T_s$ and $h=1/V$ with $V=fk$. For large $x$ the rank conditions
give $m\ge x^{1/3}$ (as $\eps<\frac1{52}$) and $1\le V\le x^{12}$, and $g_i\le b^9$ gives $\log g_i\le9x^\eps$. Hence
$\log T_s\ge c_s\,m\,x^\eps\ge c_sx^{1/3+\eps}$ and $c\sqrt{\log T_s}\ge c\sqrt{c_s}\,x^{1/6+\eps/2}$, whereas
\[
\log(4dM^2)\le\log4+3\log N+4\log g_0+4\log g_1\le\log4+9x^{\eps/5}(\log x)^2+72x^\eps,
\quad \log V\le12\log x .
\]
Each of $\log4$, $9x^{\eps/5}(\log x)^2$ and $72x^\eps$ is at most $\frac13c\sqrt{c_s}\,x^{1/6+\eps/2}$ for large $x$ (since
$(\log x)^2\le x^{1/6}$ eventually and $\eps<\frac16$), and so is $12\log x$. Hence $\log(4dM^2)\le c\sqrt{\log X}$ and
$e^{-c\sqrt{\log X}}\le1/V=h\le1$; also $M\ge N\ge d\ge100$. So the number of primes in question in
$(T_s,(1+\frac1{fk})T_s]$ is positive. This proves (LP), with $x_0$ the maximum of the thresholds of the finitely many
eventual inequalities used and of the rank conditions (with $c$ in (LP) playing the role of $c_s$).
\end{proof}

\paragraph{Points, kits and the window-share condition.}
We cite two results. The first constructs points from layered data, the second is the closed form of the window-share
requirement.
\begin{enumerate}[label=(\alph*)]
\item \lref{GlobalTheory.Fc.objCP_of_recExP} (with its statement form
\lref{GlobalTheory.T.existsObjects_of_admissibleTwistDatum}): (LP) and $\stageULay(\stageUQO)$ imply
$\stageUPts(\kappa\mapsto Q^{\mathrm{obj}}(\stageUQO(\kappa)))$, where $Q^{\mathrm{obj}}$ adds to a property of the datum and
object of a point the twisting, tameness, Minkowski-divisibility and twist-goodness conditions of the point. Its proof
chooses Harish-Chandra parameters for the words from the host (\lref{GlobalTheory.EX.OB.chGood_of_host}), a level
choice by (LP) with $T_s=(D_s/d_u)^{1/(fk)}\Gamma^{-2/(fk)}$ and \eqref{pf:stageU:eq:recexp}
(\lref{GlobalTheory.EX.OB.exists_levCh_of_levelPrime}), size data over the datum given by $\stageULay(\stageUQO)$
(\lref{GlobalTheory.EX.OB.exists_sizeData_over}) and a ledger part of its layer
(\lref{GlobalTheory.EX.Led.exists_ledPart}).
\item (W) holds (\lref{GlobalTheory.EX.G4.wClosed}, with its statement form
\lref{GlobalTheory.T.windowShareRequirementClosedForm_proof}). Its proof verifies the numerical hypotheses of
\cref{thm:families:local} in closed form.
\end{enumerate}

\begin{lemma}[Forgetting the point property]\label{pf:stageU:objc}
For every family $Q$ of point properties, $\stageUPts(Q)$ implies $\stageUPts_0(\mathrm{True})$.
\formallabel{UpperTheory.Top.objC_true}
\end{lemma}

\begin{proof}
Given the data of $\stageUPts_0(\mathrm{True})$, use the transported constant and the thresholds provided by $\stageUPts(Q)$ for
the parameter tuple formed by the given constants; the point it provides has all required properties, and its property
$Q(\kappa)$ is not needed.
\end{proof}

\begin{lemma}[Window shares]\label{pf:stageU:hostww}
Let $\lambda>1$ and let $\mathcal P$ be a host requirement. If $\stageUHosts^{\mathrm q}(\lambda,\mathcal P)$,
$\stageUPts_0(\mathrm{True})$ and \textup{(W)} hold, then $\stageUHosts^{\mathrm w}(\lambda,\mathcal P)$.
\formallabel{UpperTheory.Top.hostWW}
\end{lemma}

\begin{proof}
Fix $\eps$ below the thresholds of $\stageUPts_0(\mathrm{True})$ and (W), take the constants $C_{T2},C_8,C_5,c_6,C_b$ of
$\stageUHosts^{\mathrm q}$, the transported constant $K_0$ of $\stageUPts_0$ and the transported constant $K_1$ of (W), and use
their supremum. Let $C$ be standard and large; for $c_D,c_h$ beyond the thresholds of $\stageUPts_0$ and of the numerical
row constants (\lref{GlobalTheory.EX.szRow_numAt}, which provides the numerical data of the members), choose at every
window scale $s$ of $\stageUHosts^{\mathrm q}$ beyond a common threshold $x_1$: the host data $(R,H,q_0)$; weights $a,b$ with
the window condition (\lref{GlobalTheory.EX.exists_weights_in_windows}, existence of such weights); an injective code
$c$ of $\lceil3\Gamma^2x^{C_F}\rceil$ words (\lref{GlobalTheory.EX.exists_injective_words}, existence of such codes,
since $2^{k-2}$ exceeds this number); and, by $\stageUPts_0(\mathrm{True})$ for every $u_0<n$, points coming from
$(H,a,b,c,u_0)$, which form a kit (\lref{GlobalTheory.EX.DC.kit_of_points}). Consider the point property $Q'$: ``if a
weight vector exists at the point, then the scale of the point admits data $(R,H,q_0,a',b')$ with all the conditions of
$\stageUHosts^{\mathrm w}$''. Every point $p$ of the kit has $Q'$: if a weight vector exists,
\lref{GlobalTheory.EX.weightWindowRequirement_host_of_point} produces from it weights $a',b'$ supported on $T$ with
the window and the window-share conditions at the scale of $p$, and the host conditions are those of $(R,H,q_0)$. The
kits form a family $\stageUFm$ whose members at a window scale $x$ are the points of the kit at $x$, whose constants are
those of $C$ and whose windows are those of the set of scales (\lref{GlobalTheory.EX.exists_famP_of_kits}, the
construction of a family from kits). This family is large for $K_1$, so (W) gives $x_W$ such that every member at
$x\ge x_W$ admits a weight vector; by $Q'$, every window scale $s$ with $s.x\ge\max(x_1+1,x_W)$ admits the required data.
\end{proof}

\begin{lemma}[Kits]\label{pf:stageU:kits}
Let $\lambda>1$, $\mathcal G$ a requirement and $Q$ a family of point properties. If
$\stageUHosts^{\mathrm w}(\lambda,\mathcal G.\mathrm{req})$ and $\stageUPts(Q)$, then
$\stageUKits(\lambda,\kappa\mapsto Q^{\mathrm{du}}(Q(\kappa))\wedge Q^{\mathrm{wall}}\wedge\mathcal G.\mathrm{pt})$, where
$Q^{\mathrm{du}}$ adds the duties, $Q^{\mathrm{wall}}$ is the wall condition and $\mathcal G.\mathrm{pt}$ is the requirement
$\mathcal G$ at the point.
\formallabel{UpperTheory.Top.kitS}
\end{lemma}

\begin{proof}
Combine the two statements as in the proof of \cref{pf:stageU:hostww}. At a window scale $s$,
$\stageUHosts^{\mathrm w}$ gives $(R,H,q_0,a,b)$ with the window-share condition and the requirement; choose an injective
code $c$ (\lref{GlobalTheory.EX.exists_injective_words}); then $q_0<X'$ by the polynomial bound for $q_0$
(\lref{GlobalTheory.EX.small_prime_lt}, which states that a prime bounded by a power of $Q_Sqd$ is $<X'$ for large
$x$). For every $u_0<n$, $\stageUPts(Q)$ gives a point $p_{u_0}$ coming from $(H,a,b,c,u_0)$ with $Q(\kappa)$. These points
have the same data except the distinguished index (\lref{GlobalTheory.EX.TieSame.same_of_from}, which states this for
valued canonical points coming from the same $(H,a,b,c)$), and their words are identified through the code, so they form
a kit. Each point has the duties, the wall condition, the thick-set bound and the two choices by
\lref{GlobalTheory.EX.DC.has_withDuty_wall_and_choices} (which deduces them from the host conditions and the
window-share condition), and the requirement because it comes from $H$.
\end{proof}

\begin{lemma}[Designed families]\label{pf:stageU:designs}
Let $\lambda>1$, $\mathcal G$ a requirement, $Q$ a family of point properties and $c_T:\R\times\R\to\R$. Assume
$\sfname{MinkDiv}$ and $\stageUKits(\lambda,\kappa\mapsto Q^{\mathrm{du}}(Q^{\mathrm{obj}}(Q(\kappa)))\wedge Q^{\mathrm{wall}}\wedge
\mathcal G.\mathrm{pt})$. Then $\stageUDes((\stageUC\cap\mathcal G.\mathrm{cls}(\lambda))\cap\sfname{TwGood},3,c_T)$, where
$\stageUC$ is the class of families $\stageUFm$ for which there is $\kappa\sim\stageUFm$ such that all members satisfy
$Q^{\mathrm{du}}(Q^{\mathrm{obj}}(Q(\kappa)))\wedge Q^{\mathrm{wall}}$ and the page conditions hold.
\formallabel{UpperTheory.Top.designS}
\end{lemma}

\begin{proof}
Let $\eps\le\min(\eps_0,\frac1{53})$ and a transported constant $K$ be given. Take the constants of the kit statement and
choose a standard constant package $C$ with the prescribed $C_{T2},C_8$ which is large both for $K_0$ and for $K$ with the
entry $C_{Tu}$ replaced by $c_T(\eps,9)$ (\lref{GlobalTheory.EX.exists_std_large}, existence of such packages).
Choose $c_D,c_h$ beyond the thresholds of the kit statement and of the numerical row constants, and at every window
scale $s$ beyond the threshold choose a kit. These kits carry the numerical data, the chosen sets $T$, the chosen codes and
$\lceil3\Gamma^2x^{C_F}\rceil$ words. They form a family $\stageUFm$ (\lref{GlobalTheory.EX.exists_famP_of_kits})
whose members satisfy all point properties common to the kits, whose constants are those of $C$, and whose windows
are those of the set of scales. Then: the page conditions hold, with $\mathrm{CYL}=9$
(\lref{GlobalTheory.EX.DC.fix_and_win_of_numAt}, which reads them off the numerical data); $\stageUFm$ is large for
$(K,C_F)$ because its constants are those of $C$; twist goodness, tame distinctness, boxes, $QD$-rows, the chosen $QD$-data,
bad characters, sets $T$, codes and the thick-set bound hold because they hold at all points of the kits; divisibility of
the multipliers holds because the Minkowski point condition gives divisibility by $\mathrm{minkM}$ under $\sfname{MinkDiv}$;
the number of words is $\lceil3\Gamma^2x^{C_F}\rceil\le3\Gamma^2x^{C_F}+1$; and $\stageUFm\in\mathcal G.\mathrm{cls}(\lambda)$ by
the windows and the requirement at the points. Finally replace the constant $C_{Tu}$ of $\stageUFm$ by $c_T(\eps,9)$; this does
not affect the properties above, and it gives $c_T(\eps,\mathrm{CYL})=c_T(\eps,9)\le C_{Tu}$ and
$(\eps,C,c_D,c_h,C_5,c_6,C_b)\sim\stageUFm$, with $C_H=C_F\ge1$.
\end{proof}

\begin{proposition}[Designs from layered data]\label{pf:stageU:recs}
Let $\stageUQO$ be a family of properties of data and objects indexed by parameter tuples, and $c_T:\R\times\R\to\R$. If
$\sfname{MinkDiv}$ implies $\stageULay(\kappa\mapsto\stageUQO(\kappa)\wedge\mathrm{QD})$, then $\sfname{MinkDiv}$ implies
$\stageUEx^{\mathrm{des}}(\stageUA[\stageUQO],c_T)$.
\formallabel{UpperTheory.Top.recS}
\end{proposition}

\begin{proof}
Assume $\sfname{MinkDiv}$, and let $\lambda>1$ and $\mathcal G$ be a requirement with $\sfname{HostExW}(\lambda,\mathcal G.\mathrm{req})$.
Write $\stageUQO'(\kappa)$ for the point property ``$\stageUQO(\kappa)\wedge\mathrm{QD}$ holds at the datum and object of the
point''.
\begin{enumerate}[label=\arabic*.]
\item (L) holds (\cref{pf:stageU:smallprime}), so \cref{pf:stageU:hostq} gives
$\stageUHosts^{\mathrm q}(\lambda,\mathcal G.\mathrm{req})$.
\item (LP) holds (\cref{pf:stageU:levelprime}) and $\stageULay(\kappa\mapsto\stageUQO(\kappa)\wedge\mathrm{QD})$ holds by
hypothesis; by the cited result (a) above, $\stageUPts(\kappa\mapsto Q^{\mathrm{obj}}(\stageUQO'(\kappa)))$, and hence
$\stageUPts_0(\mathrm{True})$ by \cref{pf:stageU:objc}.
\item With (W), \cref{pf:stageU:hostww} gives $\stageUHosts^{\mathrm w}(\lambda,\mathcal G.\mathrm{req})$, and
\cref{pf:stageU:kits} gives
$\stageUKits(\lambda,\kappa\mapsto Q^{\mathrm{du}}(Q^{\mathrm{obj}}(\stageUQO'(\kappa)))\wedge Q^{\mathrm{wall}}\wedge\mathcal G.\mathrm{pt})$.
\item \Cref{pf:stageU:designs} gives $\stageUDes((\stageUC\cap\mathcal G.\mathrm{cls}(\lambda))\cap\sfname{TwGood},3,c_T)$ for the class
$\stageUC$ described there. By definition, the members of a family in $\stageUC[\stageUQO(\kappa)]$ are required to satisfy
exactly $Q^{\mathrm{du}}(Q^{\mathrm{obj}}(\stageUQO'(\kappa)))\wedge Q^{\mathrm{wall}}$, together with the page conditions, so
$\stageUC=\stageUA[\stageUQO]$.
\end{enumerate}
Since $\lambda$ and $\mathcal G$ were arbitrary, this is $\stageUEx^{\mathrm{des}}(\stageUA[\stageUQO],c_T)$.
\end{proof}

% ---------------------------------------------------------------------------
\subsection{Layered data and the window lower bound}\label{sec:pf:stageU:layers}

Write (VM) for the statement: for every datum $D$ whose hermitian volumes are normalised (the volume of every
hermitian space is the covolume of its arithmetic group for normalised Haar measures), the volume of every positive
definite hermitian space, and of every hermitian space whose signatures have at least two positive entries, equals the
mass of the corresponding standard genus.

\begin{lemma}[Volumes and masses]\label{pf:stageU:volmass}
\textup{(VR)} implies \textup{(VM)}.
\formallabel{GlobalTheory.Mem.VolMass.volumes_eq_masses_of_unfold}
\end{lemma}

\begin{proof}
Let $D$ have normalised volumes. The normalisation in the form used in (VR) follows
(\lref{Formalism.L2.Theta.VolTieG.volTie}, which compares the two normalisations), so (VR) gives the volume ratio
for $D$. Unfolding this ratio with the global arithmetic identity gives the corresponding relation between volumes and
masses (\lref{Formalism.MainNotions.StandardLattice.volumeMass_of_unfolding}). In the definite case the mass of every
definite space is that of the standard genus (\lref{GlobalTheory.genusMassStd_all}), and the definite volume
decomposes over the double cosets of the genus (\lref{GlobalTheory.DefiniteDecompObj.doubleCosets_mass}); together
with the normalisation this gives volume $=$ mass for definite spaces
(\lref{Formalism.L2.GlobalTheoryTheory.GenusMassVolumes.mass_volume_identity_of_class_decomp}). The two relations
give (VM) (\lref{Formalism.MainNotions.StandardLattice.volume_eq_mass_of_definite}, which combines them).
\end{proof}

\begin{definition}[Layers at good contexts]\label{pf:stageU:recatgdef}
For a property $\mathrm{QL}(K,D,Y)$ of contexts, data and layers, $\stageULay^{\mathrm g}(\mathrm{QL})$ is the statement
$\stageUEv^{\mathrm g}_K$ of: there are a datum $D$ and a layer $Y$ of $D$ lying over the pre-character datum and the level
choice of $K$, satisfying the record, ledger and twist conditions, with $\mathrm{QL}(K,D,Y)$. A layer $Y$ of $D$ \emph{comes
from a core in} $Q'$ if there are stages $X$ with $X.D=D$, a frame, a twisting datum and a core in $Q'$ whose object is the
object of $Y$.
\formallabel{GlobalTheory.Mem.Top.RecAtG}
\end{definition}

We use the following statement (\lref{GlobalTheory.Mem.Top.LayOfCore.lay_of_core}): $\stageUEv_K$ of: for every core
$c$ over $K$ whose datum satisfies the conclusion of (VM), there is a layer of this datum with object $c.O$, lying over the
pre-character datum and level choice of $K$, satisfying the record and ledger conditions, in which every tuple is used.

\begin{lemma}[Layers from two passes]\label{pf:stageU:recatg}
Assume \textup{(VM)}. Let $\stageUKP$ be a base predicate and $Q,Q'$ slots. If $\mathrm{M}_0(\stageUKP)$, $\mathrm{M}_1(\stageUKP,Q)$,
$\mathrm{M}_2(\stageUKP,Q)$ and $\mathrm{M}_3(\stageUKP,Q,Q')$, then $\stageULay^{\mathrm g}$ of the property: $Y$ comes from a core
in $Q'$, and $\mathrm{QD}$ holds for $D$ and $Y$.
\formallabel{GlobalTheory.Mem.Top.GroupRetwisting.recAtG_of_two_passes}
\formallabel{GlobalTheory.T.layeredAutomorphicData_exists_of_twoPasses}
\end{lemma}

\begin{proof}
The statements $\mathrm{M}_1,\mathrm{M}_2,\mathrm{M}_3$ have the form $\stageUEv^{\mathrm g}_K\stageULg_K(\cdot)$, and so has, for
every context, the statement that $T$ is forced (\lref{GlobalTheory.Mem.Top.GroupRetwisting.lg_forced}: it suffices
that the prime set contains a prescribed finite set). The base statement $\mathrm{M}_0$(i) is an
$\stageUEv^{\mathrm g}$-statement and the layer statement quoted above is an $\stageUEv$-statement. Work at a good context
$K$ at which all of them hold, and let $(S_1,P_1,\mathcal E,S_2)$ be exceptional data serving simultaneously for
$\mathrm{M}_1$, $\mathrm{M}_2$, $\mathrm{M}_3$ and the forcing. By $\mathrm{M}_0$(i) there are a pre-base and a frame $f$ over $K$
and a base $B$ on this pre-base with $\stageUKP(B,f)$ whose coefficient fields contain $S_1$ and the numbers $S_2$ attached
to the character set $\{1,\delta,\delta'\}$ of $f$, which is the character set of untwisted data. Let $T$ be the
untwisted twisting datum over $B$ with prime set $P_1$; it satisfies the exceptional data (its auxiliary sets are empty)
and is forced. By $\mathrm{M}_1$ there is a core $c$ in $Q$ over $(B,f,T)$; by $\mathrm{M}_2$ it satisfies the export
condition; by $\mathrm{M}_3$, applied with empty exclusions and no numbers, there are an enlargement $B'$, the same frame,
a twisting datum $T'$ and a core $c'$ over $(B',f,T')$ which lies in $Q$ and in $Q'$ and is twisted at all places with good
twists. The datum of $c'$ has normalised volumes (its real record contains the normalisation), so it satisfies the
conclusion of (VM), and the layer statement gives a layer $Y$ with object $c'.O$, lying over the data of $K$, with the
record and ledger conditions and all tuples used. The twist conditions of $Y$ are those of $c'$, since the norm choice
and theta data of $Y$ are those of $c'.O$, and the stability of the lattices is part of the core $c'$. Hence $Y$ comes
from a core in $Q'$ and $\mathrm{QD}$ holds.
\end{proof}

\paragraph{Main terms.} The following statement is cited (\lref{GlobalTheory.Mem.Top.exists_du_nat}): for every
$\eps\in(0,\frac1{52})$ there is $x_0$ such that for every context datum at a scale with exponent $\eps$ and $x\ge x_0$ the
main terms $d_P(u)$ of \eqref{pf:stageU:eq:dD} are positive integers. (Its proof shows that the product of the mass
$\mathrm{mass}_k$ and the column indices is a positive integer, using $g_i\equiv1\pmod{Q_Sqd}$ and $g_i-1>2fk$.)

\begin{lemma}[Mass telescoping and the row main term]\label{pf:stageU:mainterms}
\begin{enumerate}[label=(\alph*)]
\item Assume \textup{(SM)}. For every CM field $E$ unramified over its maximal real subfield and all integers $k\ge2$,
$n\ge0$,
\[
\mathrm{mass}_{k+n}=\mathrm{mass}_k\prod_{M=k+1}^{k+n}\rho_E(M),\qquad
\rho_E(M)=\frac{\bigl((M-1)!\,(v/\pi)^M\bigr)^{f}}{(2v)^{f/2}}\,L(M,\eta^M),
\]
where $f=[E^+:\Q]$, $v$ is as in \cref{prop:families:mass} and $\rho_E(M)$ is the ratio given there.
\item For every pre-character datum $P$, $D_P>0$ and
\[
D_P\ \ge\ \bigl(\tfrac23\bigr)^f\,\mathrm{mass}_r\,\bigl(\iota_{\mathrm r}(g_0)\iota_{\mathrm r}(g_1)\bigr)^f
\prod_\tau\dim(v_\tau,-b_\tau),
\]
where $(v_\tau,-b_\tau)$ is the row parameter of the base word at $\tau$ and $\dim$ is its Weyl dimension.
\end{enumerate}
\formallabel{GlobalTheory.EX.WinLow.mass_tel}\formallabel{GlobalTheory.Mem.Top.DsLow.DsP_pos_and_lower}
\end{lemma}

\begin{proof}
(a) Induction on $n$: for $n=0$ the product is empty, and $\mathrm{mass}_{k+n+1}=\mathrm{mass}_{k+n}\,\rho_E(k+n+1)$ by (SM), since
$k+n+1\ge3$. (b) The mass $\mathrm{mass}_r$ is positive (\lref{Formalism.HermitianLatticesTheory.massTelescope_std_pos},
positivity of the telescoped mass of the standard lattice), the row indices are positive
(\lref{GlobalTheory.Mem.Top.DsLow.idxRowLoc_pos}), and the degree $d_\sigma(\tau)$ is the Weyl dimension of the row
parameter, which is positive (\lref{GlobalTheory.Mem.Top.DsLow.trueDsig_eq_dimV}); hence $D_P>0$, and the inequality
holds because $(\frac23)^f\le1$.
\end{proof}

\paragraph{A lower bound for the ratio of main terms.}
For a sequence $y$ indexed by $\N$ and $N\ge1$ write $V_N(y)=\prod_{0\le i<j<N}(y_i-y_j)$ and
$\dim_N(y)=V_N(y)/\prod_{i<N}i!$. Let $x\ge1$, $\Delta\ge1$, $B=e^{x^\eps}$ and $a=xB$. Sequences $(u,v,b,w)$ form a
\emph{place datum} of type $(k,m)$ if $0\le v_i\le4k\Delta$ and $v_i-v_{i+1}=4\Delta$ for $i<k-1$ (resp.\ $i+1<k-1$);
$v_i<u_i$ and $u_{i+1}<v_i$ for $i<k-1$; $-b_0<u_{k-1}$; either $u_0\le2a$ and $-a\le u_{k-1}$, or
$4k\Delta\le u_0-v_0\le8k\Delta$ and $u_{k-1}=v_{k-2}-\frac12$; $b_j<w_j$ and $|w_j-b_j|\le b_j/x$ for $j<m$; $w_j<b_{j+1}$
and $b_{j+1}-b_j\ge1$ for $j+1<m$; $4k\Delta\le b_0/x$; and $w_{m-1}\le2B$. Write $(v,-b)$ for the sequence
$v_0,\dots,v_{k-2},-b_0,-b_1,\dots$. The following estimate is cited
(\lref{Formalism.LogVandermondeBounds.vandermonde_place_lower_bound}, proved from
\lref{Formalism.LogVandermondeBounds.StandardLattice.logRatioLowerBound_of_logVandermondeGap} and
\lref{Formalism.LogVandermondeBounds.StandardLattice.VandPlace_holds}).

\smallskip\noindent\textbf{(LRB)} For every $\eps\in(0,\frac1{52})$ and $C_0\ge0$ there is $C_1$ such that the following
holds for all $x\ge1$, with $L=x^\eps$. Let $f,m\ge1$ and $2\le k\le x^{C_0}$ be integers, $v_E\ge\pi$ with $\log v_E\le L$,
$\Delta\ge1$, $g_1,g_2\ge1$, let $(u_\rho,v_\rho,b_\rho,w_\rho)$ be place data of type $(k,m)$ for $\rho<f$, and let
$\mathrm{mass}_k,\mathrm{mass}_r,I_\sigma,I_\pi>0$ and $L_M>0$ ($k<M<k+m$) with $|\log L_M|\le f\log\zeta(M)$ satisfy
\[
\mathrm{mass}_r=\mathrm{mass}_k\prod_{M=k+1}^{k+m-1}\frac{((M-1)!\,(v_E/\pi)^M)^f}{(2v_E)^{f/2}}L_M,\qquad
\frac{I_\sigma}{I_\pi}\ge2^{-2f}(g_1g_2)^{fm(2k+m-1)/2}.
\]
If $D_*\ge(\frac23)^f\mathrm{mass}_rI_\sigma\prod_\rho\dim_{k-1+m}(v_\rho,-b_\rho)$ and
$d_*=\mathrm{mass}_kI_\pi\prod_\rho\dim_k(u_\rho)$, then
\[
\sum_{j<m}\frac1f\sum_{\rho<f}\log\frac{w_\rho(j)v_E}{\pi}+m(\log g_1+\log g_2)-C_1\Bigl(L+\frac{m^2L}k\Bigr)\ \le\
\frac1{fk}\log\frac{D_*}{d_*}.
\]
Its proof compares Vandermonde products: for interlacing sequences as above $V_k(u)\le(8a)^{2k-3}V_{k-1}(v)$, since the
differences of the $u_i$ are at most $8a$, and $V_{k-1+m}(v,-b)\ge V_{k-1}(v)\prod_jb_j^{k-1}$, since $v_i+b_j\ge b_j$ and
$b_j-b_i\ge1$; the remaining terms are sums of logarithms of factorials and of the mass factors.

\begin{lemma}[Window lower bound]\label{pf:stageU:winlow}
Assume \textup{(SM)}. Let $\eps\in(0,\frac1{52})$, $C$ a standard constant package and $c_D\ge0$. There is $c_2$ such that for
every $c_h\ge c_2$ there is $x_0$ with the following property. For every scale $s$ with exponent $\eps$ and $s.x\ge x_0$, every
pre-character datum $P$ with weights and code satisfying the conditions relating them to $s$ (with the constants
$C,c_D,c_h$), all positive integers $d_u=d_P(u)$ and every real
$D_s\ge(\frac23)^f\mathrm{mass}_r(\iota_{\mathrm r}(g_0)\iota_{\mathrm r}(g_1))^f\prod_\tau\dim(v_\tau,-b_\tau)$,
\[
m\,L\ \le\ \log\Bigl(\Bigl(\frac{D_s}{d_u}\Bigr)^{\frac1{fk}}\Gamma(C,s,R)^{-\frac2{fk}}\Bigr)\qquad\text{for every word }u.
\]
\formallabel{GlobalTheory.EX.WinLow.window_lower}
\end{lemma}

\begin{proof}
Let $C_1$ be the constant of (LRB) for $\eps$ and $C_0=5$. For $x$ large the following hold, by results cited with their
content.
\begin{enumerate}[label=(\roman*)]
\item At every real place the Harish-Chandra parameters of every word form a place datum of type $(k,m)$ with
$\Delta=\lceil x^{c_D}\rceil$, and $w_\rho(j)\ge1$ (\lref{GlobalTheory.EX.WinLow.places_ok}).
\item $k\le x^5$ (\lref{GlobalTheory.EX.WinLow.scale_sizes}), and $2|C_1|\le m$, $m^2\le k$, $2C_\Gamma m^4\le k$ and
$L\ge1$ (\lref{GlobalTheory.EX.WinLow.rank_consts}).
\item $\pi\le v_E$ (\lref{GlobalTheory.EX.WinLow.pi_le_v}), $\log v_E\le L$ and $g_i\ge b^8$, so $\log g_i\ge8L$ and $g_i\ge3$
(\lref{GlobalTheory.EX.hostBounds}, the size bounds for the host).
\item The column and row indices are positive and $I_\sigma/I_\pi\ge2^{-2f}(g_0g_1)^{fm(2k+m-1)/2}$ for
$I_\sigma=(\iota_{\mathrm r}(g_0)\iota_{\mathrm r}(g_1))^f$ and $I_\pi=(\iota_{\mathrm c}(g_0)\iota_{\mathrm c}(g_1))^f$
(\lref{GlobalTheory.EX.WinLow.idx_ratio}, from the orders of the finite general linear groups).
\item $\mathrm{mass}_r=\mathrm{mass}_k\prod_{M=k+1}^{r}\rho_E(M)$ with $r=k+m-1$ (\cref{pf:stageU:mainterms}(a)), and
$L(M,\eta^M)>0$, $|\log L(M,\eta^M)|\le f\log\zeta(M)$ for $M\ge2$
(\lref{Formalism.HermitianLattices.GlobalBoundRefinements.log_l_value_field_bound} and
\lref{Formalism.MainNotions.StandardLattice.cmField_unramified_massIdentity}, bounds for these $L$-values).
\end{enumerate}
The main term $d_u$ is the $d_*$ of (LRB) for the parameters of $u$, and the hypothesis on $D_s$ is that of (LRB). Hence
\[
\frac1{fk}\log\frac{D_s}{d_u}\ \ge\ \sum_{j<m}\frac1f\sum_\rho\log\frac{w_\rho(j)v_E}{\pi}+m(\log g_0+\log g_1)
-C_1\Bigl(L+\frac{m^2L}k\Bigr).
\]
The first sum is non-negative since $w_\rho(j)\ge1$ and $v_E\ge\pi$; $m(\log g_0+\log g_1)\ge16mL$; and
$C_1(L+m^2L/k)\le|C_1|\cdot2L\le mL$. Hence $\frac1{fk}\log(D_s/d_u)\ge15mL$ (in particular $d_u>0$). Finally
$\frac2{fk}\log\Gamma(C,s,R)=\frac{2C_\Gamma m^4}{k}\,mL\le mL$, so the logarithm in the assertion is at least $14mL\ge mL$.
\end{proof}

\begin{proposition}[Layered data from layers]\label{pf:stageU:recexp}
Assume \textup{(SM)}. Let $Q'$ be a slot and $\stageUQO$ a family of properties of data and objects indexed by parameter
tuples. If $\stageULay^{\mathrm g}$ holds for the property ``$Y$ comes from a core in $Q'$ and $\mathrm{QD}$ holds for $D$ and $Y$'',
and $\stageUEv_{\kappa,K}$(every bundle point $\beta$ over $K$ in $Q'$ satisfies $\stageUQO(\kappa)(\beta.D,\beta.O)$), then
$\stageULay\bigl(\kappa\mapsto(\kappa\text{ positive}\Rightarrow\stageUQO(\kappa))\wedge\mathrm{QD}\bigr)$.
\formallabel{GlobalTheory.Fc.recExP_of_recAtG}\formallabel{GlobalTheory.T.existsLayeredData_of_goodScaleContexts}
\end{proposition}

\begin{proof}
Both hypotheses are eventual statements in $(\eps,C_8,C_5,c_6,C_b)$, $C$, $c_D$, $c_h$ and $x$; combine them with the
thresholds of \cref{pf:stageU:winlow} (for $c_D\ge0$) and of the integrality of the main terms. Thus, given the
constants, we obtain $c_D,c_h$ and $x_0$ such that the following holds for every scale $s$ with exponent $\eps$ and
$s.x\ge x_0$; we take $c_s=1$. Let $P,a,b,c$ satisfy the goodness conditions with $\kappa$ (which include the conditions
relating them to $s$).
\begin{enumerate}[label=\arabic*.]
\item The main terms $d_u=d_P(u)$ are positive integers, and $D_s:=D_P$ is positive and satisfies the lower bound of
\cref{pf:stageU:mainterms}(b). By \cref{pf:stageU:winlow}, \eqref{pf:stageU:eq:recexp} holds with $c_s=1$.
\item Let a level choice with $d^*=d$, $D^*=D_s$ and $\Gamma=\Gamma(C,s,R)$ be given. Together with $s,C,c_D,c_h,P,a,b,c$ it
forms a context $K$ with $\kappa\triangleright K$. The hypothesis $\stageULay^{\mathrm g}$ at $K$ gives a datum $D$ and a layer $Y$
over $(P,\text{level choice})$ with the record, ledger and twist conditions, coming from a core $c'$ in $Q'$, with
$\mathrm{QD}$.
\item Let $p$ be a ledger part of $Y$ at the scale $s$, with weights $(a,b)$, valued and canonical. If $\kappa$ is positive,
then $c'$, the layer $Y$ and the ledger part $p$, together with the record and ledger conditions, the used tuples and the
twist conditions, form a bundle point $\beta$ over $K$ in $Q'$ whose datum and object are those of $p$; by the second
hypothesis, $\stageUQO(\kappa)$ holds at $p$. Moreover the datum and object of $p$ arise from $D$ and the object of $Y$ by
replacing the division lattices and the height function, and $\mathrm{QD}$ is preserved by this replacement, because
the used tuples, the twist goodness and the lattices $\widetilde M(\mu)$ only involve components which are not changed.
Hence $\mathrm{QD}$ holds at $p$.
\end{enumerate}
\end{proof}

\begin{proposition}[Layered data from member existence]\label{pf:stageU:membertop}
Assume \textup{(SM)} and \textup{(VR)}. Let $\stageUKP$ be a base predicate, $Q,Q'$ slots and $\stageUQO$ a family of properties
of data and objects indexed by parameter tuples. If $\mathrm{M}_0(\stageUKP)$, $\mathrm{M}_1(\stageUKP,Q)$, $\mathrm{M}_2(\stageUKP,Q)$,
$\mathrm{M}_3(\stageUKP,Q,Q')$ and $\stageUEv_{\kappa,K}$(every bundle point $\beta$ over $K$ in $Q'$ satisfies
$\stageUQO(\kappa)(\beta.D,\beta.O)$), then $\stageULay\bigl(\kappa\mapsto(\kappa\text{ positive}\Rightarrow\stageUQO(\kappa))\wedge\mathrm{QD}\bigr)$.
\formallabel{GlobalTheory.Fc.memberTop_final_contextParameters}
\end{proposition}

\begin{proof}
By \cref{pf:stageU:volmass}, (VR) implies (VM). \Cref{pf:stageU:recatg} gives $\stageULay^{\mathrm g}$ for the property ``$Y$ comes
from a core in $Q'$ and $\mathrm{QD}$ holds'', and \cref{pf:stageU:recexp} gives the conclusion.
\end{proof}

% ---------------------------------------------------------------------------
\subsection{The size row}\label{sec:pf:stageU:size}

For a family $\stageUFm$ write $\sfname{TuSize}(\stageUFm)$ for the statement: eventually in $x$, every member satisfies
$b^{m/2}\le T(u)$ and $\log T(u)\le\stageUFm.C_{Tu}\,m\,L$ for every word $u$, where $T(u)$ is the level prime of $u$
(\cref{def:families:closed}). Put $c_T^{\min}(y)=\max\bigl(11,\,16+3\max(5,y+1)\bigr)$. Recall that $\mathrm{mass}_n$ is the
mass of the standard hermitian lattice $\OO_E^n$. Write (GM) for the statement: for every datum $D$ and every positive
definite hermitian space $W$ of rank two over $E$, the mass of the distinguished lattice of $W$ equals $\mathrm{mass}_2$.

\begin{lemma}[Level primes are small]\label{pf:stageU:tusizemin}
Let $\stageUFm$ be a family with $c_T^{\min}(\stageUFm.\mathrm{CYL})\le\stageUFm.C_{Tu}$ such that eventually every member $D$
satisfies: the volume of the base space of $D$ (of rank $r$) equals $\mathrm{mass}_r$, and the degrees $d_\sigma(\tau)$
satisfy the two-sided bounds of \cref{prop:families:famgl} in terms of the gaps. Then $\sfname{TuSize}(\stageUFm)$.
\formallabel{Formalism.MainNotions.StandardLattice.HarishChandraReduction.tuSize_min}
\end{lemma}

\begin{proof}
Put $C_Y=\max(5,\mathrm{CYL}+1)$. By \lref{Formalism.MainNotions.StandardLattice.HarishChandraReduction.explicitConstants_family_of_volume_eq_mass}
(explicit count bounds for the Vandermonde products of the members, deduced from the mass ratio, the base volume and the
degree bounds through the level-prime family of the division lattices), eventually every member satisfies the explicit
Vandermonde count bounds with constants $(C_Y,C_Y)$. By \lref{Formalism.MainNotions.GlobalTheoryTheory.rowsAF_fam}
(the row bounds follow from the data bounds), eventually at every member the data bounds with constant $C_Y$ imply the
row bounds with constant $C_Y$. These two statements combine
(\lref{Formalism.MainNotions.StandardLattice.part_five_all_of_explicit}, which requires $11\le c$ and $16+3C_Y\le c$
for the final constant $c=c_T^{\min}(\mathrm{CYL})$) into the conclusion of the Harish-Chandra reduction with constant
$c_T^{\min}(\mathrm{CYL})$: explicit inequalities for the level primes and the Vandermonde variations of the members.
Finally \lref{Formalism.MainNotions.GlobalTheoryTheory.tuSize_fam} gives $\sfname{TuSize}(\stageUFm)$ from these
inequalities, since $c_T^{\min}(\mathrm{CYL})\le C_{Tu}$.
\end{proof}

\begin{proposition}[The size row]\label{pf:stageU:hTu}
Assume \textup{(VR)} and \textup{(GM)}. For every slot $Q$ and every linkage $\stageULk$ for the bundle conditions,
\[
\sfname{Row}\Bigl(Q,\stageULk,\ \stageUA\mapsto\sfname{Closed}\bigl(W'_1(\stageUA),\ \stageUFm\mapsto
(c_T^{\min}(\stageUFm.\mathrm{CYL})\le\stageUFm.C_{Tu}\Rightarrow\sfname{TuSize}(\stageUFm))\bigr)\Bigr).
\]
\formallabel{GlobalTheory.Fc.row_hTu_cTmin_alt}\formallabel{GlobalTheory.T.hTu_cTmin_of_genusMass}
\end{proposition}

\begin{proof}
A property of classes with selections which holds for every $\mathrm{True}$-selection is a row
(\lref{GlobalTheory.Fc.Row.of_sel}, which takes $E=\mathrm{True}$), and a property of families which holds for every
family of a class is closed on that class (\lref{Formalism.MainNotions.Pk.Closed.of_class}). So let $\sigma$ be a
$\mathrm{True}$-selection from $(Q,\stageULk)$ for a class $\stageUA$, and $\stageUFm\in W'_1(\stageUA)$ with
$c_T^{\min}(\mathrm{CYL})\le C_{Tu}$. Since $W'_1(\stageUA)\subseteq\stageUA$ and $\sigma$ is a selection, there is
$\kappa\sim\stageUFm$ such that every selected bundle $\sigma(\stageUFm,x,i)$ is built from $Q$ along $\stageULk$ at $\kappa$.
A built bundle contains a real record of its datum (\lref{GlobalTheory.Fc.built_val}, which lists the data carried
by built bundles), and a real record normalises the volumes of the datum
(\lref{Formalism.MainNotions.RT.volTie_of_realTied}). So at every member $D$ the volumes are normalised, and (VR)
gives the volume ratio formula at $D$. By the global arithmetic identity, this formula says that $\vol(W)/\vol(W')$ is the
right-hand side of the mass ratio of \cref{prop:families:mass}
(\lref{Formalism.MainNotions.StandardLattice.volumeMass_of_unfolding}, which states this reformulation). By (GM) and
the decomposition of the volume of a definite space over the classes of its genus
(\lref{Formalism.L2.GlobalTheoryTheory.GenusMassVolumes.volume_mass_of_genus_rank}), $\vol(W)=\mathrm{mass}_2$ for every
positive definite $W$ of rank two. Climbing from rank two to the rank $r$ of the base space with the volume ratios
(\lref{Formalism.L2.GlobalTheoryTheory.GenusMassVolumes.volume_eq_mass_rank_two}) gives that the volume of the base
space equals $\mathrm{mass}_r$. The two-sided degree bounds hold eventually at the members of every family
(\lref{Formalism.MainNotions.GappedMultiplicityBounds.explicit_bound_family}, a consequence of multiplicity one of the
lowest $K$-type in large gap). \Cref{pf:stageU:tusizemin} gives $\sfname{TuSize}(\stageUFm)$.
\end{proof}

% ---------------------------------------------------------------------------
\subsection{Uniform statements}\label{sec:pf:stageU:uniform}

\paragraph{The thin-scale package.}
$\sfname{LitThin}$ (\cref{prop:families:design}) is the conjunction of the following three statements about Dirichlet
$L$-functions. (Formally each is attached to a datum $D$, which does not occur in it.) Here $m_\rho$ denotes the order of a
zero $\rho$.
\begin{itemize}
\item[(Z)] There are $C_1,c>0$ such that
\begin{enumerate}[label=(Z\arabic*)]
\item for every $Q\ge3$, the number of zeros $\rho$ with $\Re\rho\ge1-\frac1{400}$ and $|\Im\rho|\le Q$ of the functions $L(s,\chi)$,
$\chi$ a primitive character of conductor $\le Q$, counted with multiplicity and summed over $\chi$, is at most
$C_1Q^{1/4}$;
\item for every primitive character $\chi$ of conductor $N\ge3$ and every zero $\rho$ of $L(s,\chi)$ with
$\Re\rho\ge1-c/\log\max(N,N|\Im\rho|)$: $\chi$ is real, $\rho$ is real and simple, and $\rho$ is the only zero of $L(s,\chi)$ in
the part of the formal proof treated here;
\item every zero $\rho$ of $\zeta$ with $0<\Re\rho<1$ satisfies $\Re\rho\le1-c/\log(|\Im\rho|+3)$.
\end{enumerate}
\item[(EF)] There is $C>0$ such that for every primitive $\chi$ of conductor $N$ and all $2\le T\le u$ the set of zeros with
$0<\Re\rho<1$, $|\Im\rho|<T$ is finite and
\[
\Bigl|\psi_0(u,\chi)+\sum_{|\Im\rho|<T}m_\rho\frac{u^\rho}\rho\Bigr|\le C\Bigl(\frac{u\log^2(NuT)}T+\log(Nu)
+\sum_{|\Im\rho|<1}\frac{m_\rho}{|\rho|}\Bigr),
\]
where $\psi_0(u,\chi)$ is the normalised Chebyshev function of $\chi$ and the sums run over zeros with $0<\Re\rho<1$.
\item[(ZC)] There is $C>0$ such that for every primitive $\chi$ of conductor $N$ and every $t\in\R$ the zeros with
$0<\Re\rho<1$ and $|\Im\rho-t|\le1$ are finite in number and $\sum m_\rho\le C\log(N(|t|+2))$.
\end{itemize}

\begin{littheorem}[Zero density near $\Re s=1$ {\cite[Ch.~12]{Montgomery1971}, \cite{Gallagher1970}}]\label{pf:stageU:zd}
Statement \textup{(Z1)} holds.
\formallabel{GlobalTheory.ZeroDensityMult.zero_density_mult}
\end{littheorem}

This is a weak form of the large-sieve density estimate
$\sum_{q\le Q}\sideset{}{^*}\sum_{\chi}N(\alpha,T,\chi)\ll(Q^2T)^{3(1-\alpha)/(2-\alpha)}\log^9(QT)$ of
\cite[Ch.~12]{Montgomery1971}: with $T=Q$ and $\alpha=1-\frac1{400}$ the exponent is $3\cdot\frac{3(1-\alpha)}{2-\alpha}=\frac9{401}<\frac14$.
The formal proof follows Montgomery's method: a mollified Dirichlet polynomial over $(Q^3,Q^6]$ has absolute value
$\ge\frac12$ at every zero in the box (zero detection); the hybrid large sieve, deduced from the multiplicative large sieve
and Gallagher's lemma \cite{Gallagher1970}, bounds the number of unit cells $\{(\chi,\lfloor\Im\rho\rfloor)\}$ containing zeros by
$CQ^{1/8}$; and each cell contains $\ll\log Q\ll Q^{1/8}$ zeros by (ZC)
(\lref{GlobalTheory.T.ncard_unitCells_dirichletLZeros_near_one_le} is the cell count).

\begin{lemma}[Landau--Page region]\label{pf:stageU:lp}
There is $c>0$ such that for every modulus $q$, every Dirichlet character $\chi\bmod q$ and every zero $s$ of $L(s,\chi)$
(with $s\ne1$ if $\chi$ is principal) in the region $\Re s\ge1-c/\log(q(|\Im s|+2))$: $\chi^2=1$, $\chi\ne1$, $s$ is real and
$L'(s,\chi)\ne0$; and for every character $\chi'\bmod q$ and every zero $s'$ of $L(s,\chi')$ in the part of the formal proof treated here, $\chi'=\chi$ and
$s'=s$.
\formallabel{GlobalTheory.LandauPageRegion.main}
\end{lemma}

\begin{proof}
We use three classical statements \cite[Ch.~14]{Davenport2000}:
(A) \lref{GlobalTheory.ZeroFreeComplex.zero_free_region}: there is $c_1>0$ such that if $\chi\ne1$, $L(s,\chi)=0$ and
$\chi^2\ne1$ or $\Im s\ne0$, then $\Re s<1-c_1/\log(q(|\Im s|+2))$;
(B) \lref{GlobalTheory.RealZeroNearOne.main}: there is $c_2>0$ such that if $\chi,\psi\bmod q$ are non-principal real
characters and $\beta,\beta'\ge1-c_2/\log2q$ are real zeros of $L(s,\chi)$ and $L(s,\psi)$, then $\chi=\psi$, $\beta=\beta'$ and
$L'(\beta,\chi)\ne0$;
(C) \lref{GlobalTheory.ZeroFreePrincipal.zero_free}: there is $c_3>0$ such that every zero $s\ne1$ of $L(s,\chi_0)$, $\chi_0$
principal modulo $q$, satisfies $\Re s<1-c_3/\log(q(|\Im s|+2))$.
Let $c=\min(c_1,c_2,c_3)$; the region grows with $c$. If $s$ is a zero of $L(s,\chi)$ in the region, then $\chi\ne1$ by (C),
and $\chi^2=1$ and $\Im s=0$ by (A); hence $s=\beta$ is real and $\beta\ge1-c_2/\log2q$, and (B) with $\psi=\chi$ gives
$L'(\beta,\chi)\ne0$. If $s'$ is a zero of $L(s,\chi')$ in the region, the same argument shows that $\chi'$ is real and
non-principal and $s'=\beta'\ge1-c_2/\log2q$ is real, and (B) gives $\chi'=\chi$ and $\beta'=\beta$.
\end{proof}

\begin{lemma}[Zero-free regions in the form used]\label{pf:stageU:zfhalf}
If $c$ is as in \cref{pf:stageU:lp}, then \textup{(Z2)} and \textup{(Z3)} hold with the constant $c/2$.
\formallabel{GlobalTheory.ZeroFreeHalf.main}
\end{lemma}

\begin{proof}
(Z2) A primitive character $\chi$ of conductor $N\ge3$ is not principal. For $t\in\R$ we have $N(|t|+2)\le\max(N,N|t|)^2$:
if $|t|\le1$ then $N(|t|+2)\le3N\le N^2$, and if $|t|>1$ then $N(|t|+2)\le3N|t|\le(N|t|)^2$. Hence
$\log(N(|t|+2))\le2\log\max(N,N|t|)$, and a zero $\rho$ with $\Re\rho\ge1-\frac c2/\log\max(N,N|\Im\rho|)$ lies in the region of
\cref{pf:stageU:lp}. Thus $\chi^2=1$, so $\chi$ is real, $\rho$ is real, $L'(\rho,\chi)\ne0$, so the zero is simple, and every
other zero in the region of (Z2) lies in the region of the lemma and equals $\rho$. (Z3) Let $\rho$ be a zero of $\zeta$ with
$0<\Re\rho<1$, and suppose $\Re\rho>1-\frac c2/\log(|\Im\rho|+3)$. Since $\log(|t|+2)\le2\log(|t|+3)$, $\rho$ lies in the region of
\cref{pf:stageU:lp} for $q=1$; as $\zeta(s)=L(s,\chi_0)$ for the character $\chi_0$ modulo $1$ and $\rho\ne1$, the lemma gives
$\chi_0\ne1$, a contradiction.
\end{proof}

\begin{theorem}[The thin-scale package]\label{pf:stageU:litthin}
$\sfname{LitThin}$ holds.
\formallabel{GlobalTheory.ThinScaleLiterature.litThin}\formallabel{GlobalTheory.ZerosNearOne.main}
\formallabel{GlobalTheory.ZerosNearOne.density_and_regions}\formallabel{GlobalTheory.ZeroBounds.explicit_formula_and_window_count}
\end{theorem}

\begin{proof}
(Z1) is \cref{pf:stageU:zd}, and (Z2), (Z3) follow from \cref{pf:stageU:lp,pf:stageU:zfhalf}. (EF) is the truncated explicit
formula \lref{GlobalTheory.ExplicitFormula.truncated_explicit_formula} \cite[Ch.~17, 19]{Davenport2000}, and (ZC) is the local
zero count \lref{GlobalTheory.ZeroWindow.count_zeros_window} \cite[Ch.~16]{Davenport2000}. Since the three statements do not
involve the datum, they hold for every datum.
\end{proof}

\paragraph{The volume ratio.}
Write (VR) for the statement: for every datum $D$ with normalised volumes, every hermitian space $W$ over $E$ of rank $n\ge3$
with signature $(P,Q)$ at every complex embedding of $E$, and every sign $\pm$ with $P\ge1$ (for $+$) or $Q\ge1$ (for $-$),
there is $x\in W$ with $\langle x,x\rangle=\pm1$ such that $x^\perp$ is the image of an isometric embedding of a hermitian
space $W'$ of signature $(P-1,Q)$ (resp.\ $(P,Q-1)$) with
\[
\frac{\vol(W')}{\vol(W)}=\Bigl(\frac{(2\pi)^n}{(n-1)!}\Bigr)^{f}\,(2v)^{f(\frac12-n)}\,L(n,\eta^n)^{-1},
\]
where $f=[F:\Q]$, $v$ is the invariant of $E$ in \cref{prop:families:mass} and $L(n,\eta^n)$ is the $L$-value there.

\begin{proposition}[The volume ratio]\label{pf:stageU:volratio}
The Siegel--Weil formula in the convergent range (\cref{lit:families:sw}, in the form of the field of $P$ asserting finite
volume and the Siegel density identity for every CM extension and every non-degenerate hermitian matrix) implies
\textup{(VR)}.
\formallabel{GlobalTheory.VolRatioFromSiegelWeil.volumeRatio}
\end{proposition}

\begin{proof}
The proof has three steps, of which the last is cited.
\begin{enumerate}[label=\arabic*.]
\item Specialising the formula to one vector ($n=1$, $T=(\beta)$) gives, for $M\ge3$, $\beta\in F^\times$ and
$\Phi=\Phi_\infty\otimes\Phi_f$ with $\Phi_\infty$ smooth and compactly supported and $\Phi_f$ a pure finite test function:
the mean over $H(F)\backslash H(\A_F)$ of $h\mapsto\sum_{\langle x,x\rangle=\beta}\Phi(h^{-1}x)$ equals
$\Phi^\sharp_\infty(\beta)\prod_v\Phi^\sharp_v(\beta)$, where $\Phi^\sharp_\infty,\Phi^\sharp_v$ are the continuous densities of the
push-forwards of $\Phi_\infty,\Phi_v$ under $x\mapsto\langle x,x\rangle$ for self-dual measures (\cref{pf:stageU:onevec}).
\item By approximation, the same identity holds for non-negative measurable $\Phi_\infty$ and non-negative pure $\Phi_f$, the
densities being replaced by integrals against fibre measures (disintegrations of the Haar measures along the hermitian
norm) and the sum running over $x\ne0$ (\lref{GlobalTheory.OneVecMeas.sw1b2_of_cc}, with its statement form
\lref{GlobalTheory.T.siegel_one_vector_measurable_of_compactSupport}).
\item This one-vector mean-value formula implies (VR) (\lref{GlobalTheory.VolRatioFromSecondClause.volumeRatio}, by
unfolding the mean of a Gaussian over the orbit of a vector of norm $\pm1$, as in \cref{prop:families:volratio}).
\end{enumerate}
\end{proof}

\paragraph{The main-term bound for equal words.}
Write (MT) for the statement of the main-term field $\mathsf{sM}$ of the package: for all constants
$C_g,c_b,C_w,c_T,C_T,C_0,A$ with $c_b,c_T>0$ there are $C_B\ge0$ and $C'_0,C_{l,0}$ such that for all $C'\ge C'_0$ and
$C_l\ge C_{l,0}$, every datum $D$ with size data, every Hida setting of $D$ with full data and every word $u$ the
following holds: if $D$ satisfies its ties, the pair $(u,u)$ satisfies the counting hypotheses with the constants
$(C',C_l,C_g,c_b,C_w,c_T,C_T,C_0,A)$, the concentration tie holds and the diagonal operator of $u$ converges absolutely,
then the sum of the generic terms of the orbital expansion of the diagonal operator is at most
$\gamma_B\,d_u^2/D_*$. Here $\gamma_B=\exp(C_B\,m_R^5\,f\log b)$ with $m_R=2r_S+2$ determined by the ranks of $D$, and
$d_u$, $D_*$ are the dimension and the main term of the Hida setting (\cref{thm:theorem-S}).

\begin{proposition}[The main-term bound]\label{pf:stageU:mainrow}
The Siegel--Weil formula in the convergent range (in the form of \cref{pf:stageU:volratio}) implies \textup{(MT)}.
\formallabel{GlobalTheory.rowMainEqP_of_siegelWeil}\formallabel{GlobalTheory.rowMainEqP_of_siegelWeil_of_stepChoice}
\end{proposition}

\begin{proof}
The proof combines eleven step statements by a cited combination theorem.
\begin{enumerate}[label=\arabic*.]
\item The Siegel--Weil formula gives the step-choice constants
(\lref{GlobalTheory.St.stepChoiceDataC_of_siegelWeil_inInfinity}), given the convergence at infinity of the energy
integrals for real objects (\lref{Formalism.L2.SG.RealObj.EnergyConvJ.BM.realObj_inInfinityX_jsz}, applied with the
uniqueness statement \lref{CompanionTheory.Lime.BM.uniqGap_holds} for large gap).
\item The Siegel--Weil formula gives the volume identities $\sfname{SWXS}$ of \cref{prop:families:volratio}(ii)
(\lref{GlobalTheory.St.swxs_of_siegelWeil}), and these give the generic family step
(\lref{GlobalTheory.T.genericFamilyStep_of_siegelWeil}).
\item The chain of step estimates for equal words (\lref{GlobalTheory.T.equalWords_nonGeneric_estimates_chain}) gives
the non-generic chain, the class, point, norm-one, tie, pin and vanishing steps, and, from the generic family step, the
step of local generation for pinned data; one of its nine components is not used. The normal step is the normal count
for generic rows (\lref{GlobalTheory.T.normalCount_genericRow}).
\item Theorem B in the pinned form required here follows from the global Theorem B
(\lref{Formalism.MainStatements.RowRest.bDat_theoremB_global}), which has the same conclusion under one hypothesis
fewer.
\end{enumerate}
The combination theorem \lref{GlobalTheory.T.equalWordGenericBound_of_normalisedSteps} (proved in
\lref{GlobalTheory.SnPos.rowMainEqP_of_pinned} by a cut bound, a height bound from local estimates and an analytic
bound from the choice step) turns these eleven statements into (MT).
\end{proof}

\begin{lemma}[Finite dimensionality and the genus mass in rank two]\label{pf:stageU:findim}
\begin{enumerate}[label=(\alph*)]
\item For every automorphic instance of a unitary group, the space of automorphic forms cut out by a locally constant
level function supported in the level subgroup and by the matrix coefficients, at the infinite places, of irreducible
unitary representations with formal degree and integrable matrix coefficients, is finite-dimensional.
\item \textup{(GM)} holds.
\end{enumerate}
\formallabel{GlobalTheory.rowsFinDim_link}\formallabel{GlobalTheory.genusMassAt_two}
\formallabel{Formalism.L2.GlobalTheoryTheory.GenusMassVolumes.genusMassAt_of}
\end{lemma}

\begin{proof}
(a) is \lref{CompanionTheory.D.gk.RF.automorphic_forms_finite_dimensional}, which states exactly this (compare
\cite{HarishChandra1968}). (b) By \lref{GlobalTheory.genusMassStd_all}, for every datum $D$ and every positive definite
hermitian space $W$ over $E$, the mass of the distinguished lattice of $W$ equals $\mathrm{mass}_n$, $n$ the rank of $W$;
specialise to $n=2$.
\end{proof}

% ---------------------------------------------------------------------------
\subsection{Assembly}\label{sec:pf:stageU:assembly}

\paragraph{Three further uniform statements.}
Besides $\sfname{LitThin}$, (VR), (GM) and (MT), the package contains the following uniform statements, which are proved
elsewhere: (Wu) for every datum whose rank $r$ is at least $1801$ and whose Harish-Chandra parameters have all gaps at least
$16r^{19}$, and every Hida setting $X$ of it, Wu's decay estimate holds for $X$ and the archimedean test function of $X$ is
integrable (\lref{GlobalTheory.wuDecayOfGap}, which proves it); (TY) the type bounds: for all constants
$C_g,c_b,C_w,c_T,C_T,C_0$ with $c_b,c_T>0$ there are $C_H,C_h,C_p\ge0$ such that for $C_s\ge C_p+4$, all large $C$ and all
large $C',C_l$, every word of every datum with size data, for every Hida setting with full data satisfying the ties, the
counting hypotheses with further parameter conditions, and the concentration tie, satisfies the type bound with the
constants $(C_H,C_s,C,C_h,C_p)$ (\lref{GlobalTheory.rowsPos_typeP}, which proves it); and (NG) the
bound $\varepsilon_{\mathrm{NG}}\,d_ud_{u'}/D_*$ for the non-generic terms of \cref{prop:families:volratio}(iii)
(\lref{GlobalTheory.rowNGP_of_siegelWeil}, which deduces it from the Siegel--Weil formula). We write (RF) for
\cref{pf:stageU:findim}(a).

\paragraph{The fields of the package.}
An element of $\sfname{HypsAlt}(\stageUNN,\stageUDR)$ (\cref{def:hyps}, \cref{tab:families:fields}) consists of a class
$\stageUA$, a raw selection $\sigma$ for $\stageUA$, a function $c_T:\R\times\R\to\R$ and $43$ fields of five kinds:
\begin{enumerate}[label=(F\arabic*)]
\item $\mathsf{hEx}$: $\sfname{MinkDiv}\Rightarrow\sfname{ExS}(\stageUA,c_T)$; $\mathsf{hduty}$: every family in $\stageUA$
fulfils the duties of the package; $\mathsf{hMt}$: for every family in $\stageUA$, eventually the lattices
$\widetilde M(\mu)$ of the members are stable under $\OO_{\mathit{LL}}$;
\item eight properties ($\mathsf{hRB}$, $\mathsf{hTw}$, $\mathsf{hPA}$, $\mathsf{hLamP}$, $\mathsf{hDivG}$, $\mathsf{hDcPtG}$,
$\mathsf{hInstR}$, $\mathsf{selJ}$) of the selected bundles or of the members, required for every $\stageUFm\in\stageUA$,
every $x>x_1$ which is at least the threshold of the selection and every $i\in I_x$; the last one asks for a datum (a join
of the indefinite realisation), the others for propositions;
\item the closedness of properties on $W'_1(\stageUA)$ ($\mathsf{hTu}$, $\mathsf{hn4}$) and on $D3(3,\stageUA)$ (fourteen
fields $\mathsf{hTF}$, $\mathsf{sO}$, $\mathsf{h30}$, $\mathsf{jOn}$, $\mathsf{aPt0}$, $\mathsf{aPtTw}$, $\mathsf{aPer0}$,
$\mathsf{aPerTw}$, $\mathsf{aP10}$, $\mathsf{aAtG}$, $\mathsf{aC2}$, $\mathsf{irrC}$, $\mathsf{zL}$, $\mathsf{tie}$, some of
which involve the selection), two statements about the selection on $D3(3,\stageUA)$ ($\mathsf{sLp}$, $\mathsf{hServ}$), the
definite rows along the selection ($\mathsf{rRows}$, with the definite-side condition $\stageUDR$), and
$\stageUNN(D3(3,\stageUA))$ ($\mathsf{nRest}$); here $\mathsf{hTu}$ is
$\sfname{Closed}\bigl(W'_1(\stageUA),\ \stageUFm\mapsto(c_T(\eps,\stageUFm.\mathrm{CYL})\le\stageUFm.C_{Tu}\Rightarrow
\sfname{TuSize}(\stageUFm))\bigr)$;
\item constants: a function $c_{30}$ from lists of reals to counting constants such that, for every family, the
counting constants of its output list satisfy $0\le C$ and $2(\mathrm{CYL}+1)\le A$ ($\mathsf{c30}$, $\mathsf{hc30}$), and a
real $c_{Ic}>0$ ($\mathsf{cIc}$, $\mathsf{hcI}$); the fields $\mathsf{h30}$ and $\mathsf{irrC}$ involve these constants;
\item eight uniform statements: $\sfname{LitThin}$, (Wu), (RF), (TY), (NG), (MT), (VR) and (GM).
\end{enumerate}
The row package $\sfname{AllRowsAlt}(\stageUKP,Q,Q',\stageULk,\stageUNN,\stageUDR)$ (\cref{prop:families:allrows}) is the
conjunction of twenty-six linked statements. Twenty-five of them are linked rows $\stageULin_{\stageUKP}(Q,Q';\stageULk,H)$,
one for each field of (F2) and (F3) other than $\mathsf{hTu}$, $\mathsf{h30}$ and $\mathsf{irrC}$, whose property $H$ is the
property of the field with the following differences: for the fields of (F2) it is required at all selected bundles for all
$x>x_1$ (for $\mathsf{selJ}$: there exists a join); for $\mathsf{hn4}$ the closedness is asserted on $W_1(\stageUA)$; and
for the nine fields $\mathsf{jOn}$, $\mathsf{aPt0}$, $\mathsf{aPtTw}$, $\mathsf{aPer0}$, $\mathsf{aPerTw}$, $\mathsf{aP10}$,
$\mathsf{aAtG}$, $\mathsf{aC2}$, $\mathsf{tie}$ it is asserted on $W_3(\stageUA)$ (\cref{sec:pf:stageU:calculus}). The
twenty-sixth, the \emph{counting statement}, is $\stageULin_{\stageUKP}(Q,Q';r_{30})$, where $r_{30}(S)$ says: there are
$c_{30}$ and $c_{Ic}$ as in (F4) such that $\sfname{Row}(S,\stageULk,H_{30})$, $H_{30}$ being the conjunction of the
properties of $\mathsf{h30}$ and $\mathsf{irrC}$ for these constants.

\begin{theorem}[The package from the rows]\label{pf:stageU:allrows}
Assume \textup{(SM)}, $\sfname{LitThin}$, \textup{(Wu)}, \textup{(RF)}, \textup{(TY)}, \textup{(NG)}, \textup{(MT)},
\textup{(VR)} and \textup{(GM)}. Let $\stageUKP$ be a base predicate, $Q,Q'$ slots, $\stageULk$ a linkage for the bundle
conditions, $\stageUNN$ a property of classes and $\stageUDR$ a definite-side condition. Suppose
\begin{enumerate}[label=(\roman*)]
\item $\sfname{MinkDiv}$ implies $\mathrm{M}_0(\stageUKP)$, $\mathrm{M}_1(\stageUKP,Q)$, $\mathrm{M}_2(\stageUKP,Q)$ and
$\mathrm{M}_3(\stageUKP,Q,Q')$;
\item $\sfname{ExBL}(\stageUKP,Q,Q',\stageULk)$;
\item $\sfname{AllRowsAlt}(\stageUKP,Q,Q',\stageULk,\stageUNN,\stageUDR)$.
\end{enumerate}
Then $\sfname{HypsAlt}(\stageUNN,\stageUDR)$ is non-empty.
\formallabel{GlobalTheory.Fc.nonempty_hyps_of_allRows_alt}
\end{theorem}

\begin{proof}
\emph{Step 1: one linked row.} The counting statement gives a slot $R_{30}$ with $\stageULB_{\stageUKP}(Q,R_{30})$, constants
$c_{30},c_{Ic}$ as in (F4) and $\sfname{Row}(Q'\stageUand R_{30},\stageULk,H_{30})$; thus
$\stageULin_{\stageUKP}(Q,Q';\stageULk,H_{30})$. By \cref{pf:stageU:hTu}, which uses (VR) and (GM),
$\sfname{Row}(Q',\stageULk,H_{Tu})$ holds for the property $H_{Tu}$ of $\mathsf{hTu}$ with $c_T(\eps,y)=c_T^{\min}(y)$. The slot
of all cores satisfies the large-base condition $\stageULB_{\stageUKP}(Q,\cdot)$
(\lref{GlobalTheory.Fc.largeBaseProperty_true}, which states this), so \cref{pf:stageU:rowcalc}(a) gives
$\stageULin_{\stageUKP}(Q,Q';\stageULk,H_{Tu})$. Combining the twenty-five linked rows of (iii), the counting row and the size
row by \cref{pf:stageU:rowcalc}(b), we obtain $\stageULin_{\stageUKP}(Q,Q';\stageULk,H)$, where $H$ is the conjunction of
these twenty-seven properties.

\emph{Step 2: class selection.} \Cref{pf:stageU:classsel}, applied with (SM), (VR), the bundle conditions, the linkage
$\stageULk$, the property $H$, the function $c_T(\eps,y)=c_T^{\min}(y)$ and the hypotheses (i), (ii) and step 1, gives a class
$\stageUA$ and a raw selection $\sigma$ such that $\sfname{MinkDiv}$ implies $\stageUEx^{\mathrm{des}}(\stageUA,c_T)$, every
family in $\stageUA$ fulfils the seven duties and the lattice stability, and $H(\stageUA,\sigma)$. (The fact that $\sigma$ is a
$\mathrm{True}$-selection is not used.)

\emph{Step 3: the fields.} We take $\stageUA$, $\sigma$ and $c_T(\eps,y)=c_T^{\min}(y)$, and obtain the fields as follows.
\begin{itemize}
\item (F1) $\mathsf{hEx}$ follows from the output of step 2 by \cref{pf:stageU:frozen}(c), applied for every $\lambda>1$ and
every requirement $\mathcal G$ with $\sfname{HostExW}(\lambda,\mathcal G)$; $\mathsf{hduty}$ follows from the seven duties
by \cref{pf:stageU:frozen}(b); $\mathsf{hMt}$ is the lattice stability.
\item (F2) The corresponding rows in $H(\stageUA,\sigma)$ give the eight properties at all selected bundles for all $x>x_1$;
the restriction to $x$ beyond the threshold of the selection is not needed. For $\mathsf{selJ}$ the row gives the
existence of a join, and we choose one.
\item (F3) The row of $\mathsf{hn4}$ gives closedness on $W_1(\stageUA)$, and $W'_1(\stageUA)=W_1(\stageUA)$ by
\cref{pf:stageU:frozen}(a); similarly the nine rows on $W_3(\stageUA)$ give closedness on $D3(3,\stageUA)=W_3(\stageUA)$.
(Closedness on a class implies closedness on every subclass, \lref{Formalism.MainNotions.Pk.Closed.sub}, so the
inclusions suffice.) The fields $\mathsf{hTF}$, $\mathsf{sO}$, $\mathsf{zL}$, $\mathsf{sLp}$, $\mathsf{hServ}$, $\mathsf{rRows}$,
$\mathsf{nRest}$ are the corresponding rows, $\mathsf{h30}$ and $\mathsf{irrC}$ are the two parts of $H_{30}$, and
$\mathsf{hTu}$ is $H_{Tu}(\stageUA)$.
\item (F4) The constants $c_{30},c_{Ic}$ of step 1, with their bounds.
\item (F5) The hypotheses of the theorem.
\end{itemize}
\end{proof}

\paragraph{Parameters.}
For a proof $h_L$ of $\sfname{RecTExists}$, $c\in\R$ and an at-property $A$, let $\stageUpd(h_L,c)$ and $\stageUpA(h_L,c)$ be
the parameters of \cref{sec:families:nonempty}: the norm-choice extension of $\stageUpd(h_L,c)$ asks for a law and a grid
datum at level $(1,c,b)$ with full $\lambda$-tie, and that of $\stageUpA(h_L,c)$ adds a saturated state of the strict-henselian
acyclic level (the \emph{origin} data) and the at-property $A$; both have no spectral extension and the geometric predicate
of all geometries. Thus every datum of the extension of $\stageUpA(h_L,c)$ yields one of the origin extension, hence one of
$\stageUpd(h_L,c)$, hence a law. We use the following monotonicity
(\lref{GlobalTheory.Glob.open_and_linked_of_stronger_extension_alt}): \emph{if every datum of a norm-choice extension
$N'$ yields a datum of $N$, then every linked row and the open parameter statements pass from the parameter of $N$ to the
parameter of $N'$} (with the same spectral extension and geometric predicate). The \emph{functional row} and the
\emph{covolume row} are the linked statements for $\sfname{FUN}$ and $\sfname{L3}$ of \cref{sec:families:nonempty}, and
the \emph{norm row} of a property $\stageUNN$ of classes is $\sfname{Row}(\cdot,\cdot,\stageUA\mapsto\stageUNN(D3(3,\stageUA)))$.

\begin{theorem}[The package from unramified exponents]\label{pf:stageU:unram}
Assume $P$ and $\sfname{RowRat}$. Let $h_L$ be a proof of $\sfname{RecTExists}$, $c\in\R$ and $A$ an at-property with
$\sfname{EvLgG}(A)$; write $\stageUpA=\stageUpA(h_L,c)$. Suppose
\begin{enumerate}[label=(\roman*)]
\item $\sfname{LinkedAt}\bigl(\stageUpd(h_L,c),\ \sfname{Row}(\cdot,\cdot,\stageUA\mapsto
\sfname{Closed}(D3(3,\stageUA),\sfname{BoundAtTwoThree}))\bigr)$;
\item $\stageUpA$ is linked for the functional row;
\item $\stageUpA$ is linked for the covolume row.
\end{enumerate}
Then $\sfname{HypsAlt}(\sfname{NL},\sfname{DefiniteSideCondition})$ is non-empty.
\formallabel{GlobalTheory.Fc.S0.StageU.hypotheses_nonempty_of_unramified_exponents_alt}
\end{theorem}

\begin{proof}
Write $Q_{\mathrm b},Q_{\mathrm g},\stageULk$ for the slots of built and good points and the bundle linkage of $\stageUpA$
(\cref{pf:stageU:linked}(d)).
\begin{enumerate}[label=\arabic*.]
\item By \cref{thm:member-bundles} (\lref{GlobalTheory.T.large_unramified_exponent_bundles_exist_alt}: $P$,
$\sfname{RowRat}$ and $\sfname{EvLgG}(A)$ imply (SM), member existence $\sfname{MemberE}(\stageUpA)$, and $\sfname{ExBL}$
for the base predicate $\stageUKP_0$, the slots $Q_{\mathrm b}$, $Q_{\mathrm g}$ and the linkage $\stageULk$), these three
statements hold.
\item From the fields of $P$ asserting the non-endoscopic hypotheses and the compatibility of the Galois representations
attached to polarized cuspidal representations, \lref{GlobalTheory.T.hyps_e_third_layer_ser_n_of_geometric_galois}
deduces four statements: an elliptic hypothesis, the third-layer inputs, the rationality and integrality of twist series,
and a local Langlands statement. We use the second and the third.
\item By \cref{prop:families:points} (\lref{GlobalTheory.Fc.S0.ontoX_and_lamLow_of_unitaryGroupHypotheses'}: $P$
implies $\sfname{SurjectiveOntoX}$ and $\sfname{LambdaLowerBound}_2$), these two statements hold. By
\cref{thm:families:open} (\lref{GlobalTheory.T.open0_of_unitary_group_hypotheses}: $P$ implies the open parameter
statements for $\stageUpd(h_L,c)$), the open statements hold for $\stageUpd(h_L,c)$.
\item By \cref{thm:linkN}(b) (\lref{UpperTheory.GalC.final_alt_alt}: under $\sfname{Hypotheses}_2$, the compatibility of
Galois representations, depth-zero support, local finiteness, local Tate duality and the local Euler characteristic formula,
which are fields of $P$, and the two statements of step 2, hypothesis (i) implies that $\stageUpd(h_L,c)$ is linked for the
norm row of $\sfname{NL}$), $\stageUpd(h_L,c)$ is linked for the norm row of $\sfname{NL}$.
\item By the monotonicity, applied to the extensions of $\stageUpd(h_L,c)$ and $\stageUpA$, the open statements and the norm
row of $\sfname{NL}$ hold for $\stageUpA$.
\item By \cref{prop:families:allrows} (\lref{GlobalTheory.Fc.AllAtMemberStage0.allRows_of_open0_alt}: global Artin
reciprocity, the Siegel--Weil formula, five statements on real unitary groups and their automorphic forms,
$\sfname{SurjectiveOntoX}$, $\sfname{LambdaLowerBound}_2$ and $h_L$ imply that a parameter satisfying the open statements and
linked for the covolume, functional and norm rows of $\stageUNN$ satisfies
$\sfname{AllRowsAlt}(\stageUKP_0,Q_{\mathrm b},Q_{\mathrm g},\stageULk,\stageUNN,\sfname{DefiniteSideCondition})$), applied with
the fields of $P$, steps 3 and 5 and the hypotheses (ii), (iii), the row package
$\sfname{AllRowsAlt}(\stageUKP_0,Q_{\mathrm b},Q_{\mathrm g},\stageULk,\sfname{NL},\sfname{DefiniteSideCondition})$ holds.
\item The uniform statements hold: $\sfname{LitThin}$ by \cref{pf:stageU:litthin}; (Wu), (TY) and (NG) as stated above, (NG)
from the second Siegel--Weil field of $P$; (RF) and (GM) by \cref{pf:stageU:findim}; (MT) and (VR) by
\cref{pf:stageU:mainrow,pf:stageU:volratio}, from the same field of $P$.
\item \Cref{pf:stageU:allrows}, applied to $\stageUKP_0$, $Q_{\mathrm b}$, $Q_{\mathrm g}$, $\stageULk$, $\sfname{NL}$ and
$\sfname{DefiniteSideCondition}$, gives the conclusion: its hypothesis (i) holds because $\sfname{MemberE}(\stageUpA)$ is the
conjunction of $\mathrm{M}_0(\stageUKP_0)$, $\mathrm{M}_1(\stageUKP_0,Q_{\mathrm b})$, $\mathrm{M}_2(\stageUKP_0,Q_{\mathrm b})$
and $\mathrm{M}_3(\stageUKP_0,Q_{\mathrm b},Q_{\mathrm g})$ (\cref{pf:stageU:memberE}) and holds by step 1, (ii) is step 1 and
(iii) is step 6.\qedhere
\end{enumerate}
\end{proof}

\begin{theorem}[\cref{thm:hyps-nonempty}]\label{pf:stageU:main}
Let $P$ and $F$ be the two bundles of classical results of \cref{sec:families:theta}, and let $h_A$, $h_M$, $h_X$, $h_L$ be
proofs of $\sfname{ThetaIntegralsAgree}$, $\sfname{LocalMatching}$, $\sfname{CovolumeBoundOfX}$ and $\sfname{RecTExists}$.
Assume that for every $c\in\R$
\[
\sfname{LinkedAt}\bigl(\stageUpd(h_L,c),\ \sfname{Row}(\cdot,\cdot,\stageUA\mapsto
\sfname{Closed}(D3(3,\stageUA),\sfname{BoundAtTwoThree}))\bigr).\tag{$h_N$}
\]
Then $\sfname{HypsAlt}(\sfname{NL},\sfname{DefiniteSideCondition})$ is non-empty.
\formallabel{UpperTheory.Lit.nonempty_hyps_definite_side_of_closed_bound}
\end{theorem}

\begin{proof}
Let $C$ be the field of $P$ containing the join statement and the functional equation of Hecke $L$-functions.
\begin{enumerate}[label=\arabic*.]
\item \emph{Separation.} Howe's unramified statement holds (\lref{GlobalTheory.Hk3.howeUnram_holds}, which proves it).
By \cref{prop:families:fs} (\lref{UpperTheory.Lit.functionalsSeparate_of_howeUnram_of_siegelWeil}: $F$, Howe's
unramified statement and the Siegel--Weil formula imply $\sfname{FunctionalsSeparate}$), applied with the first Siegel--Weil
field of $P$, $\sfname{FunctionalsSeparate}$ holds.
\item \emph{Rationality.} By \cref{prop:families:rowrat} (\lref{GlobalTheory.T.test_vector_assembly_holds}: the Siegel
set and real Howe duality fields of $F$, Green's theorem on cuspidal characters from $P$, $F$, the two statements of $C$, $h_A$,
$h_M$ and $\sfname{FunctionalsSeparate}$ imply the row test at all starting configurations; its proof rests on the
properties of theta-lift rows, \lref{GlobalTheory.RowClaim.Join.thetaLift_row_properties_conj}), the row test holds at
all starting configurations, and this is $\sfname{RowRat}$ (\lref{GlobalTheory.Mem.LiteralAutomorphy.rowRatW_of_rowTest},
which records it).
\item \emph{Fit.} By \cref{prop:families:points} (\lref{GlobalTheory.BigTop.Fun.Close.fourier_coefficient_fit_of_join_alt}:
$F$, the two statements of $C$, $h_A$, $h_M$, $\sfname{FunctionalsSeparate}$ and $\sfname{LambdaLowerBound}_2$ imply
$\sfname{ConstantFunctionFits}$), with $\sfname{LambdaLowerBound}_2$ from
\lref{GlobalTheory.Fc.S0.ontoX_and_lamLow_of_unitaryGroupHypotheses'}, $\sfname{ConstantFunctionFits}$ holds.
\item \emph{Grids.} The twisted Weil representations of the square configurations converge, their local factors match and
their theta integrals agree (\lref{GlobalTheory.WeilConv.twisted_weil_representation_convergence},
\lref{GlobalTheory.T.twisted_local_factors_match} and \lref{GlobalTheory.T.twisted_square_theta_integrals_eq}, which
deduce these three statements from $F$, $h_A$ and $h_M$). Hence, by
\lref{GlobalTheory.Mem.LiteralAutomorphy.evaluation_grid_of_join} (real Howe duality, $h_M$, $F$, the two statements of
$C$, $h_A$, these three statements and $\sfname{FunctionalsSeparate}$ imply that there is $C_{F1}$ such that
$\sfname{EvLgG}(\sfname{GridA}(C_F))$ for all $C_F\ge C_{F1}$), there is such a $C_{F1}$: full grids at level $C_F$ exist
eventually and for large twisting data.
\item \emph{Functional rows.} By \lref{GlobalTheory.Fc.StageOne.rowFun_at_gridA_alt} (five fields of $P$ and
$\sfname{ConstantFunctionFits}$ imply: there is $C_{F0}$ such that for all $C_F\ge C_{F0}$, every $h_L$ and every norm-choice
extension $N$ whose data yield grid data at level $C_F$ as in $\stageUpd(h_L,C_F)$, the open statements for the parameter of
$N$ extended by the at-property $\sfname{GridA}(C_F)$ imply that this parameter is linked for the functional row), there is
such a $C_{F0}$.
\item Put $c^*=\max(C_{F0},C_{F1})$ and $A=\sfname{GridA}(c^*)$; then $\sfname{EvLgG}(A)$ by step 4.
\item \emph{Open statements.} By \cref{thm:families:open} (\lref{GlobalTheory.T.open0_of_unitary_group_hypotheses}) the open
statements hold for $\stageUpd(h_L,c^*)$, and by the monotonicity they hold for $\stageUpA(h_L,c^*)$.
\item \emph{Functional and covolume rows.} The origin data at $c^*$ yield the grid data of $\stageUpd(h_L,c^*)$, and
$c^*\ge C_{F0}$; so step 5, applied to the origin extension, shows that $\stageUpA(h_L,c^*)$ is linked for the functional row.
By \lref{GlobalTheory.Glob.CovolRow.linked_bundles_l3_of_open_alt} (five fields of $P$, $h_X$ and $h_L$ imply: for every
norm-choice extension whose data yield a law, the open statements imply that the parameter is linked for the covolume row),
applied to the extension of $\stageUpA(h_L,c^*)$, this parameter is linked for the covolume row.
\item \Cref{pf:stageU:unram}, applied with $c=c^*$, the at-property $A$, $\sfname{RowRat}$ from step 2, $\sfname{EvLgG}(A)$,
the hypothesis $(h_N)$ at $c^*$ and step 8, gives the conclusion.\qedhere
\end{enumerate}
\end{proof}

In the application in \cref{sec:families}, $h_L$ is \lref{UpperTheory.Zero.LLC.recTExists}, which constructs the local
parameter maps, and $(h_N)$ holds by \cref{lem:families:at23} (see \cref{rem:families:supply}).

\begin{remark}[On the formal statements]\label{pf:stageU:rem}
\begin{enumerate}[label=(\arabic*)]
\item The rationality statement produced in step 2 of the proof of \cref{pf:stageU:main} and $\sfname{RowRat}$ have the same
definition.
\item In \cref{pf:stageU:allrows} member existence is required only under the hypothesis $\sfname{MinkDiv}$; in the proof of
\cref{pf:stageU:unram} it is supplied unconditionally, so that this hypothesis is not used there.
\item Of the four statements of step 2 of the proof of \cref{pf:stageU:unram}, the elliptic hypothesis and the local
Langlands statement are not used.
\item The eight fields of (F2) carry the hypothesis that $x$ is at least the threshold of the selection; it is not used.
\item The hypothesis $(h_N)$ of \cref{pf:stageU:main} is used only at $c=c^*$.
\item $P$ contains two fields stating the Siegel--Weil formula in the convergent range, in two formulations; the first is used
for $\sfname{FunctionalsSeparate}$, the second for the row package, (NG), (MT) and (VR).
\item The function $c_T$ of the package does not depend on $\eps$, and the property ``$\sigma$ is a $\mathrm{True}$-selection''
produced by class selection is not used.
\item Each of the three components of $\sfname{LitThin}$ is stated for every datum $D$, which does not occur in it.
\end{enumerate}
\end{remark}

% ===========================================================================
\subsection{Results used in other sections}\label{sec:pf:stageU:other}

The results of this subsection are ingredients of other sections of this part; they are grouped by topic, and the
notation is that of the section using them, recalled briefly. Ingredients which are not proved here are cited with
their content.

% ---------------------------------------------------------------------------
\subsubsection{Downward theta lifts}\label{sec:pf:stageU:down}

We use the real data of \cref{pf:liftinj:realdefs}: a global datum $D$ with ranks $r=\dim W_i$, $\dim V_i=r+t$, a real datum $R$
over $(D,\mathcal A,K_s)$, a face $F$, the standing hypotheses (G) and the fullness condition (F); for $0\le i<m$ and a tuple
$\mu$ the space $H_V(i,\mu)$ with extension map $\mathrm{ext}_V$, the Weil datum $\omega=\omega_P(i,\mu)$ with characters
$(\chi_W,\chi_V)=(\chi_W(i,\mu),\chi_V(i,\mu))$, the theta maps $\theta^F=\theta^F_{i,\mu}$ and $\theta^G=\theta^G_{i,\mu}$, the
embedding $e_W$, the downward integrals $\mathrm{down}(y,f,\ell,g)$ of \cref{eq:pf:liftinj:downint}, good places, Hecke
families, systems, Satake families and segments. Let $L_V$ be the level group of $V_i$ at the level primes $P_g$, with level
character $\chi_{\mathrm{lev}}$, and for $g$ in the finite group $G$ let $\kappa(g)\in L_V$ be a level lift of $e_G(g)$.

\begin{definition}[Level characters, finite type, the downward properties]\label{pf:stageU:downdefs}
Fix $i$ and $\mu$.
\begin{enumerate}[label=(\alph*)]
\item For a level prime $P$, $0\ne a\in k_P^r$, $\mathrm{tr}\in\{0,1\}$ and $\epsilon=\pm1$ let $m_{P,a,\mathrm{tr},\epsilon}:G\to\C$ be
$g\mapsto\chi_{i+1,P}(\det e_G(g)_P)^{\epsilon}$ if $e_G(g)$ is trivial at the level primes other than $P$ and $e_G(g)_P$ either moves
every vector by a multiple of $a$ ($\mathrm{tr}=1$) or fixes the hyperplane $a^\perp$ pointwise ($\mathrm{tr}=0$), and $0$ otherwise.
A vector $y\in H_V(i,\mu)$ has \emph{no level-character component} if, for all such data and every linear form $\ell$, the
group-algebra operator of $m_{P,a,\mathrm{tr},\epsilon}$ (through the level lifts) kills $h\mapsto\ell(\mathrm{ext}_V(y)(h))$.
\item A theta function $f$ for $\omega$ is of \emph{finite type} if every finitely supported combination of finite-adelic
Heisenberg translates which kills the distinguished vector of $\omega$ kills $f$. It is of \emph{level type} if
$\omega(\kappa,1)f=\chi_{\mathrm{lev}}(\kappa)^{-1}f$ for all $\kappa\in L_V$.
\item For a theta function $f$ and data as in (a), the \emph{mirabolic average} of $f$ is
$\sum_{g\in G}m_{P,a,\mathrm{tr},\epsilon}(g)\,\chi_{\mathrm{lev}}(\kappa(g))^{-1}\,\omega(\kappa(g)^{-1},1)f$; $f$ is \emph{clean} if all its
mirabolic averages vanish.
\item The properties of $(i,\mu)$ used in \cref{sec:pf:liftinj:realdata} are: $(\mathrm{Gen}_f)$ every $y$ without level-character
component such that $\mathrm{down}(y,f,\ell,g)=0$ for all $f$ of finite type and all $\ell,g$ has no lift; $(\mathrm{Gen}_\infty)$ if
$\theta^F(y)=0$ then $\mathrm{down}(y,f,\ell,g)=0$ for all $f$ of finite type and all $\ell,g$; $(\mathrm{GM})$ $e_W(\theta^G(y))=\theta^F(y)$
for good $y$; $(\mathrm{GC})$ good vectors have no level-character component; $(\mathrm{EC})$ for every Hecke family and every system
$s$ there is a system $c$ with $\theta^F(y)\in\mathrm{part}_W(c)$ for all eigenvectors $y$ with system $s$; $(\mathrm{PE})$ and
$(\mathrm{UC}(\mathrm{ok}))$ as in \cref{sec:pf:liftinj:realdata}, \eqref{eq:pf:liftinj:unr}; $(\mathrm{NV})$ and $(\mathrm{TNZ})$ as in
\cref{pf:liftinj:realdefs}.
\end{enumerate}
\formal{GlobalTheory.DownLift.NoChi}\formal{GlobalTheory.DownLift.FinLike}\formal{GlobalTheory.DownLift.GenFinite.TypeVec}
\formal{GlobalTheory.DownLift.GenFinite.Clean}\formal{GlobalTheory.DownLift.GenFin}\formal{GlobalTheory.DownLift.GenInf}
\formal{GlobalTheory.DownLift.GoodMem}\formal{GlobalTheory.DownLift.GoodNoChi}\formal{GlobalTheory.DownLift.DownEig}
\formal{GlobalTheory.DownLift.ParamsEx}\formal{GlobalTheory.DownLift.UnrCorr}
\end{definition}

In the following results $R$ is a real datum over $D$ and $F$ a face such that (G) and (F) hold, and $i,\mu$ are arbitrary.

\begin{lemma}[Generation at finite places]\label{pf:stageU:genfin}
Assume Howe's unramified statement. Then $(\mathrm{Gen}_f)$ holds.
\formallabel{GlobalTheory.DownLift.GenFinite.genFin}
\end{lemma}

\begin{proof}
Let $y$ have no level-character component and suppose that $\mathrm{down}(y,f,\ell,g)=0$ for all $f$ of finite type. Let
$\mathrm{Pr}(\varphi)$ mean: $\varphi$ is a theta function and $\mathrm{down}(y,\varphi,\ell,g)=0$ for all $\ell,g$. The downward integrand
is integrable for every theta function (\lref{GlobalTheory.T.integrableOn_downwardThetaLift_integrand}), so $\mathrm{down}$ is
linear in $\varphi$; theta functions form a vector space stable under $\omega$; hence $\mathrm{Pr}$ holds on the span of any set on
which it holds. We show $\mathrm{Pr}$ on successively larger sets.
\begin{enumerate}[label=\arabic*.]
\item \emph{Translates.} Since $\theta(h,g;\omega(1,g')f)=\theta(h,gg';f)$, we have
$\mathrm{down}(y,\omega(1,g')f,\ell,g)=\mathrm{down}(y,f,\ell,gg')$, which vanishes for $f$ of finite type.
\item \emph{Mirabolic averages.} The downward integrals of a mirabolic average of a theta function vanish for vectors without
level-character component (\lref{GlobalTheory.T.level_averages_theta_pair_zero_without_character}).
\item \emph{Level type.} A function of level type lies in the span of the clean functions of level type and of the mirabolic
averages (\lref{GlobalTheory.T.thetaFunction_levelType_mem_span_averages}), and a clean function of level type lies in the span
of the translates $\omega(1,g')f$ of functions $f$ of finite type
(\lref{GlobalTheory.T.thetaFunction_levelType_mem_span_translates}, which uses Howe's unramified statement). By steps 1 and 2,
$\mathrm{Pr}$ holds for functions of level type.
\item \emph{Differences.} For $\kappa\in L_V$ and a theta function $f$,
$\mathrm{down}(y,\omega(\kappa,1)f,\ell,g)=\chi_{\mathrm{lev}}(\kappa)^{-1}\mathrm{down}(y,f,\ell,g)$
(\lref{GlobalTheory.T.downwardLift_level_acts_by_inverse_character}), so $\mathrm{Pr}$ holds for
$\omega(\kappa,1)f-\chi_{\mathrm{lev}}(\kappa)^{-1}f$.
\item Every theta function lies in the span of the functions of level type and of these differences
(\lref{GlobalTheory.T.thetaFunction_mem_span_levelType_differences}).
\end{enumerate}
Hence $\mathrm{Pr}(f)$ for every theta function $f$, i.e.\ $y$ has no lift.
\end{proof}

\begin{lemma}[Generation at real places]\label{pf:stageU:geninf}
$(\mathrm{Gen}_\infty)$ holds.
\formallabel{GlobalTheory.GenReal.genInf}
\end{lemma}

\begin{proof}
Let $\theta^F(y)=0$, $f$ of finite type, and $\ell,g$ given. Choose a measure of covolume one with a fundamental domain on the
space of $\omega$ (\lref{GlobalTheory.Size.exists_isCovolOne}). Since $\theta^F(y)=0$, every downward integral of $y$ against the
Fock vectors attached to a low-type vector and a coefficient vanishes
(\lref{GlobalTheory.T.theta_lift_vanishes_single_coefficient_of_kernel_kills}), and therefore so does every downward integral of
$y$ against the Fock vectors $\mathrm{pt}(i,\mu,P)$ of all polynomials $P$
(\lref{GlobalTheory.T.realFockThetaLift_vanishing_of_generating_vectors}: these Fock vectors are generated, for the downward
integrals, by the former). Suppose $\mathrm{down}(y,f,\ell,g)\ne0$ and put $\varepsilon=|\mathrm{down}(y,f,\ell,g)|$. The downward integral
is continuous for a Sobolev norm of order $d$: there are $d$ and $\delta>0$ such that $|\mathrm{down}(y,f,\ell,g)-\mathrm{down}(y,\mathrm{pt}
(i,\mu,P),\ell,g)|<\varepsilon$ whenever the Sobolev distance of $f$ and $\mathrm{pt}(i,\mu,P)$ is $<\delta$
(\lref{GlobalTheory.T.downward_theta_integral_sobolev_continuous}); and a function of finite type is a Sobolev limit of the Fock
vectors $\mathrm{pt}(i,\mu,P)$ (\lref{GlobalTheory.T.thetaFunction_fixedFinitePart_sobolevLimit_fockVectors}). For $P$ with Sobolev
distance $<\delta$ we get $\varepsilon=|\mathrm{down}(y,f,\ell,g)-0|<\varepsilon$, a contradiction.
\end{proof}

\begin{lemma}[Good vectors and the idempotent]\label{pf:stageU:goodmem}
$(\mathrm{GM})$ holds: $e_W(\theta^G(y))=\theta^F(y)$ for every good $y\in H_V(i,\mu)$.
\formallabel{GlobalTheory.GoodLift.goodMem_of_given}
\end{lemma}

\begin{proof}
For $g\in G$ choose a level lift $\kappa(g)$ of $e_G(g)$ (\lref{GlobalTheory.GoodLift.liftsV_of_core}) and a translate
$y_g\in H_V(i,\mu)$ with $\mathrm{ext}_V(y_g)(h)=\chi_{\mathrm{lev}}(\kappa(g))^{-1}\mathrm{ext}_V(y)(h\kappa(g))$
(\lref{GlobalTheory.GoodLift.translate_of_core}). Let $e_C$ be the idempotent of the face. Since $y$ is good,
$\mathrm{ext}_V(\sum_ge_C(g)y_g)=\mathrm{ext}_V(y)$ (expand the group-algebra operator of $e_C$,
\lref{GlobalTheory.GoodLift.gaL_apply_choose}, and use \lref{GlobalTheory.GoodLift.sum_translates}). The difference
$z=e_W(\theta^G(y))-\theta^F(y)$ lies in the image of $e_W$ (the core conditions), and elements of this image vanish if their values
on all low-type vectors $v$ vanish (\lref{GlobalTheory.LiftRows.spanP_of_given}). For a low-type vector $v$, the core conditions
give $e_W(\theta^G(y))(v)=\sum_ge_C(g)\,\mathrm{act}_W(e_G(g))\,\theta^F(y)(v)$ (\lref{GlobalTheory.GoodLift.eWf_apply_sum}); the
translation rule $\theta^F(y_g)(v)=\mathrm{act}_W(e_G(g))\,\theta^F(y)(v)$ (\lref{GlobalTheory.GoodLift.thF_translate_low}) and
linearity turn this into $\theta^F(\sum_ge_C(g)y_g)(v)$, which equals $\theta^F(y)(v)$ because $\theta^F(\cdot)(v)$ only depends on
$\mathrm{ext}_V$ (\lref{GlobalTheory.GoodLift.thF_low_eq}). Hence $z(v)=0$ for all $v$ and $z=0$.
\end{proof}

\begin{lemma}[Good vectors have no level-character component]\label{pf:stageU:goodnochi}
$(\mathrm{GC})$ holds.
\formallabel{GlobalTheory.GoodTypeNoChi.goodNoChi}
\end{lemma}

\begin{proof}
Let $y$ be good, $\ell$ a linear form and $u=\ell\circ\mathrm{ext}_V(y)$. The function $u$ has the principal level type
(\lref{GlobalTheory.T.definiteSpace_forms_functional_principalType}), and level lifts exist
(\lref{GlobalTheory.GoodTypeNoChi.p_lifts}). For functions of this type the group-algebra operators compose like
convolution: the operator of $m=m_{P,a,\mathrm{tr},\epsilon}$ applied to the operator of $e_C$ is the operator of
$x\mapsto\sum_ym(y)e_C(y^{-1}x)$ (\lref{GlobalTheory.GoodTypeNoChi.p_comp}). Since $y$ is good, the operator of $e_C$ fixes $u$,
so the operator of $m$ applied to $u$ is the operator of $m\ast e_C$ applied to $u$. Finally $m\ast e_C=0$
(\lref{GlobalTheory.T.goodProjector_conv_mirabolicAvg_eq_zero}, deduced from the vanishing of mirabolic averages on the series of
characters of the family, \lref{GlobalTheory.GoodTypeNoChi.seriesKillsAvg}). Hence the operator of $m$ kills $u$.
\end{proof}

\begin{lemma}[Eigen-compatibility]\label{pf:stageU:downeig}
Assume Howe's unramified statement. Then $(\mathrm{EC})$ holds.
\formallabel{GlobalTheory.DownEigen.downEig}
\end{lemma}

\begin{proof}
Let $h$ be a Hecke family and $s$ a system. At a good place $P$ the characters $\chi_W,\chi_V$ are unramified
(\lref{GlobalTheory.T.splittingCharacters_unramified_offLevels}), and for an eigenvector $y$ with system $s$ the Hecke operator
of $V_i$ with coefficients $c_V$ acts on $y$ by the scalar $\mathrm{ev}_P(s(P),c_V)$. If the vacuum vector of $\omega$ is fixed at $P$
and the Hecke elements $c_V$ of $V_i$ and $c_W$ of $W_i$ \emph{match} under $\omega$ at $P$ (they act in the same way on the
vacuum, $c_V$ through its dual), then the theta map intertwines them: the right action of the $W_i$-Hecke operator of $c_W$
on $\theta^F(y)$ is $\mathrm{ev}_P(s(P),c_V)\,\theta^F(y)$ (\lref{GlobalTheory.DownEigen.thF_match'}). It remains to see that every
Hecke operator of $W_i$ at $P$ acts on $\{\theta^F(y)\}$ by a scalar independent of $y$.
\begin{enumerate}[label=(\roman*)]
\item If $P$ is self-conjugate (inert), the vacuum is fixed (\lref{GlobalTheory.VacFixInert.vacFix_of_inert}, using that $r$ and
$t$ are odd and the characters are unramified), and by Howe's unramified statement every Hecke element of $W_i$ at $P$ matches
some Hecke element of $V_i$ (\lref{GlobalTheory.Hk3.CorrInert.sphCorr_of_inert}).
\item If $P$ is not self-conjugate, the explicit matching for $\GL$ at split places
(\lref{GlobalTheory.Hk3.SplitMatch.match_split}) expresses the action of $V_i$-Hecke elements on the vacuum through a
triangular system in the $W_i$-Hecke elements of labels $0,\dots,r$, with non-zero diagonal coefficients (powers of $\Nm P$ and
of the character values, \lref{GlobalTheory.DownEigen.diag_ne_zero}); solving it
(\lref{GlobalTheory.DownEigen.tri_solve'}) shows that these $W_i$-Hecke operators act by scalars, and the operators of the
other labels at $P$ are then scalars as well (\lref{GlobalTheory.DownEigen.scalar_of_split'}).
\end{enumerate}
The level operators fix the vectors $\theta^F(y)$ at good places (\lref{GlobalTheory.LevOp.level_operators}). Hence the
vectors $\theta^F(y)$, $y$ an eigenvector with system $s$, are simultaneous eigenvectors of the Hecke operators of $W_i$ at all good
places with eigenvalues independent of $y$; this system $c$ satisfies $\theta^F(y)\in\mathrm{part}_W(c)$.
\end{proof}

\begin{lemma}[Existence of parameters]\label{pf:stageU:paramsex}
Assume rapid decay of cusp forms. Then $(\mathrm{PE})$ holds.
\formallabel{GlobalTheory.DownLift.Params.paramsEx}
\end{lemma}

\begin{proof}
The first half is \lref{GlobalTheory.DownLift.Params.paramsV} (the system of a non-zero eigenvector in $H_V(i,\mu)$ has a
Satake family of rank $\dim V_i$ off the bad places, from the unramified Hecke algebras of the definite group), and the second
half is \lref{GlobalTheory.DownLift.Params.paramsW} (a system $c$ with $\mathrm{part}_W(c)\ne0$ has a Satake family of rank $r$ whose
elements satisfy $|b|\le\Nm w^{(r-1)/2}$ at good $w$, from rapid decay of cusp forms on $W_i$).
\end{proof}

\begin{lemma}[The unramified correspondence at inert places]\label{pf:stageU:unrinert}
Assume Howe's unramified statement. Then $(\mathrm{UC}(\mathrm{ok}))$ holds for $\mathrm{ok}(w)$: ``$w$ is self-conjugate''.
\formallabel{GlobalTheory.DownInert.unrCorr_inert}
\end{lemma}

\begin{proof}
Let $y$ be an eigenvector with system $s$, $0\ne\theta^F(y)\in\mathrm{part}_W(c)$, $B$ and $B'$ Satake families of ranks $r+t$ and $r$
for $s$ and $c$, and $w$ a good self-conjugate place; $w$ is inert (\lref{GlobalTheory.DownInert.isInertAt_of_conj}) and $\chi_W,\chi_V$
are unramified at $w$ (\lref{GlobalTheory.UnrOffLevels.unrPl_of_good}, \lref{GlobalTheory.DownInert.unrAt_of_unrPl}).
\begin{enumerate}[label=\arabic*.]
\item There is a low-type vector $v_0$ with $\theta^F(y)(v_0)\ne0$, since $\theta^F(y)\ne0$ is determined by its values on low-type
vectors (\lref{GlobalTheory.LowTypeDetermined.eq_zero_of_lowT}). The function $z=\theta^F(y)(v_0)$ on $G_{W_i}$ is spherical at
$w$ with Satake parameter $B'(w)$ (\lref{GlobalTheory.LFun.isParamFunV_iff_of_mem_partFull}).
\item Expanding the coefficient space in a basis $(b_k)$, $z(g)=\sum_k\int_{\mathfrak F_V}\theta(h,g;f_k)\,x_k(h)\,d\nu_V(h)$, where
$f_k$ are Fock vectors which are vacuum vectors at $w$ (\lref{GlobalTheory.T.fockVectors_chain_isVacuum_off_thetaLevels},
\lref{GlobalTheory.DownInert.vacFin_of_isVacAt}) and $x_k=b_k^*\circ\mathrm{ext}_V(y)$ are continuous automorphic functions on
$G_{V_i}$, each zero or spherical at $w$ with Satake parameter $B(w)$ (\lref{GlobalTheory.DownInert.paramV_coord}).
\item The pairing $(f,x)\mapsto\int_{\mathfrak F_V}\theta(h,\cdot;f)x(h)\,d\nu_V(h)$ is bilinear and equivariant for
$\Uni(V_i)(F_w)\times\Uni(W_i)(F_w)$ (right translation invariance of integrals over fundamental domains,
\lref{GlobalTheory.T.setIntegral_fundamentalDomain_right_translate_invariant}; integrability,
\lref{GlobalTheory.DownInert.integrable_contAut}).
\item The abstract form of Howe's unramified correspondence at inert places
(\lref{GlobalTheory.DownInert.param_down}: for such an equivariant pairing of spherical vacuum vectors and spherical vectors with
parameter $B(w)$, a non-zero image vector with parameter $B'(w)$ forces $(-1)^{\dim W}B(w)=(-1)^{\dim V}B'(w)+\mathrm{str}(q_w,\dim V-\dim W)$)
gives the relation between $B(w)$ and $B'(w)$, with $\dim V_i-\dim W_i=t$.
\end{enumerate}
Writing the unramified character values at the inert place through the segment of length $t$
(\lref{GlobalTheory.ChainOnto.val_seg_of_isInertAt}, \lref{GlobalTheory.DownInert.seg_coe}) turns this relation into
\eqref{eq:pf:liftinj:unr}.
\end{proof}

\begin{lemma}[The unramified correspondence at the other places]\label{pf:stageU:unrsplit}
$(\mathrm{UC}(\mathrm{ok}))$ holds for $\mathrm{ok}(w)$: ``$w$ is not self-conjugate''.
\formallabel{GlobalTheory.SplitCorr.unrCorr_split}
\end{lemma}

\begin{proof}
This is \lref{GlobalTheory.SplitCorr.unrCorr_at} (the relation \eqref{eq:pf:liftinj:unr} at good places $w$ with $\bar w\ne w$,
from the correspondence for $\GL_n$ at split places), since $w$ is not self-conjugate iff $\bar w\ne w$
(\lref{GlobalTheory.GoodPrimes.isConjPl_iff}).
\end{proof}

\begin{lemma}[Type non-vanishing from generation]\label{pf:stageU:typenz}
If $(\mathrm{NV})$, $(\mathrm{Gen}_\infty)$, $(\mathrm{Gen}_f)$, $(\mathrm{GC})$ and $(\mathrm{GM})$ hold for $(i,\mu)$, then $(\mathrm{TNZ})$ holds for
$(i,\mu)$.
\formallabel{GlobalTheory.DownLift.typeNZ_of_nonVan_gen}
\end{lemma}

\begin{proof}
Let $y$ be a good eigenvector whose system has a segment, with $\theta^G(y)=0$, and suppose $y\ne0$. By $(\mathrm{GM})$,
$\theta^F(y)=e_W(\theta^G(y))=0$; by $(\mathrm{Gen}_\infty)$ all downward integrals of $y$ against functions of finite type vanish; by
$(\mathrm{GC})$, $y$ has no level-character component; so $y$ has no lift by $(\mathrm{Gen}_f)$. This contradicts $(\mathrm{NV})$.
\end{proof}

\begin{lemma}[Segments from non-zero lifts]\label{pf:stageU:hasstring}
Let $\mathrm{ok}$ be a predicate on places, and assume $(\mathrm{UC}(\mathrm{ok}))$, $(\mathrm{UC}(\neg\mathrm{ok}))$, $(\mathrm{PE})$ and $(\mathrm{EC})$ for
$(i,\mu)$. If $y$ is an eigenvector with system $s$ for a Hecke family and $\theta^F(y)\ne0$, then $s$ has a segment.
\formallabel{GlobalTheory.LiftInjProof.hasString_of_lift}
\end{lemma}

\begin{proof}
By $(\mathrm{EC})$ there is $c$ with $\theta^F(y)\in\mathrm{part}_W(c)$; by $(\mathrm{PE})$ there are a Satake family $B$ of rank $\dim V_i$ for
$s$ and, since $\mathrm{part}_W(c)\ni\theta^F(y)\ne0$, a Satake family $B'$ of rank $r$ for $c$ with $|b|\le\Nm w^{(r-1)/2}$. At a good $w$,
$(\mathrm{UC}(\mathrm{ok}))$ or $(\mathrm{UC}(\neg\mathrm{ok}))$ gives \eqref{eq:pf:liftinj:unr}, so $B(w)=B''(w)+\mathrm{str}(\Nm w,\chi_V(w),t)$
with $B''(w)=\chi_V(w)\chi_W(w)^{-1}B'(w)$, whose elements have absolute value $\le\Nm w^{(r-1)/2}$ because
$|\chi_V(w)|=|\chi_W(w)|=1$ (\lref{GlobalTheory.LFun.norm_val}).
\end{proof}

\begin{lemma}[Injectivity from type non-vanishing]\label{pf:stageU:liftinj}
If for every $i$ the properties $(\mathrm{TNZ})$, $(\mathrm{EC})$, $(\mathrm{UC}(\mathrm{ok}))$ for ``$w$ is not self-conjugate'' and for ``$w$ is
self-conjugate'', and $(\mathrm{PE})$ hold at $(i,\mu_0)$, then $(\mathrm{INJ})(R)$ holds.
\formallabel{GlobalTheory.LiftInjProof.liftInj_of_typeNZ}
\end{lemma}

\begin{proof}
Let $x,y$ be eigenvectors with system $s$ with $\theta^G(x)\ne0$, $y$ good and $\theta^G(y)=0$. The core conditions give
$e_W\circ\theta^G=\theta^F$ on the low types with $e_W$ injective, so $\theta^F(x)\ne0$. By \cref{pf:stageU:hasstring} (with
$\mathrm{ok}$ = ``self-conjugate''), $s$ has a segment, and $(\mathrm{TNZ})$ applied to $y$ gives $y=0$.
\end{proof}

% ---------------------------------------------------------------------------
\subsubsection{The doubled space and its Schr\"odinger model}\label{sec:pf:stageU:dbl}

Let $A,B$ be hermitian spaces over $E$ for a datum $D$, $n=\dim A$, $A^\square=A\oplus(-A)$ the doubled space with its Siegel
parabolic $P^\square$ (the stabiliser of the diagonal), $G_n$ the standard unitary group of the skew-hermitian form
$\left(\begin{smallmatrix}0&1\\-1&0\end{smallmatrix}\right)$ with its Siegel parabolic $P_n$, Levi elements $m(a)$, unipotent elements
$n(b)$, Weyl element $w_n$ and standard maximal compact subgroup $K$. For a majorant $a$ of $A$ let $\hat K(a)$ be the maximal
compact subgroup of $G_{A^\square}$ attached to $a$ and the self-dual frame lattice. An element $\mu\in E$ is \emph{skew} if
$\mu\ne0$ and $\bar\mu=-\mu$; a \emph{chart} for $\mu$ is an isomorphism $\kappa:G_{A^\square}\to G_n$ with $T_\mu\kappa(g)=gT_\mu$
for a fixed rational matrix $T_\mu$ built from $\mu$. Let $\omega^\square$ be a Weil datum for $(A^\square,B)$ with characters
$(1,\chi_B)$, and $\nu$ a measure of covolume one with a fundamental domain on $\mathrm{Mat}_{\dim B\times n}(\A_E)$. The
\emph{Schr\"odinger transform} $\mathrm{sch}(f)$ of a theta function $f$ for $\omega^\square$ is the function on
$\mathrm{Mat}_{\dim B\times n}(\A_E)$ obtained by the partial Fourier transform (normalised by $\nu$) from the mixed model of the doubled
space. A representation $\omega_S$ of $G_n$ on the Schwartz space $\mathcal S(B^n)$ of \cref{pf:liftinj:schrdef} is a
\emph{transport} of $\omega^\square$ through $\kappa$ if $\omega_S(\kappa(g))\mathrm{sch}(f)=\mathrm{sch}(\omega^\square(g,1)f)$ for all $g$
and all theta functions $f$.

\begin{lemma}[Theta kernels as lattice sums]\label{pf:stageU:thetaS}
For every theta function $f$ for $\omega^\square$ and all $g\in G_{A^\square}$, $h\in G_B$,
$\theta(g,h;f)=\sum_{\xi\in\mathrm{Mat}_{\dim B\times n}(E)}\mathrm{sch}(\omega^\square(g,1)f)(h^{-1}\xi)$.
\formallabel{GlobalTheory.T.thetaKernel_eq_tsum_schrodingerFunction}
\formallabel{GlobalTheory.StdUnn.thetaKO_eq_tsum_sch}
\end{lemma}

\begin{proof}
We have $\omega^\square(g,h)f=\omega^\square(1,h)\,\omega^\square(g,1)f$, and $\omega^\square(g,1)f$ is a theta function. The value at $0$
of a theta function is the lattice sum of its Schr\"odinger transform (the mixed-model formula,
\lref{GlobalTheory.HeisMixSource.heis_mixS}), and the transform intertwines the action of $h\in G_B$ with the substitution
$x\mapsto h^{-1}x$ (\lref{GlobalTheory.T.schrodingerFunction_unitary_levi_radical_actions}).
\end{proof}

\begin{lemma}[Iwasawa decomposition of the doubled group]\label{pf:stageU:iwasawa}
For every majorant $a$ of $A$, $G_{A^\square}$ has an Iwasawa decomposition $G_{A^\square}=P^\square\hat K(a)$ (adelically, with the
local factors at all places).
\formallabel{GlobalTheory.Core.Rk.IP.hasIwasawa_Khat}
\end{lemma}

\begin{proof}
The frame lattice of $A$ is self-dual (\lref{GlobalTheory.WeilAtPlace.isSelfDualC_frameLat}). For a self-dual lattice and a
majorant the adelic decomposition is assembled from local decompositions at the split finite places, at the other finite
places, and at the real places (\lref{GlobalTheory.Core.Rk.IP.hasIwasawa_KAd_selfDual}, which combines
\lref{GlobalTheory.T.hasIwasawaDecomposition_doubled_split_place},
\lref{GlobalTheory.T.iwasawa_doubled_hermitian_space_nonsplit_place} and
\lref{GlobalTheory.T.iwasawa_doubled_hermitian_space_real_place}).
\end{proof}

\begin{lemma}[Standard charts]\label{pf:stageU:chart}
For every skew $\mu$ there is a chart $\kappa:G_{A^\square}\to G_n$ for $\mu$.
\formallabel{GlobalTheory.T.exists_standardUnitaryChart_doubledSpace}
\formallabel{GlobalTheory.StdUnn.chart_exists}
\end{lemma}

\begin{proof}
The matrix $T_\mu$ is invertible and carries the Gram matrix of $A^\square$ to the standard skew-hermitian form, up to the unit
$\mu$ (\lref{GlobalTheory.StdUnn.Tm_gram_inv}). Hence $g\mapsto T_\mu^{-1}gT_\mu$ maps the unitary group of $A^\square$ isomorphically
onto $G_n$ (adelically), and $T_\mu\kappa(g)=gT_\mu$ by construction.
\end{proof}

\begin{lemma}[Conjugating the compact subgroup]\label{pf:stageU:conjcompact}
There is a skew $\mu\in E$ such that for every $A$, every majorant $a$ of $A$ and every chart $\kappa$ for $\mu$ there is
$p\in P^\square$ with $k\in\hat K(a)\iff\kappa(p^{-1}kp)\in K$.
\formallabel{GlobalTheory.T.exists_siegel_parabolic_conj_doubled_compact}
\formallabel{GlobalTheory.StdUnn.ConjElt.exists_isConj}
\end{lemma}

\begin{proof}
Choose $\mu$ with auxiliary scalars (\lref{GlobalTheory.StdUnn.scalData_exists}) and half data
(\lref{GlobalTheory.StdUnn.exists_halfData}), independently of $A$. Given $A$ and $a$, choose a lattice frame $b$ adapted to $a$
(\lref{GlobalTheory.StdUnn.exists_isLatticeFrame_bc_eq}); its Gram matrix $H$ is unimodular and hermitian
(\lref{GlobalTheory.StdUnn.Unimod.gram_unimodular}) with $H^2=1$ at the real places, because $a$ is a majorant. These data give
frames $M,M'$ of $A^\square$ (\lref{GlobalTheory.StdUnn.Frame.exists_frame}) and an element $p$ of the Siegel parabolic transporting
the frame (\lref{GlobalTheory.StdUnn.ConjElt.exists_siegel_frame}); conjugation by $p$ then carries $\hat K(a)$ onto
$\kappa^{-1}(K)$ (\lref{GlobalTheory.StdUnn.Conj.isConj_of_frame}).
\end{proof}

\begin{lemma}[The shifted additive character]\label{pf:stageU:shiftchar}
Let $\psi$ be a standard additive character of $\A_F/F$ (of parameter $-1$), $\theta$ a family of archimedean embeddings and $\mu$
skew. Then $\lambda=-\xi_0/\mu$ lies in $F$, the character $\psi_\lambda(x)=\psi(\lambda x)$ is a non-trivial character of $\A_F/F$, and
its real parameters $l(\tau)$ satisfy $|l(\tau)|=|\tau(\xi_0)|/|\tau(\mu)|$.
\formallabel{GlobalTheory.T.shiftedAdditiveCharacter_archimedeanParameters}
\formallabel{GlobalTheory.StdUnn.shift_char_realPar}
\end{lemma}

\begin{proof}
Since $\bar\xi_0=-\xi_0$ and $\bar\mu=-\mu$, $\lambda=-\xi_0/\mu$ is fixed by complex conjugation, hence lies in $F$, and $\lambda\ne0$. The
character $\psi_\lambda$ is continuous, trivial on $F$ (as $\lambda F=F$) and non-trivial (if $\psi(x)\ne1$ then $\psi_\lambda(\lambda^{-1}x)\ne1$).
At a real place $\tau$ the parameter of $\psi_\lambda$ is $\tau(\lambda)$ times that of $\psi$, which has absolute value $1$; hence
$|l(\tau)|=|\tau(\lambda)|=|\tau(\xi_0)|/|\tau(\mu)|$.
\end{proof}

\begin{lemma}[Theta functions and Schwartz functions]\label{pf:stageU:schschwartz}
For every standard $\psi$, family $\theta$, spaces $A,B$ and $\nu$ as above: the Schr\"odinger transform of a theta function for
$(A^\square,B)$ lies in $\mathcal S(B^n)$; every function in $\mathcal S(B^n)$ is the transform of a theta function; and the transform
is injective on theta functions.
\formallabel{GlobalTheory.T.schrodingerTheta_eq_schwartzBruhat}
\formallabel{GlobalTheory.SchSchwartzAll.schSchwartz}
\end{lemma}

\begin{proof}
The transform of a theta function is a Schwartz--Bruhat function
(\lref{GlobalTheory.T.isSchwartzBruhat_doubledPairSchrodingerFunction}), and the Schwartz--Bruhat functions on
$\mathrm{Mat}_{\dim B\times n}(\A_E)$ are the elements of $\mathcal S(B^n)$ (\lref{GlobalTheory.T.isSchwartzBruhat_iff_mem_adelicMatrixSchwartzSpace}).
Conversely every Schwartz--Bruhat function is the transform of a theta function
(\lref{GlobalTheory.T.exists_thetaFunction_schrodingerFunction_eq}). If $\mathrm{sch}(f)=\mathrm{sch}(f')$, then $f-f'$ is a theta function
with transform $0$ (the transform is linear), hence $f-f'=0$ (\lref{GlobalTheory.T.doubledPairThetaFunction_eq_zero_of_schrodinger_eq_zero}).
\end{proof}

\begin{lemma}[Transport to the standard group]\label{pf:stageU:transport}
Let $\omega^\square,\nu$ be as above, $\theta$ an archimedean family and $\kappa:G_{A^\square}\to G_n$ an isomorphism. Then there is a
transport $\omega_S$ of $\omega^\square$ through $\kappa$ on $\mathcal S(B^n)$.
\formallabel{GlobalTheory.T.exists_weilDatum_transport_schwartzBruhat_standard}
\formallabel{GlobalTheory.StdUnn.exists_transport_std}
\end{lemma}

\begin{proof}
By \cref{pf:stageU:schschwartz} the transform is a linear bijection from theta functions onto $\mathcal S(B^n)$
(\lref{GlobalTheory.T.schrodingerFunction_linear_thetaFunctions} for linearity), and theta functions are stable under
$\omega^\square(g,1)$. Hence $\omega_S(\kappa(g))=\mathrm{sch}\circ\omega^\square(g,1)\circ\mathrm{sch}^{-1}$ defines a representation of $G_n$ with
the required property.
\end{proof}

\begin{lemma}[The three formulas]\label{pf:stageU:schrformulas}
Let $\mu$ be skew, $\kappa$ a chart for $\mu$, $\omega_S$ a transport of $\omega^\square$ through $\kappa$ and $\lambda=-\xi_0/\mu$. Then
$\omega_S$ satisfies \eqref{eq:pf:liftinj:schr} for $(\psi_\lambda,\chi_B^{-1})$ and $\nu$:
$\omega_S(m(a))\Phi(x)=\chi_B^{-1}(\det a)|\det a|^{\dim B/2}\Phi(xa)$, $\omega_S(n(b))\Phi(x)=\psi_\lambda(\frac12\tr(b\langle x,x\rangle))\Phi(x)$ and
$\omega_S(w_n)\Phi(x)=\int\Phi(y)\psi_\lambda(\tr\langle x,y\rangle)\,d\nu(y)$.
\formallabel{GlobalTheory.T.transportedSchrodinger_standardLeviAction}
\formallabel{GlobalTheory.T.transportedSchrodinger_siegelUnipotentAction}
\formallabel{GlobalTheory.T.weyl_formula_standardUnitaryGroup_weilDatum}
\formallabel{GlobalTheory.StdUnn.transport_levi}\formallabel{GlobalTheory.StdUnn.Unip.transport_unip}\formallabel{GlobalTheory.StdUnn.WeylStd.main}
\end{lemma}

\begin{proof}
Write $\Phi=\mathrm{sch}(f)$ (\cref{pf:stageU:schschwartz}); then $\omega_S(g)\Phi=\mathrm{sch}(\omega^\square(\kappa^{-1}(g),1)f)$. Under the chart,
$m(a)$ corresponds to the Levi element of $a$ for the Witt decomposition of $A^\square$
(\lref{GlobalTheory.T.standardUnitaryChart_blockDiagonal_eq_levi}), $n(b)$ to a unipotent element of the Siegel parabolic with
parameter $c$ (\lref{GlobalTheory.T.standardUnitaryChart_unipotent_eq_wittUnipotent}), and $w_n$ to a rational Weyl element
(\lref{GlobalTheory.T.standardUnitaryChart_weylElement_rational}). The Schr\"odinger transform turns the action of the Levi and
unipotent elements of the doubled space into the first two formulas
(\lref{GlobalTheory.T.schrodingerFunction_unitary_levi_radical_actions}; the character of the unipotent element is rewritten
through $\psi_\lambda$ by \lref{GlobalTheory.StdUnn.Unip.chr_eq_shift}), and the action of the Weyl element into the Fourier
transform for $\psi$, which is the third formula after the substitution $\lambda=-\xi_0/\mu$
(\lref{GlobalTheory.T.weylElement_acts_on_schrodinger_by_fourier}). 
\end{proof}

\begin{lemma}[Inverses of adapted Hecke characters]\label{pf:stageU:heckeinv}
If $\chi$ is a Hecke character of $\A_E^\times$ (continuous, trivial on $E^\times$) whose restriction to $\A_F^\times$ is $\eta^m$, then
$\chi^{-1}$ has the same two properties.
\formallabel{GlobalTheory.T.isHeckeChar_restrictionEta_inv}
\end{lemma}

\begin{proof}
$\chi^{-1}$ is continuous and $\chi^{-1}(x)=\chi(x)^{-1}=1$ for $x\in E^\times$; and $\chi^{-1}|_{\A_F^\times}=\eta^{-m}=\eta^m$ since $\eta^2=1$
(\lref{GlobalTheory.StdUnn.isHeckeChar_resEta_inv}).
\end{proof}

% ---------------------------------------------------------------------------
\subsubsection{Hecke point laws and the third layer}\label{sec:pf:stageU:points}

The next two results concern the geometric theta functions of the twisted operators: for a CM field $E$ with a host, a
datum, a level $l$, a moduli space $X_o$ of the corresponding abelian schemes with a coefficient datum $c$, and a prime $P$ of $E$
split and off the level, a Hecke kernel $C$ of type $(P,p,i)$ and a Hecke isogeny $\Psi$ with kernel $C$ between the abelian
schemes over two points $x,x'$, a \emph{theta trace} of a section $s$ along $\Psi$ is a section $s'$ related to $s$ by the trace
of $\Psi$, and a \emph{theta generator} is a section which does not vanish identically in the relevant sense. The \emph{theta
transport law} says that for all such data the theta transport is total and single-valued; the \emph{Hecke point laws} are
finiteness of the Hecke correspondences, totality and single-valuedness of the Hecke relation on frames, and detection of
points by their Hecke images.

\begin{lemma}[The theta transport law]\label{pf:stageU:transportlaw}
The theta transport law holds.
\formallabel{GlobalTheory.ThetaTransport.thetaTransportLaw}
\formallabel{GlobalTheory.ThetaTransport.existsUnique_trace_isThGen}
\end{lemma}

\begin{proof}
Let $x$ be a point with the given coefficient datum, $C$ a Hecke kernel of type $(P,p,i)$, $\Psi$ a Hecke isogeny with kernel $C$, and
$s$ a theta generator at $x$. A section $s$ has a unique theta trace along $\Psi$ (\lref{GlobalTheory.ThetaTransport.existsUnique_trace},
from the splitting of $P$ off the level), and $s'$ is a theta generator iff it does not vanish, which holds for the trace of a
generator, because the trace of $s$ vanishes only if $s$ does (\lref{GlobalTheory.ThetaTransport.isThGen_iff}). Hence there is
a unique theta trace, and it is a theta generator; this is the
totality and the uniqueness of the transport.
\end{proof}

\begin{lemma}[Hecke point laws]\label{pf:stageU:pointlaws}
The theta transport law implies the Hecke point laws; in particular the Hecke point laws hold.
\formallabel{GlobalTheory.HeckePointLaw.heckePointLaws_of_transport}
\formallabel{GlobalTheory.HeckePointLaw.heckePointLaws}
\end{lemma}

\begin{proof}
Finiteness of the Hecke correspondences and detection of points are \lref{GlobalTheory.HeckePointLaw.heckeFinite} and
\lref{GlobalTheory.HeckePointLaw.pointsDetect}. For a Hecke prime, the totality part of the transport law gives a frame
related to a given frame by the Hecke relation (\lref{GlobalTheory.HeckePointLaw.exists_hkRelFr}), and the uniqueness part shows
that two frames related to the same frame coincide (\lref{GlobalTheory.HeckePointLaw.hkRelFr_unique}); this gives totality and
single-valuedness. The second statement follows with \cref{pf:stageU:transportlaw} (in the form
\lref{GlobalTheory.T.heckeIsogeny_thetaGenerator_transport}).
\end{proof}

For the following results we use the notions of the third layer (\cref{sec:families:nonempty}): the published third-layer inputs
$\Pi$, the multiplicity statement $\mathrm{(SMt)}$ (module inputs exist at all states), the parameter statement for $h_L$
(stability of the conjugate-dual local parameters at primes not dividing the level), the eventual statements ``at every member of
a layer point'', character choice data $P$ for $(h_L,D)$ with its character frame, norm choice, tuple, module family and
representation, real expansions $X$ agreeing with $P$, and joins of $P$ (a module form, canonical row modules at the local indices
and theta modules).

\begin{lemma}[The multiplicity statement]\label{pf:stageU:smt}
$\mathrm{(SMt)}$ holds; in particular $\Pi$ implies $\mathrm{(SMt)}$.
\formallabel{GlobalTheory.MtInMain.prN_sMt_alt_alt}
\formallabel{GlobalTheory.MtId.sMt}
\end{lemma}

\begin{proof}
At every state the pinned data provide a module input (\lref{GlobalTheory.MtId.inMt_state}). The hypothesis $\Pi$ is not used.
\end{proof}

\begin{lemma}[Stability of local parameters]\label{pf:stageU:conjdual}
$\Pi$ and $\mathrm{(SMt)}$ imply the parameter statement for every $h_L$.
\formallabel{GlobalTheory.HeightCounting.Up.conjDualParam_stable_of_prime_not_dividing_level_alt}
\end{lemma}

\begin{proof}
The parameter statement is a step statement (eventually, for large twisting data, at all states). Combine (by the calculus of
steps, \cref{pf:stageU:steps}) the following eventual statements: the lifting of all parameters from parallel inputs with module
inputs (\lref{GlobalTheory.T.liftAll_of_parIn_of_mtInput}), the local components of parallel inputs
(\lref{GlobalTheory.T.local_component_of_par_in_third_layer}), the sizes, congruences, host facts and column choices
(\lref{GlobalTheory.T.eventually_sizes_congruences_hostFacts_columnChoices}), and $\mathrm{(SMt)}$. At a state $x$ with a pin, the
parity conditions of the row hold (\lref{GlobalTheory.HeightCounting.Up.ParityConditions.ofSt}); the module input gives the lift of
all parameters and the local component; the host facts give the CM and definiteness conditions; the parameters satisfy their
laws (\lref{UpperTheory.GalN.par_laws_bare}) and the base change identities
(\lref{UpperTheory.BaseChangeIdentities.Conclusion1_alt_alt_alt}, from base change, the Bernstein decomposition and the generic
and standard modules in $\Pi$). These data give the cuspidal parameter condition, which is the parameter statement at $x$
(\lref{GlobalTheory.T.cuspParam_of_definite_liftAll}).
\end{proof}

\begin{lemma}[Cuspidal parameters at members]\label{pf:stageU:cuspparam}
Assume $\Pi$. For every $h_L$, eventually at every member of a layer point the row sides with the local representation of the
member have good cuspidal parameters.
\formallabel{GlobalTheory.HeightCounting.Up.cuspParameter_at_member_alt}
\end{lemma}

\begin{proof}
As in \cref{pf:stageU:conjdual}, combine the lifting statement, the local component statement and the size and host facts.
At a member, the norm-choice data give a parallel copy of the row with a parallel input
(\lref{GlobalTheory.HeightCounting.Up.ParCopy.parIn_of_nTrue_alt}) and a module input
(\lref{GlobalTheory.HeightCounting.Up.ParCopy.exists_mtInput_of_nTrue}); the lift and the local component of the copy, with the host
facts, give good cuspidal parameters for every row side with the same local representation
(\lref{GlobalTheory.HeightCounting.Up.ParCopy.good_cusp_param_of_definite_heq_locR}).
\end{proof}

\begin{lemma}[Joins of character choices]\label{pf:stageU:join}
Assume Harish-Chandra finiteness and rapid decay of cusp forms. Let $P$ be character choice data for $(h_L,D)$ and $X$ a real
expansion agreeing with $P$. Then $P$ has a join.
\formallabel{GlobalTheory.T.nonempty_join_of_agree}
\formallabel{GlobalTheory.JoinExists.main_alt_alt}
\end{lemma}

\begin{proof}
For every local index $i$ there is a canonical row module $Y_i$ at $j_i=j(i)$ for $X$
(\lref{GlobalTheory.T.exists_canonical_module_at_successor}, applied at $j_i-1<m$). The module form is obtained from the
agreement (\lref{GlobalTheory.T.mtForm_of_realAt_agree}), and for each canonical module a theta module exists
(\lref{GlobalTheory.T.nonempty_theta_mod_of_is_canon}). These data form a join.
\end{proof}

\begin{lemma}[Canonical row modules]\label{pf:stageU:rowmodcanon}
In the situation of \cref{pf:stageU:join}, for every $j<m$ and $j'=j+1$ there is a row module at $j'$ which is canonical for $X$.
\formallabel{GlobalTheory.RowModCanon.main_alt_alt}
\end{lemma}

\begin{proof}
Let $M$ be the space of $K$-finite automorphic vectors of the principal level frame of $W_j$ (the space $M_{\inf}$ of the
real datum of $X$), with the right action of $G_{W_j}$, and put $c=\prod_\rho\dim K_\rho/n_P>0$, the product over the archimedean
places of the dimensions of the relevant $K$-types divided by a positive index. Take the row module with underlying space $M$
and form $c\langle x,y\rangle$, where $\langle\,,\rangle$ is the Petersson inner product over the fundamental domain of $W_j$. By
Harish-Chandra finiteness and rapid decay the inner products of elements of $M$ exist and satisfy Schur orthogonality
(\lref{GlobalTheory.MultFinite.hcFinite_and_schurMult}); hence the form is sesquilinear, hermitian and positive definite on $M$
(\lref{GlobalTheory.FinTypesRows.minf_sesq}), and $M$ is stable under the congruence subgroups
(\lref{GlobalTheory.FinTypesRows.minf_stable}, using \lref{GlobalTheory.T.unitary_congruence_conj_contains_congruence}). The
module of the datum at $j+1$ maps into the good part of the multiplicity space, which lies in $M$
(\lref{GlobalTheory.FinTypesRows.levOK_levS}); this is the canonicity. (The statement form
\lref{GlobalTheory.T.exists_canonical_module_at_successor} is used in \cref{pf:stageU:join}.)
\end{proof}

% ---------------------------------------------------------------------------
\subsubsection{Galois-theoretic conditions at members}\label{sec:pf:stageU:galois}

We use the conditions on a global datum $D$ with its norm-choice data $(\mathrm{nd},\mathrm{naut})$ named in \cref{pf:linkN:GP}
(the Galois preconditions): the ties $\mathrm{tie}_0$ (local parameter, automorphy and Galois ties), the Galois closure condition
$\mathrm{adm}_2$, the isobaric multiplicity-one and Chebotarev conditions $\mathrm{adm}_5$, the tame label condition, the linked
inputs and the point conditions.

\begin{lemma}[Admissibility of the residual tame eigenvalues]\label{pf:stageU:admtame}
For every datum $D$, the conditions $\mathrm{tie}_0$, $\mathrm{adm}_2$, $\mathrm{adm}_5$ and the tame label condition imply the admissibility of the
residual tame eigenvalues.
\formallabel{GlobalTheory.AdmTame.main_alt_alt_alt}
\end{lemma}

\begin{proof}
The tame label condition gives that the wild inertia at the type primes acts trivially; $\mathrm{adm}_2$ gives that the Galois
representations are unramified away from the relevant set (\lref{GlobalTheory.AdmTame.galAwayOfS'}); $\mathrm{adm}_5$ gives the
density of Frobenius elements (Chebotarev) and multiplicity one for the isobaric parts: if two one-part isobaric sums have the same
local components almost everywhere, their parts coincide. The local-parameter tie gives the laws of the characters at the type
primes. Multiplicity one and the character laws give invariance under twisting
(\lref{GlobalTheory.TypeLeaf.twistInv_of_charLaws_multOne}). Since $t$ is odd, the archimedean infinitesimal characters shifted by
$-t/2$ are integers (a direct check on the two possible parameters); with the twist invariance and the automorphy tie this gives the
point parameters (\lref{GlobalTheory.TypeLeaf.pointPar_of_twist}). The first character is trivial on the wild inertia at the type
primes (\lref{GlobalTheory.T.firstCharacter_trivial_wildInertia_typePrime}). These conditions and the full labels (from the tame
label condition, \lref{UpperTheory.GalS.fullLabS_of_fullLabS_trivial}) give the tame sums, i.e.\ the admissibility
(\lref{GlobalTheory.TypeLeaf.tameSum_of_inputs_alt}).
\end{proof}

\begin{lemma}[Residual irreducibility in rank $r$]\label{pf:stageU:resirr}
Let $D$ be a datum satisfying the three-prime admissibility condition, $\mathrm{adm}_5$, the residual representation condition and
the depth-zero condition, and assume the Brauer--Nesbitt irreducibility statement for the relevant Galois group. Then the residual
representations of rank $r$ are irreducible.
\formallabel{Formalism.MainNotions.AdmissibleThreePrime.residual_irreducible_of_theorem_three_prime_alt}
\end{lemma}

\begin{proof}
By the three-prime theorem (\lref{Formalism.MainNotions.AdmissibleThreePrime.theorem_three_prime_holds_alt_alt}), for every local
index $i$ with residue characteristic $\ell_i$ and every reduction $\bar\rho$ the residual representation is irreducible, also on
every normal subgroup of finite solvable index containing the type elements. Its proof: the bad sets of the two type primes
$g_0,g_1$ are disjoint, so $\ell_i$ avoids one of them, say that of $g_k$; for this $k$ the tame inertia at $g_k$ acts through a good
tuple of characters, and the irreducibility criterion for good tuples
(\lref{Formalism.CyclotomicNorms.irreducibility_of_good_tuple}) applies. Residual irreducibility in rank $r$ then follows with
$\mathrm{adm}_5$, the Brauer--Nesbitt statement \cite{CurtisReiner1962} and the residual and depth-zero conditions
(\lref{Formalism.MainNotions.N.residual_irreducible_of_theorem}).
\end{proof}

\begin{lemma}[The first provision from the summands]\label{pf:stageU:parts}
For every datum $D$, the summands suffice: the admissibility at all local indices, the label tie, the row sizes, the Hecke split
condition, the positivity of the $\lambda$-values, the theta condition and the unramifiedness of the cuspidal Galois
representations imply the first provision.
\formallabel{UpperTheory.GalS.Parts.main_alt_alt}
\end{lemma}

\begin{proof}
The hypotheses give the parts statement (\lref{UpperTheory.GalS.Parts.p41'_alt}) and, through the parts, the conclusion on
reduced parameters (\lref{GlobalTheory.NRows.Parts.conclusion42_of_parts}); the latter is equivalent to the parameter
proposition (\lref{Formalism.MainNotions.SumRepresentations.proposition_parameters_iff_reduced}). Together these are the first
provision.
\end{proof}

\begin{lemma}[Conditions (A)]\label{pf:stageU:aclose}
The theta-local inputs imply, for every $h_L$, every datum $D$ with character choice data $P$ and the jump-count condition, the
conditions (A) of the rows: the label tie, the Hecke split condition, $\mathrm{tie}_0$, $\mathrm{adm}_2$, the two isobaric multiplicity-one
statements, $\mathrm{adm}_5$, $\mathrm{adm}_7$, the trivial local condition and the splitting of the automorphisms.
\formallabel{UpperTheory.GalS.a_close_main}
\end{lemma}

\begin{proof}
The label tie comes from the character choice (\lref{GlobalTheory.TieLabAt.main_alt_alt}). The pin of $P$ gives isobaric
rigidity (\lref{GlobalTheory.NPin.isobaricRigid_of_pin_alt}), hence multiplicity one for isobaric cuspidal sums
(\lref{GlobalTheory.RowClose.isobaricCuspidalMultiplicityOne_of_twisted}) and for isobaric discrete self-dual sums
(\lref{GlobalTheory.SpehRigid.isobaricDiscreteSelfDualMultOne_of_isobaricSumMultOne}). The Chebotarev density of split primes
(\cref{pf:stageU:chebsplit}) and these two statements give $\mathrm{adm}_5$
(\lref{Formalism.MainNotions.SumRepresentations.condition_iff_additive_maps}). The automorphism tie and the splitting come from
\lref{GlobalTheory.AutomorphismSplitting.tieAut_and_split_alt}; $\mathrm{tie}_0$ combines the local parameter tie
(\lref{GlobalTheory.RowTieLP.main}), the automorphism tie and the Galois tie, which comes from the jump count and the de Rham
and potentially semistable inputs (\lref{GlobalTheory.T.tieGal_of_jumpsCard_of_deRhamPotSemistable}). Finally $\mathrm{adm}_2$,
$\mathrm{adm}_7$, the Hecke split condition and the trivial local condition are \lref{UpperTheory.GalS.galClose_alt},
\lref{UpperTheory.GalN.allCuspidalOpenBound_of_cuspU_character_choice},
\lref{GlobalTheory.RowClose.heckeSplitOne_of_charTied} and \lref{UpperTheory.GalS.trivialLocalCondition_of_inputs}.
\end{proof}

\begin{lemma}[Reduced multiplicity]\label{pf:stageU:redmult}
For every $h_L$ and datum $D$, the linked inputs and the point conditions imply the reduced multiplicity statement.
\formallabel{GlobalTheory.T.reducedMultiplicity_of_inputs}
\formallabel{GlobalTheory.PtSt.reducedMultiplicity_of_rows'}
\end{lemma}

\begin{proof}
Fix a local index $i$. The linked inputs give internal isotropy and, with the point conditions and \cref{pf:stageU:parts}, the
conclusion on reduced parameters. With the parallel-up condition, the membership of the downward lifts and type non-vanishing
(components of the point conditions), the multiplicity spaces $M_j^\sigma$ have the predicted image
(\lref{GlobalTheory.PtSt.map_MjSig_eq}), and the good part of the image is as required
(\lref{Formalism.MainNotions.SumRepresentations.good_part_image_iff}).
\end{proof}

\begin{lemma}[Character choices at members]\label{pf:stageU:pin}
Assume $\Pi$. For every $h_L$, eventually at every member of a layer point there is character choice data $P$ for the datum of the
member at which the member's character choice is pinned.
\formallabel{GlobalTheory.T.exists_character_choice_pin_at_of_third_layer}
\formallabel{GlobalTheory.PinLayer.specialization_at_member_of_layer_alt}
\end{lemma}

\begin{proof}
At a member, the types of the datum are true (from the group condition of the row data), and the datum is elementary: complex
conjugation preserves the coefficient field (which is Galois), the Harish-Chandra parameters are integral, and the places above the
level primes split (since the level primes are $\equiv1\pmod{Q_Sqd}$, \lref{GlobalTheory.HeightCounting.Up.isSplitPl_of_mem_levelPrimes}).
With Green's cuspidal characters (a field of $\Pi$), the injectivity of the row representations, the Chebotarev density of split
primes (\cref{pf:stageU:chebsplit}) and the label and zeta data of the real expansion, the pin exists
(\lref{GlobalTheory.PinLayer.pinAt_of_side_alt}).
\end{proof}

\begin{lemma}[Base change parameters up to permutation]\label{pf:stageU:bcvperm}
If a list $\mathrm{Pa}$ of parts is a base change parameter for a point at a local index, then so is every permutation of
$\mathrm{Pa}$.
\formallabel{GlobalTheory.RowThetaParam.isBCV_of_perm}
\end{lemma}

\begin{proof}
Each of the six defining conditions is invariant under permutations of the list: membership, the absence of repetitions, the
total rank (a sum over the list), the local isobaric sums at the places (\lref{GlobalTheory.RowThetaParam.isoLoc_perm}) and the
archimedean sums.
\end{proof}

\begin{lemma}[Characters to finite groups ignore the infinite part]\label{pf:stageU:finpart}
Let $E$ be totally complex, $B$ a finite group and $\psi$ a homomorphism from the id\`eles of $E$ to $B$. Then $\psi(a)=\psi(a_f)$ for
every id\`ele $a$, where $a_f$ is $a$ with all archimedean components replaced by $1$.
\formallabel{GlobalTheory.RingCls.map_finPartI}
\end{lemma}

\begin{proof}
Let $N=|B|$. Since every archimedean completion is $\C$, each archimedean component of $a$ has an $N$-th root; let $u$ be the id\`ele
with these roots at the archimedean places and $1$ elsewhere. Then $a=u^Na_f$, and $\psi(u)^N=1$, so $\psi(a)=\psi(a_f)$.
\end{proof}

% ---------------------------------------------------------------------------
\subsubsection{Steps, member statements and positions}\label{sec:pf:stageU:stepcalc}

\begin{lemma}[Calculus of steps and of member statements]\label{pf:stageU:steps}
Let $\stageUKP$ be a base predicate.
\begin{enumerate}[label=(\alph*)]
\item The steps (\cref{sec:pf:stageU:vocab}) from $p$ to $q$ and from $p$ to $r$ give the step from $p$ to $q\wedge r$; the steps from
$p$ to $q$ and from $q$ to $r$ give the step from $p$ to $r$. The same holds for the steps between properties of stage data.
\item If $p$ and $q$ hold eventually at every member of a layer point, so does $p\wedge q$; if $p$ does and $p$ implies $q$ pointwise,
so does $q$.
\item A property which holds over all rows holds at all states.
\item If $P$ implies the slot $R$ at all cores of layers, then every core in the position of $P$ lies in $R$.
\end{enumerate}
\formallabel{GlobalTheory.Mem.Top.GroupRetwisting.Step.and}\formallabel{GlobalTheory.Mem.Top.GroupRetwisting.Step.trans}
\formallabel{GlobalTheory.Mem.Top.Bl.StepW.and}\formallabel{GlobalTheory.HeightCounting.Up.Model.AtMem.and_z}
\formallabel{GlobalTheory.HeightCounting.Up.Model.AtMem.at_member_imp}
\formallabel{GlobalTheory.HeightCounting.Up.RowModel.overRows_allSt}\formallabel{GlobalTheory.Mem.Top.Bl.positionOf_ind}
\end{lemma}

\begin{proof}
(a) A step is an $\sfname{EvLgG}$-statement; two such statements hold simultaneously (minimum of the exponents, supremum of the
transported constants, maxima of the thresholds, unions of the exceptional sets), and at a context where both hold the
implications compose, resp.\ combine. For stage data the step is defined through a step of the first kind followed by a pointwise
implication, so the same argument applies. (b) The member statements are eventual statements of the same kind; combine as in
(a), resp.\ apply the implication at each member. (c) At a state, the pinned core provides the data of a row over which the
property holds. (d) A core in the position of $P$ is, up to the identification of the stage tuple with $w.\mathrm{st}(Z)$, a core
$c'$ of a layer with $P(w,Z,c')$; the hypothesis gives $R$ at it, and the identification transports $R$.
\end{proof}

% ---------------------------------------------------------------------------
\subsubsection{Moving layers}\label{sec:pf:stageU:moving}

For a layer $Z$ over column data $w$ (\cref{pf:linkN:moving}), a family $N$ of division lattices $N_u$ and a height function
$h$ on the words, $\mathrm{setDiv}(Z,N,h)$ is the layer with these division lattices and heights, and $\mathrm{upL}(Z,S',N,h)$ is the layer
with in addition the $S$-stage $S'$. Two row data over the layers are \emph{rows-same} if their rows coincide.

\begin{lemma}[Rows pass to moved layers]\label{pf:stageU:rowsmove}
Let $t$ be row data over $Z$ with object $O$ and $O'$ the transported object over $\mathrm{setDiv}(Z,N,h)$. There are row data $t'$ over
$\mathrm{setDiv}(Z,N,h)$ with object $O'$, rows-same with $t$, whose geometric part and packet are the transports of those of $t$, such that
the object condition and the non-emptiness of the reading data pass from $t$ to $t'$.
\formallabel{GlobalTheory.T.rows_pass_to_moved_divisionLatticesHeights}
\end{lemma}

\begin{proof}
The components of $t$ which do not involve the division lattices or the heights are kept, and the others are transported along
the identification of the stages; each listed property is checked componentwise
(\lref{GlobalTheory.Mem.R3.Rw.DivMv.rows_move_geo}).
\end{proof}

\begin{lemma}[Kept norm choices]\label{pf:stageU:keptchoice}
Let $\mathrm{nch}$ be a norm choice and $\mathrm{xd}$ theta data at $Z$. There are a norm choice $\mathrm{nch}'$ and theta data $\mathrm{xd}'$ at
$\mathrm{upL}(Z,S',N,h)$ such that $\mathrm{nch}'$ keeps $\mathrm{nch}$; if $S'$ is the $S$-stage of $Z$, they are the transports of $\mathrm{nch}$
and $\mathrm{xd}$; the norm-choice data and the goodness of $(\mathrm{nch},\mathrm{xd})$ pass to $(\mathrm{nch}',\mathrm{xd}')$; and for rows-same row data
whose geometric parts correspond, the norm choice with law of the extension passes from the old to the new layer.
\formallabel{GlobalTheory.T.exists_kept_choice_upper_lattice_alt}
\end{lemma}

\begin{proof}
Define $\mathrm{nch}'$ and $\mathrm{xd}'$ by transporting the components of $\mathrm{nch}$ and $\mathrm{xd}$ along the stage identification, replacing
only the parts which depend on the $S$-stage, the division lattices and the heights; the listed properties are then verified
component by component (\lref{GlobalTheory.LayerMove.exists_kept_nChoice_thData_upL}).
\end{proof}

\begin{lemma}[Transport of list pins, spectral data and division geometry]\label{pf:stageU:movetransport}
\begin{enumerate}[label=(\alph*)]
\item For rows-same row data $t$ over $Z$ and $t'$ over $\mathrm{upL}(Z,S',N,h)$, a list pin of level $(c,b)$ for $t$ is one for $t'$.
\item The spectral stage data move to the set-division lattices and heights.
\item For a geometric property $\Gamma$ and rows-same row data $t$ over $Z$ and $t'$ over $\mathrm{setDiv}(Z,N,h)$ with objects $O$ and
$O.\mathrm{tr}(N,h)$, the geometric closure of $\Gamma$ passes from $(t,O)$ to $(t',O.\mathrm{tr}(N,h))$.
\item For slot data $\mathcal S$, if positions move along set division, then witnesses move along set division.
\end{enumerate}
\formallabel{GlobalTheory.Mem.GridMove.listPin_rowsSame}
\formallabel{GlobalTheory.T.spectralStageData_transfer_lattices_heights}
\formallabel{GlobalTheory.T.divisionGeometryScalars_transfer_setDivision}
\formallabel{GlobalTheory.T.witnesses_pass_along_set_division}
\end{lemma}

\begin{proof}
(a) Rows-same row data have the same grid, transfer and core components, and the algebra structures agree (both are determined by
the embedding of the coefficient field); the list pin then transfers along the lattice move
(\lref{GlobalTheory.Mem.GridMove.listPin_upL}). (b) is \lref{GlobalTheory.Mem.SRe.SxPass.sxMv}: the spectral stage data are
transported componentwise. (c) is \lref{GlobalTheory.Mem.OxMove.oAt_setDiv}: the division geometry scalars only depend on
components which are transported. (d) is \lref{GlobalTheory.Mem.Top.Bl.witMv_of_positionMv}: a witness is a position together
with data which move with it.
\end{proof}

\begin{definition}[Moved grid data]\label{pf:stageU:gdc}
For grid data $G$ at a level choice $L$ with constants $\mathrm{sc}$ and a real number $b$, the moved grid data at the moved level choice
$\mathrm{levCh}(b,L)$ have the moved places, characters and vectors of $G$.
\formallabel{GlobalTheory.Mem.GridMove.gdC}
\end{definition}

% ---------------------------------------------------------------------------
\subsubsection{Hecke operators, class sums and characters}\label{sec:pf:stageU:hecke}

\begin{lemma}[Chebotarev density of split Frobenius elements]\label{pf:stageU:chebsplit}
Assume the Chebotarev density theorem for degree-one places. Then for every datum $D$ with norm-choice data, Galois group $G$ of
the relevant extension $\Omega/E$ and every finite set $S$ of places of $E$, the Frobenius elements at the places $w\notin S$ which
split over $F$ are dense in $G$.
\formallabel{GlobalTheory.ChebSplitDense.chebSplit_nd}
\formallabel{GlobalTheory.ChebSplitDense.chebSplit}
\end{lemma}

\begin{proof}
Outside a finite set $A$, every place of $E$ of degree one is split over $F$ (\lref{GlobalTheory.FrobDense.exists_finite_split}).
By the Chebotarev theorem for degree-one places, the Frobenius elements at the degree-one places outside $A\cup S$ are dense; these
places are split and outside $S$.
\end{proof}

\begin{lemma}[Components of Hecke characters at split places]\label{pf:stageU:compat}
Let $L$ be a CM field, $\chi$ a Hecke character of $L$ and $s$ a split place (a place $s.w$ of $L$ above a place $v$ of $L^+$ with
$L_{s.w}=L^+_v$). If $\xi,\xi'$ are characters of $(L^+_v)^\times$ which are trivial on $1+\mathfrak p_v^N$ for some $N$ and agree with the
component of $\chi$ at $s.w$ on $L^\times$ (embedded through $s$), then $\xi=\xi'$.
\formallabel{GlobalTheory.T.heckeCharacter_splitPlaceComponent_unique}
\formallabel{GlobalTheory.CompAtUniq.main}
\end{lemma}

\begin{proof}
Let $N$ be a common level and $\pi$ a uniformiser of $L^+_v$. Every $u\in(L^+_v)^\times$ can be written $u=\iota_s(k)u'$ with $k\in L^\times$ and
$u'\in1+\mathfrak p_v^N$, since $L$ is dense in $L_{s.w}=L^+_v$ (write $u=\pi^au_0$ with a unit $u_0$, choose $k\in L$ close to $u$). Then
$\xi(u)=\chi_{s.w}(k)=\xi'(u)$.
\end{proof}

\begin{lemma}[Refining class sums]\label{pf:stageU:classrefine}
Let $G$ be a group, $Q\le G$ a subgroup, $J'\le J$ subgroups, $T$ and $T'$ sets of representatives of $Q\backslash G/J$ and $Q\backslash G/J'$,
$F\subseteq J$ a set of representatives of $J/J'$, and $\varphi:G\to\C$ left $Q$-invariant and right $J'$-invariant. Write
$Q_t(J)=Q\cap tJt^{-1}$. Then
\[
\sum_{t'\in T'}|Q_{t'}(J')|^{-1}\varphi(t')=\sum_{t\in T}|Q_t(J)|^{-1}\sum_{u\in F}\varphi(tu).
\]
\formallabel{GlobalTheory.ClassRefine.sum_classReps_refine}
\end{lemma}

\begin{proof}
Group the pairs $(t,u)\in T\times F$ according to the representative $t'\in T'$ of $Qtu J'$. For fixed $t'$, let $t_0\in T$ be the
representative of $Qt'J$; a pair $(t,u)$ lies over $t'$ iff $t=t_0$ and $u\in U=\{u\in F:t_0u\in Qt'J'\}$, by uniqueness of the
representatives. For $u\in U$, $\varphi(t_0u)=\varphi(t')$ by the invariances. Hence the right side equals
$\sum_{t'}|Q_{t_0}(J)|^{-1}|U|\varphi(t')$, and it remains to show $|Q_{t_0}(J)|=|U|\cdot|Q_{t'}(J')|$; $U\ne\emptyset$ because $t'\in Q t_0J$. This
is an orbit--stabiliser count for the action of $Q_{t_0}(J)$ on $t_0J/J'$, whose orbit through the class of $t'$ corresponds to $U$
and whose stabiliser is conjugate to $Q_{t'}(J')$ (\lref{GlobalTheory.ClassRefine.card_stabQ_split}).
\end{proof}

\begin{lemma}[Standard facts on Hecke double coset operators]\label{pf:stageU:heckestd}
For the Hecke double coset operators $T$ of a unitary group at a prime $\mathfrak p$ of norm $q$ prime to the level, acting on the
spaces of automorphic forms of a given level and type, and their normalisations $T_N=q^{-k(n-1)}T$: the operators preserve the
spaces; they are linear; any two of them commute; and every eigenvalue $\mu$ of $T_N$ on a non-zero eigenvector is $\ell$-integral
up to a power of $q$: there is a power $Q=q^{m+k(n-1)}$, prime to $\ell$, with $Q\mu$ integral.
\formallabel{GlobalTheory.HeckeStdFacts.heckeStdPlain}
\end{lemma}

\begin{proof}
Preservation, linearity and commutativity are the corresponding properties of the double coset operators
(\lref{GlobalTheory.HeckeStdFacts.heckeD_mem}, \lref{GlobalTheory.HeckeStdFacts.heckeD_lin},
\lref{GlobalTheory.HeckeStdFacts.heckeD_comm}). For integrality, choose double coset data for $T$
(\lref{GlobalTheory.HeckeStdFacts.exists_data}); then $T$ acts through a sum of translates with integral coefficients, and an
eigenvalue of $q^{-k(n-1)}T$ on a non-zero form with integral values becomes integral after multiplication by a power of $q$
(\lref{GlobalTheory.HeckeStdFacts.eigen_integral}); since $\ell\nmid q$, this power is prime to $\ell$.
\end{proof}

\begin{lemma}[Hecke operators commute with the group action]\label{pf:stageU:commrho}
Under the frame, family and model conditions of a real datum, let $j$ be the reading at a base embedding $\sigma$, injective, and
$\rho$ the action of $G$ transported through $j$. Then every Hecke operator $h$ of Hecke data at $(P,k)$ commutes with every
$\rho(g)$.
\formallabel{GlobalTheory.HeckeSup.comm_rho_cl}
\end{lemma}

\begin{proof}
Since $j$ is injective, it suffices to compare $j(h\rho(g)x)$ and $j(\rho(g)hx)$. Both equal the Hecke operator on the model
applied to the action of $g$ on $j(x)$, up to the same scalar, because the Hecke operator on the model commutes with the action of
$G$ (\lref{GlobalTheory.HeckeSup.heckeOn_actW_cl}).
\end{proof}

\begin{lemma}[Laws of forcing]\label{pf:stageU:evolaws}
Let a datum with character frame, host embedding, module family, norm choice and tuple be given. Forcing by automorphic
semisimple data (\cref{sec:pf:baselat2}) is closed under conjunction and monotone; if an automorphic side exists and $P_1$ is forced,
then for every family $\mathrm{big}_0$ of finite sets there are $\mathrm{big}\supseteq\mathrm{big}_0$ and a local datum built for $\mathrm{big}$ at which
$P_1$ holds; and being built for $\mathrm{big}$ is antitone in $\mathrm{big}$.
\formallabel{GlobalTheory.CoefDuty.LargeRing.evO_laws_alt}
\end{lemma}

\begin{proof}
Antitonicity: a specialisation for $\mathrm{big}'\supseteq\mathrm{big}$ is one for $\mathrm{big}$ (restrict the condition on the sets). If $P_1$ is
forced by $b$ and $P_2$ by $c$, both are forced by $b\cup c$ by antitonicity, hence so is $P_1\wedge P_2$; monotonicity is immediate. If
$P_1$ is forced by $b$ and an automorphic side exists, there is a local datum built for $\mathrm{big}=\mathrm{big}_0\cup b$, and $P_1$ holds at
it since it is also built for $b$.
\end{proof}

\begin{lemma}[Complements of Galois representations]\label{pf:stageU:complement}
Assume the Chebotarev statement for the twisting systems and that the cuspidal automorphic representations of a datum $D$ have
attached Galois representations. Then for every local index $i$, every representation $\pi$ in the theta locus at $i$ and every
continuous Galois representation $r_0$ of dimension $n_V=r+t$ attached to $\pi$, there is a representation $\rho$ of dimension $r$
such that $r_0$ is the sum of $\rho$ and the $t$ characters $\bar\chi_{i,c+1}$, $c<t$ (with the stated compatibilities).
\formallabel{GlobalTheory.BlockComplement.complement_at}
\end{lemma}

\begin{proof}
The tie of $\pi$ gives an isobaric decomposition $\pi=\boxplus_{p\in\mathrm{Pa}}p$ containing the part $(1,\mu\eta,t)$, and the infinity types
of the other parts are regular and integral up to the shift $(n_V-d_p)/2$ (\lref{GlobalTheory.BlockComplement.inf_part}). The other
parts have traces of attached Galois representations, and their sum $T$ is a continuous trace
(\lref{GlobalTheory.BlockComplement.hasTrace_part}). At almost all places the trace of $r_0(\mathrm{Frob}_w)$ equals $T(\mathrm{Frob}_w)$ plus
the values of the $t$ characters $\bar\chi_{i,c+1}$ (comparison of Frobenius eigenvalues with the local parameters of the parts and of
the segment); both sides are continuous, so by the Chebotarev statement they agree everywhere. A representation whose trace is
$T+\sum_c\bar\chi_{i,c+1}$ has a complement $\rho$ to these characters (\lref{GlobalTheory.BlockComplement.exists_complement}).
\end{proof}

\begin{lemma}[Real idempotents]\label{pf:stageU:starcoeff}
Suppose that for both type primes $g_k$ the series idempotents satisfy $\overline{e_{\mathrm{ser},k}(f)(x)}=e_{\mathrm{ser},k}(f(0)^{-1}f^{\mathrm{rev}})(x)$
for every monic $f\in\F_{g_k}[X]$ with $f(0)\ne0$ and every $x$. Then all coefficients of the idempotent $e_0$ of the types are real.
\formallabel{GlobalTheory.GoodIdemReal.star_coeff_of_reverse}
\end{lemma}

\begin{proof}
The coefficient of $e_0$ at $\gamma$ is a sum over good tuples $a$ (families of monic polynomials of degree $r$ with non-zero constant
term, indexed by the places of the type set) of the products $\prod e_{\mathrm{ser},k}(a_{k,v})(\gamma_{k,v})$. The map $f\mapsto
f(0)^{-1}f^{\mathrm{rev}}$ is an involution of the monic polynomials of degree $r$ with non-zero constant term
(\lref{GlobalTheory.SeriesSupportTest.dl_dl}), and the induced involution of tuples preserves goodness. By hypothesis the complex
conjugate of the term of $a$ is the term of the image of $a$; reindexing the sum by the involution shows that the sum equals its
complex conjugate.
\end{proof}

\begin{lemma}[Hecke statements for the localised datum]\label{pf:stageU:heckeopen}
In the setting of the specialisation of a datum $D$ along a localisation datum $L$ (side data, a choice side, a module input, the
conjugation stability of the coefficient field, integrality of the exponents, splitting of the places above $p_{de},p'_{de}$ and the
Chebotarev splitting statement), for every specialisation and every automorphic specialisation there are Hecke open statements for
the localised datum, the localised module family and the tuple, whose automorphic class data are those of the host embedding.
\formallabel{GlobalTheory.FitHecke.heckeOpen_withLoc_alt}
\end{lemma}

\begin{proof}
The specialisation provides a normalised local compatibility (cover) for the localised datum: its Hecke base, the
coefficient data, the maps $\Psi$, the points and the compatible embeddings $\iota_C$ (coherent by the automorphic specialisation).
The Hecke fields are then obtained from the cover: the hypotheses outside the bad places at each local index are transported from the
module input and are monotone in the set of places (\lref{GlobalTheory.FitHecke.hypOutside_transport'},
\lref{GlobalTheory.FitHecke.hypOutside_mono}); the Hecke tie, the equality of orders, the separation, conjugation and sameness
of the sets $Q$ are \lref{GlobalTheory.FitHecke.hT_of}, \lref{GlobalTheory.FitHecke.sameO_of}, \lref{GlobalTheory.FitHecke.sep_of},
\lref{GlobalTheory.FitHecke.Q_conj_of} and \lref{GlobalTheory.FitHecke.Q_same_of}; the automorphic data are those of the
automorphic specialisation.
\end{proof}

% ---------------------------------------------------------------------------
\subsubsection{Rows, columns and bundles}\label{sec:pf:stageU:rowsmisc}

\begin{lemma}[The theta hypotheses of a column inclusion]\label{pf:stageU:hypcolin}
Let $I$ be an ok column inclusion over a base $B$ with stages $(C,W,V,S,N,G)$. For every index $a$ and word $u$, the theta hypothesis
holds for the datum of $I$, the seam of the stages, the jet data and character of $I$, its selection, the tuple and constants at
$a$, and $u$.
\formallabel{GlobalTheory.KerT.hyp_of_colIn}
\end{lemma}

\begin{proof}
The components of the hypothesis are: admissibility of the definite base data
(\lref{GlobalTheory.T.definiteBaseData_admissible}); the seam conditions (\lref{GlobalTheory.Mem.Col2.seamOf_ok}), the jet
conditions from the okness of the ambient datum, and the character ideal of the column character; the character values
(\lref{GlobalTheory.Mem.ChiVal.chi_values}); the selection; the membership of $2,q,d,g_0,g_1$ and of the level primes $T_u$ in the
prime set; smallness; the constants at $a$; $\mu_0$, $\eta$, the character tie and the condition on $\mu q G$. Each is a component of
the okness of $I$ or of the cited lemmas.
\end{proof}

\begin{lemma}[Chain theta data are ok]\label{pf:stageU:chthok}
Let $I$ be an ok column inclusion whose jets lie in the stages and whose record satisfies the equality and vector properties for its
character. Then for every index $a$ and word $u$ the chain theta datum of $u$ at $(a,m-2)$ is ok.
\formallabel{GlobalTheory.ChainTheta.chTh_ok_of_colIn'}
\end{lemma}

\begin{proof}
This is \lref{GlobalTheory.ChainTheta.chTh_ok_of_colIn}, applied with the multiplicity condition of $I$: the chain theta datum is
built from the normalised theta family of the column choice, whose equivariance laws and lattice conditions hold by
\cref{pf:stageU:hypcolin} and the cited properties of normalised families.
\end{proof}

\begin{lemma}[Equivariance laws of normalised families]\label{pf:stageU:natequi}
For a curve ideal $\mathfrak c$, a twelfth root $\delta$ of the discriminant, a kernel base and a character $\chi$, the normalised
theta family satisfies the equivariance laws.
\formallabel{GlobalTheory.T.normalisedFamily_satisfies_equivarianceLaws}
\end{lemma}

\begin{proof}
This is \lref{GlobalTheory.Mem.NatEqui.eqP} (the normalised family transforms under the CM action and the twisting by the
equivariance laws, from the transformation law of the theta functions of the curve lattice).
\end{proof}

\begin{lemma}[Genericity generators are units]\label{pf:stageU:normgen}
Let $\ell$ be a prime which is not small and $\ne q$, $\lambda:\bar\Q\to\C_\ell$ an embedding, and $M\subset\C$ a number field containing the
values of the embeddings of $D$. For every $\tau$ and every generator $s$ of the ideal $\mathfrak l_\tau$ of $\OO_M$, $|\lambda(s)|=1$.
\formallabel{GlobalTheory.EigenGenerator.norm_generator_eq_one}
\end{lemma}

\begin{proof}
$|\lambda(s)|\le1$ because the coordinates of elements of $\mathfrak l_\tau$ are $\lambda$-integral
(\lref{GlobalTheory.EigenGenerator.norm_coord_le}). Conversely there is $v$, vanishing at all $\tau'\ne\tau$, with $|\lambda(v_\tau)|\ge1$
(\lref{GlobalTheory.EigenGenerator.exists_vanishing}); then $v_\tau\in\mathfrak l_\tau=(s)$, $v_\tau=cs$ with $c\in\OO_M$, and
$1\le|\lambda(c)||\lambda(s)|\le|\lambda(s)|$.
\end{proof}

\begin{lemma}[Bounded continuous representatives]\label{pf:stageU:contrep}
Let $B$ be row automorphic data with discrete series data $d$ whose lowest $K$-types are of the given type, $n\ge1$ and level $g\ge1$.
There is $C$ such that for every intertwiner $A$ of norm $\le1$ whose image of $\otimes w$ is fixed by $K(g)$, the vector $A(\otimes w)$
has a continuous representative $f$ with $|f(x)|\le C\,\tau_1(x)^{-2n^2f}\|A(\otimes w)\|$.
\formallabel{GlobalTheory.EvalBound.exists_contRep_bound_of_isKTypeOf}
\end{lemma}

\begin{proof}
There is a smooth compactly supported function $a$ on the archimedean group whose operator reproduces $\otimes w$
(\lref{GlobalTheory.StepAverage.exists_smooth_average_of_isKTypeOf}); its operator $T$ on $L^2$ is defined by absolute integration
(\lref{GlobalTheory.absOpInt}), and for the convolution with a smooth compactly supported function there is $C$ such that
$K(g)$-fixed vectors of the form $T(v)$ have continuous representatives with the stated bound
(\lref{GlobalTheory.SmoothGrowth.convolutionSmoothL2}). Since $A$ intertwines, $T(A(\otimes w))=A(\otimes w)$.
\end{proof}

\begin{lemma}[The Siegel--Weil formula for one vector]\label{pf:stageU:onevec}
The Siegel--Weil formula in the convergent range implies the one-vector mean-value formula used in \cref{pf:stageU:volratio}.
\formallabel{GlobalTheory.GaussMean.oneVectorCc_of_siegelWeil}
\end{lemma}

\begin{proof}
The finite-volume part is a component of the Siegel--Weil formula. For the mean-value formula take $n=1$ and $T=(\beta)$, $\beta\ne0$, so
$\det T\ne0$ and the Gram condition $(h_J(x_a,x_b))=T$ is $h_J(x,x)=\beta$ (\lref{GlobalTheory.GaussMean.OneVec.gram_form}). The
densities for the self-dual measures of the hermitian $1\times1$ matrices are carried to densities for the pushforwards under
$x\mapsto h_J(x,x)$ through the identification of $1\times1$ hermitian matrices with $F$ (\lref{GlobalTheory.GaussMean.OneVec.dens_of_herm},
\lref{GlobalTheory.GaussMean.OneVec.sd_of_herm}).
\end{proof}

\begin{lemma}[Tied kernel entries]\label{pf:stageU:kerent}
Let $I$ be an ok column datum with pins, laws $l$ and later stages $Z$, and let size data, ties, instances, constants, division data,
extra data, lift objects $O$ and kernel entries $E$ be given such that the $Q_p$-condition and the standard object condition hold, the
unfolding is complete, the entries are ok and every kernel is served by $E$. Then there is a tied kernel entry with objects $O$ and
entries $E$ which is standard, adapted and top-pinned.
\formallabel{GlobalTheory.Mem.Col2.TC.kerEntTied_true_ent}
\end{lemma}

\begin{proof}
The kernel entries are compatible with the true columns (\lref{GlobalTheory.T.kernelEntries_compatible_trueColumns}), which gives a
tied kernel datum with the standard tie, the objects $O$ and the top pin; replacing its entries by $E$ keeps these properties and
makes it adapted (\lref{GlobalTheory.Mem.Col2.TC.kerEntTied_withEnt}).
\end{proof}

\begin{lemma}[Fitting stages form states]\label{pf:stageU:fitstate}
Assume the Weil representations of unitary pairs, existence of discrete series, Harish-Chandra finiteness and rapid decay. For a
pre-state $P$ and an over-datum $o$ which fits, has the group of $P$, $\eta_F=6x^{-C_F}$ and the law of the datum, there is a state with
over-datum $o$.
\formallabel{GlobalTheory.T.fittingLaterStages_row_is_constructionState}
\formallabel{GlobalTheory.Mem.R3.Rw.st_of_fits}
\end{lemma}

\begin{proof}
The moved row lift data exist for fitting stages (\lref{GlobalTheory.T.exists_movedRowLiftData_to_fittingStages}), and a reading of
the row lift family of the moved data exists from the inputs provided by the four published statements
(\lref{GlobalTheory.T.exists_reading_of_rowLiftFamily_movedData}, \lref{GlobalTheory.LinksInputs.inputs_of_hypotheses}); together
with $o$ these form a state.
\end{proof}

\begin{lemma}[Injectivity of adjoints on eigen-parts]\label{pf:stageU:adjinj}
Let $T:H'\to H$ be a linear map of finite-dimensional Hilbert spaces, $(P_c)$ subspaces of $H$ and $s_c\ne0$ scalars with $TT^*x=s_cx$
for $x\in P_c$. If $x\in\sum_cP_c$ and $T^*x=0$, then $x=0$.
\formallabel{GlobalTheory.AdjointInjectiveOfParts.adjoint_inj_of_parts}
\end{lemma}

\begin{proof}
Each $P_c$ lies in the range of $T$, since $y=T(T^*(s_c^{-1}y))$ for $y\in P_c$; hence $x=Tw$ for some $w$. Then
$\langle x,x\rangle=\langle Tw,x\rangle=\langle w,T^*x\rangle=0$, so $x=0$.
\end{proof}

% ---------------------------------------------------------------------------
\subsubsection{Two further lemmas}\label{sec:pf:stageU:misc}

\begin{lemma}[Absolute values on the maximal unramified extension]\label{pf:stageU:kurval}
Let $x\ne0$ lie in the maximal unramified extension of $\Q_\ell$ in $\C_\ell$. Then $|x|=\ell^m$ for some $m\in\Z$.
\formallabel{ReciprocityTheory.ChartRestriction.L1.kur_val}
\end{lemma}

\begin{proof}
There is $y$ in a finite unramified extension of $\Q_\ell$ with $|x-y|<|x|$ (\lref{ReciprocityTheory.ChartRestriction.L1.kur_approx}); then
$|x|=|y|$ by the ultrametric inequality and $y\ne0$. In a finite unramified extension $\ell$ is a uniformiser, so $y=\ell^{-m}a/b$ with
units $a,b$ (\lref{ReciprocityTheory.ChartRestriction.L1.kmu_shape}), and $|y|=\ell^m$.
\end{proof}

\begin{lemma}[Explicit Vandermonde counts for families]\label{pf:stageU:divlevel}
Let $\stageUFm$ be a family such that eventually every member satisfies the mass ratio, the base-volume identity and the two-sided
degree bounds of \cref{pf:stageU:tusizemin}. Then $\stageUFm$ satisfies the explicit Vandermonde count bounds for all constants
$(C_0,C_Y)$.
\formallabel{Formalism.MainNotions.StandardLattice.division_level_prime_family}
\end{lemma}

\begin{proof}
Fix $(C_0,C_Y)$. There are thresholds beyond which: the ramified rank lines are free (\lref{Formalism.MainNotions.StandardLattice.ramified_rank_lines_of_free}),
the part-five bounds follow from the mass and volume identities (\lref{Formalism.MainNotions.StandardLattice.part_five_family_of_mass_volume}),
and the gaps satisfy the bound $\mathrm{gap}_B$ with constant $8192$ (\lref{Formalism.MainNotions.StandardLattice.gapB_ev}). Beyond the
maximum of these thresholds and those of the three hypotheses, every member satisfies the sharp conclusion five (both components:
the line part and the part-five part), and with the Vandermonde variation bound
(\lref{Formalism.MainNotions.StandardLattice.vandermonde_variation_holds}) this gives the explicit count bounds
(\lref{Formalism.MainNotions.StandardLattice.part_five_explicit_of_propositions}).
\end{proof}


